\documentclass[12pt]{article}
\usepackage{expl3}
\makeatletter
\usepackage{cmap}
\usepackage[T1,T2A]{fontenc}
\usepackage[utf8]{inputenc}
\usepackage[ukrainian]{babel}
\NewDocumentCommand{\HComp}{}{\mathbin{\ast}}
\usepackage{amssymb}
\usepackage{bm}
\usepackage{xcolor}
\usepackage{fontawesome}
\usepackage{mleftright}
\usepackage{xparse}
\usepackage{ifthen}
\usepackage{scalerel}
\usepackage{epigraph}
\usepackage{datetime}
\usepackage{mathpartir}
\usepackage{array}
\usepackage{mathtools}
\let\Bbbk\relax
\usepackage{etoolbox}
\usepackage{amsthm}
\usepackage{thmtools}
\usepackage{thm-restate}
\usepackage{enumitem}
\usepackage{hyperref}
\usepackage{euscript}
% lmcs compatibility shims
\newcommand{\lmcsorcid}[1]{}
\newcommand{\address}[1]{}
\providecommand{\email}[1]{}
% Theorem environments (normally provided by lmcs.cls)
\declaretheoremstyle[spaceabove=6pt,spacebelow=6pt,headfont=\bfseries,notefont=\normalfont,notebraces={(}{)},bodyfont=\itshape]{plain}
\declaretheoremstyle[spaceabove=6pt,spacebelow=6pt,headfont=\bfseries,notefont=\normalfont,notebraces={(}{)},bodyfont=\normalfont]{definition}
\declaretheoremstyle[spaceabove=6pt,spacebelow=6pt,headfont=\bfseries\itshape,notefont=\normalfont,notebraces={(}{)},bodyfont=\normalfont]{remark}
\declaretheorem[style=plain,name=Theorem]{thm}
\declaretheorem[style=plain,name=Lemma,sibling=thm]{lem}
\declaretheorem[style=plain,name=Proposition,sibling=thm]{prop}
\declaretheorem[style=plain,name=Corollary,sibling=thm]{cor}
\declaretheorem[style=definition,name=Definition,sibling=thm]{defi}
\declaretheorem[style=definition,name=Definition,sibling=thm]{defiC}
\declaretheorem[style=definition,name=Example,sibling=thm]{exa}
\declaretheorem[style=definition,name=Notation,sibling=thm]{nota}
\declaretheorem[style=remark,name=Remark,sibling=thm]{rem}
\declaretheorem[style=remark,name=Convention,sibling=thm]{conv}
\definecolor{Matterhorn}{RGB}{77,77,77}
\definecolor{RegalBlue}{RGB}{3,69,117}
\definecolor{RedDevil}{RGB}{134,1,17}
\hypersetup{breaklinks=true,colorlinks=true,linkcolor=RedDevil,urlcolor=RegalBlue,citecolor=RedDevil,%
            hypertexnames=false}

\newcommand{\textcite}{\cite}
\newcommand{\parencite}{\cite}

\NewDocumentCommand\Mute{m}{%
  {\color{gray} #1}%
}
\NewDocumentCommand\GrmSep{}{\Mute{\;\mid\;}}
\newcommand\GrmContinue{\multicolumn{2}{l}{}} % for weird reasons this must be \newcommand
\NewDocumentCommand{\Alert}{m}{\textcolor{red}{#1}}
\NewDocumentCommand\JdgFrame{m}{%
  \fbox{$\displaystyle #1$}
}
\NewDocumentCommand{\Pow}{}{\mathcal{P}}
\declaretheorem[style=plain]{axiom}
\renewcommand*{\theaxiom}{\Alph{axiom}}
%\newtheoremstyle{special}% name
%  {\topsep}%   Space above
%  {\topsep}%   Space below
%  {\itshape}%  Body font
%  {}%          Indent amount
%  {\bfseries}% Theorem head font
%  {}%          Punctuation after theorem head -- blank
%  {0.5em}%     Space after theorem head (0.5em is the default)
%  {{\thmname{#1}\thmnumber{ #2$^{\bm *}\!\!\!$.}\thmnote{\ \textmd{(#3)}}}}
\NewDocumentCommand{\etc}{}{etc.}
\NewDocumentCommand{\eg}{}{e.g.}
\NewDocumentCommand{\ie}{}{i.e.}
\NewDocumentCommand{\viz}{}{viz.}
\NewDocumentCommand{\opcit}{}{\emph{op.\ cit.}}
\NewDocumentCommand{\loccit}{}{\emph{loc.\ cit.}}
\NewDocumentCommand{\etal}{}{et al.}
%% From https://tex.stackexchange.com/questions/42726/align-but-show-one-equation-number-at-the-end
\NewDocumentCommand{\TagEq}{o}{\IfValueTF{#1}{\tag{#1}}{\stepcounter{equation}\tag{\theequation}}}
\newcommand{\Sub}[1]{_{\DelimProtect{#1}}}
\newcommand{\Sup}[1]{^{\DelimProtect{#1}}}
\allowdisplaybreaks%
% Misc
\NewDocumentCommand{\defeq}{}{\coloneq}
\NewDocumentCommand{\Mode}{}{\mathcal{M}}
\NewDocumentCommand{\Restrict}{mm}{#1\vert\Sub{#2}}
\NewDocumentCommand{\Image}{}{\mathsf{im}}
% Names of type theories
\NewDocumentCommand{\NamedTypeTheory}{m}{\ensuremath{\mathsf{#1}}}
\NewDocumentCommand{\MLTT}{}{\NamedTypeTheory{MLTT}}
\NewDocumentCommand{\MTT}{}{\NamedTypeTheory{MTT}}
\NewDocumentCommand{\HoTT}{}{\NamedTypeTheory{HoTT}}
\NewDocumentCommand{\STT}{}{\NamedTypeTheory{STT}}
\NewDocumentCommand{\BCTT}{}{\NamedTypeTheory{BCTT}}
\NewDocumentCommand{\TTT}{}{\NamedTypeTheory{TT}\textsubscript{$\boxslash$}}
\NewDocumentCommand{\Interp}{m}{\DelimSurround{\llbracket}{\rrbracket}{#1}}
\NewDocumentCommand{\Pushout}{m}{\mathrel{\DelimMin{1}\textstyle\amalg\Sub{#1}}}
\NewDocumentCommand{\Inj}{m}{\mathsf{in}\Sub{#1}}
\NewDocumentCommand{\BluntJoin}{}{\mathrel{\diamond}}
\NewDocumentCommand{\Proj}{O{}}{\pi\Sub{#1}}
% Judgments
\NewDocumentCommand{\IsCx}{mo}{\vdash #1\IfValueT{#2}{{} \mathrel{@} #2}}
\NewDocumentCommand{\EqCx}{mmo}{\vdash #1 = #2\IfValueT{#3}{{} \mathrel{@} #3}}
\NewDocumentCommand{\IsDCx}{mm}{\vdash #1; #2\,\mathsf{dcx}}
\NewDocumentCommand{\IsTy}{O{\Gamma}mo}{#1 \vdash #2\IfValueT{#3}{{} \mathrel{@} #3}}
\NewDocumentCommand{\EqTy}{O{\Gamma}mmo}{#1 \vdash #2 = #3\IfValueT{#4}{{} \mathrel{@} #4}}
\NewDocumentCommand{\IsTm}{O{\Gamma}mmo}{#1 \vdash #2 : #3\IfValueT{#4}{{} \mathrel{@} #4}}
\NewDocumentCommand{\EqTm}{O{\Gamma}mmmo}{#1 \vdash #2 = #3 : #4\IfValueT{#5}{{} \mathrel{@} #5}}
\NewDocumentCommand{\IsSb}{O{\Gamma}mmo}{#1 \vdash #2 : #3\IfValueT{#4}{{} \mathrel{@} #4}}
\NewDocumentCommand{\EqSb}{O{\Gamma}mmmo}{#1 \vdash #2 = #3 : #4\IfValueT{#5}{{} \mathrel{@} #5}}
\NewDocumentCommand{\DeclVar}{mmm}{#1 \mathrel{:_{#2}} #3}
\NewDocumentCommand{\At}{}{\cdot}
\NewDocumentCommand{\TotalType}{m}{\widetilde{#1}}
% Types and terms
\NewDocumentCommand{\Uni}{o}{\mathcal{U}\IfValueT{#1}{_{#1}}}
\NewDocumentCommand{\Dec}{o}{\mathsf{El}\IfValueT{#1}{_{#1}}}
\NewDocumentCommand{\Nat}{}{{\mathsf{Nat}}}
\NewDocumentCommand{\Unit}{}{{\ObjTerm{}}}
\NewDocumentCommand{\Bool}{}{{\mathsf{Bool}}}
\NewDocumentCommand{\MFn}{O{\mu}mm}{\prn{#1 \mid #2} \to #3}
\NewDocumentCommand{\Modify}{O{\mu}m}{\gl{#1 \mid #2}}
\NewDocumentCommand{\ModifyCode}{O{\mu}G{-}}{\mathsf{Mod}_{#1}\prn{#2}}
\NewDocumentCommand{\StrictModify}{O{\mu}m}{#1 \To #2}
\NewDocumentCommand{\Refl}{}{{\mathsf{refl}}}
\NewDocumentCommand{\Ax}{}{{\star}}
\NewDocumentCommand{\Lam}{om}{\lambda \IfValueTF{#1}{#1.\,}{} #2}
\NewDocumentCommand{\App}{mm}{#1\prn{#2}}
\NewDocumentCommand{\Zero}{}{\mathsf{z}}
\NewDocumentCommand{\Succ}{}{\mathsf{suc}}
\NewDocumentCommand{\True}{}{{\mathsf{tt}}}
\NewDocumentCommand{\False}{}{{\mathsf{ff}}}
\NewDocumentCommand{\MkMod}{O{\mu}m}{\mathsf{mod}\IfValueT{#1}{_{#1}}\prn{#2}}
\NewDocumentCommand{\StrictMkMod}{om}{\brc{#2}\IfValueT{#1}{_{#1}}}
\NewDocumentCommand{\StrictUnMod}{om}{{!}\IfValueT{#1}{_{#1}}\,#2}
\NewDocumentCommand{\UnMod}{om}{\mathsf{unmod}\IfValueT{#1}{_{#1}}\prn{#2}}
\NewDocumentCommand{\UnModApp}{omm}{\mathsf{unmod}\IfValueT{#1}{_{#1}}\prn{#2, #3}}
\NewDocumentCommand{\LetMod}{d<>omom}{%
  \mathsf{let}\IfValueT{#1}{_{#1}}\ \MkMod[#2]{\IfValueTF{#4}{#4}{-}} \gets #3\ \mathsf{in}\ #5%
}
\NewDocumentCommand{\IdRec}{o}{\mathsf{J}\IfValueT{#1}{_{#1}}}
\NewDocumentCommand{\BoolRec}{o}{\mathsf{if}\IfValueT{#1}{_{#1}}}
\NewDocumentCommand{\NatRec}{o}{\mathsf{rec}\IfValueT{#1}{_{#1}}}
\NewDocumentCommand{\Sb}{mm}{#1\brk{#2}}
\NewDocumentCommand{\ECx}{mm}{#1.#2}
\NewDocumentCommand{\MECx}{mO{\mu}m}{#1.\prn{#2 \mid #3}}
\NewDocumentCommand{\MECxV}{mmO{\mu}m}{#1, #2 :_{#3} #4}
\NewDocumentCommand{\EmpCx}{O{}}{\ObjTerm{#1}}
\NewDocumentCommand{\EmpSb}{O{}}{\ArrTerm{#1}}
\NewDocumentCommand{\ESb}{mm}{#1.#2}
\NewDocumentCommand{\ISb}{}{\ArrId{}}
\NewDocumentCommand{\Wk}{o}{{\uparrow}\IfValueT{#1}{^{#1}}}
\NewDocumentCommand{\Var}{o}{\mathbf{v}}
\NewDocumentCommand{\LockCx}{mO{\mu}}{#1.\brc{#2}}
\NewDocumentCommand{\LockCxV}{mO{\mu}}{#1, \brc{#2}}
\NewDocumentCommand{\LockSb}{mO{\mu}}{#1.\brc{#2}}
\NewDocumentCommand{\KeySb}{mm}{\ifthenelse{\equal{#2}{}}{}{#2.}\brc{#1}}
\NewDocumentCommand{\Open}{O{}}{\medcirc_{#1}}
\NewDocumentCommand{\Closed}{}{\medbullet}
\NewDocumentCommand{\Prop}{}{\mathsf{HProp}}
\NewDocumentCommand{\HSet}{}{\mathsf{HSet}}
\NewDocumentCommand{\Trunc}{m}{\vvrt{#1}}
\NewDocumentCommand{\FibTo}{}{\twoheadrightarrow}
%\NewDocumentCommand\TypedHole{m}{
%  \colorbox{RegalBlue!10}{${\color{black!75}#1}$}
%}
\NewDocumentCommand{\Con}{m}{\mathsf{#1}}
\NewDocumentCommand{\Kwd}{m}{\mathbf{#1}}
\NewDocumentCommand{\Ap}{}{\Con{ap}}
\NewDocumentCommand{\Coe}{}{\Con{coe}}
\NewDocumentCommand{\ZApp}{}{\mathbin{\circledast}}
\NewDocumentCommand{\Simp}{}{\boxslash}
\NewDocumentCommand{\Hom}{omm}{\mathsf{hom}\IfValueT{#1}{_{#1}}\prn{#2, #3}}
\NewDocumentCommand{\DepSimp}{}{\hat{\boxslash}}

% Useful Propositions
\NewDocumentCommand{\IsContr}{}{\Con{isContr}}
\NewDocumentCommand{\IsEquiv}{}{\Con{isEquiv}}
\NewDocumentCommand{\IsProp}{}{\Con{isProp}}
\NewDocumentCommand{\IsSet}{}{\Con{isSet}}
\NewDocumentCommand{\IsConn}{}{\Con{isConn}}
\NewDocumentCommand{\IsSegal}{}{\Con{isSegal}}
\NewDocumentCommand{\IsSegalGraph}{}{\Con{isSegalGraph}}
\NewDocumentCommand{\IsRezk}{}{\Con{isRezk}}
\NewDocumentCommand{\IsDisc}{}{\Con{isDisc}}
\NewDocumentCommand{\IsSimp}{}{\Con{isSimp}}
\NewDocumentCommand{\IsIso}{}{\Con{isIso}}
\NewDocumentCommand{\IsCocartArr}{}{\Con{isCocartArr}}
\NewDocumentCommand{\IsCocartFam}{}{\Con{isCocartFam}}
\NewDocumentCommand{\IsCovFam}{}{\Con{isCov}}
\NewDocumentCommand{\ACov}{}{\Con{ACov}}
\NewDocumentCommand{\IsACov}{}{\Con{isACov}}
\NewDocumentCommand{\IsInnFam}{}{\Con{isInnFam}}
\NewDocumentCommand{\Spec}{}{\Con{Spec}}
\let\IsCov\IsCovFam
\NewDocumentCommand{\Transpose}{sm}{\IfBooleanTF{#1}{{#2}^\rightarrow}{{#2}^\leftarrow}}
\NewDocumentCommand{\Int}{}{\mathbb{I}}
\NewDocumentCommand{\Sp}{}{\mathsf{Sp}}
\NewDocumentCommand{\Amazing}{m}{{#1}_{\mathbb{I}}}
\newcommand{\Space}{\mathcal{S}}
\newcommand{\SSimp}{\mathsf{S}}
\newcommand{\Cat}{\mathsf{Cat}}
\newcommand{\CAT}{\mathsf{CAT}}
\newcommand{\SET}{\mathsf{SET}}
\newcommand{\GRPD}{\mathsf{GRPD}}
\newcommand{\SIMP}{\Delta}
\newcommand{\PSH}[1]{\widehat{#1}}
\newcommand{\Yo}{\mathsf{Yo}}
\newcommand{\id}{\mathsf{id}}
\newcommand{\arr}{\mathsf{arr}}
\newcommand{\iso}{\mathsf{iso}}
\newcommand{\Iso}[2]{\mathsf{Iso}\prn{#1, #2}}
\newcommand{\Equiv}{\mathsf{Equiv}}
\newcommand{\Map}[2]{\mathsf{Map}\prn{#1, #2}}
\NewDocumentCommand\Compr{mm}{\DelimBetweenSurround{\lbrace}{\vert}{\rbrace}{#1}{#2}}
\NewDocumentCommand{\Ext}{mmm}{\DelimBetweenSurround{\lbrace}{\vert}{\rbrace}{#1}{#2 \mapsto #3}}
\NewDocumentCommand{\Kan}{}{\mathsf{Kan}}
\NewDocumentCommand{\Segal}{}{\mathsf{Segal}}
\NewDocumentCommand{\Rezk}{}{\mathsf{Rezk}}
\NewDocumentCommand{\IsKan}{}{\Con{isKan}}
\NewDocumentCommand{\ArrId}{O{}}{\mathsf{id}\Sub{#1}}
\NewDocumentCommand{\ObjTerm}{O{}}{\mathsf{1}\Sub{#1}}
\NewDocumentCommand{\ArrTerm}{O{}}{\mathsf{!}\Sub{#1}}
\NewDocumentCommand{\IdIso}{O{}}{\mathsf{id}\Sub{#1}}
\NewDocumentCommand{\Com}{O{}}{\mathsf{comp}\Sub{#1}}
\NewDocumentCommand{\Cst}{m}{\mathsf{const}\Sub{#1}}
\NewDocumentCommand{\UU}{}{\Uni}
\NewDocumentCommand{\OO}{}{\mathcal{O}}
\NewDocumentCommand{\Sharp}{}{\sharp}
\NewDocumentCommand{\Flat}{}{\flat}
\NewDocumentCommand{\yo}{}{\Yo}
\NewDocumentCommand{\yon}{}{\Yo}
\newcommand*{\Dist}{\mathrm{Dist}}
\newcommand*{\Core}{\mathrm{Core}}
\NewDocumentCommand{\Glue}{}{\mathsf{Gl}}
\RenewDocumentCommand{\Modify}{om}{\mathsf{#2}\IfValueT{#1}{_{#1}}}
\NewDocumentCommand{\Adjoint}{}{\dashv}
\NewDocumentCommand{\Lan}{}{\mathsf{Lan}}
\NewDocumentCommand{\Ran}{}{\mathsf{Ran}}
\NewDocumentCommand{\HM}{}{\mathsf{H}}
\NewDocumentCommand{\CC}{}{\mathcal{C}}
\NewDocumentCommand{\DD}{}{\mathcal{D}}
\NewDocumentCommand{\EE}{}{\mathcal{E}}
\NewDocumentCommand{\FF}{}{\mathcal{F}}
\NewDocumentCommand{\pb}{}{\text{pb}}
\NewDocumentCommand{\Op}{m}{{#1}^{\mathsf{op}}}
\NewDocumentCommand{\Whisker}{}{\cdot}
\NewDocumentCommand{\prn}{m}{\left(#1\right)}
\NewDocumentCommand{\brk}{m}{\left[#1\right]}
\NewDocumentCommand{\brc}{m}{\left\{#1\right\}}
\NewDocumentCommand{\ang}{m}{\left\langle#1\right\rangle}
\NewDocumentCommand{\vrt}{m}{\left\vert#1\right\vert}
\NewDocumentCommand{\vvrt}{m}{\left\Vert#1\right\Vert}
\NewDocumentCommand{\MorToFun}{}{\mathsf{mor2fun}}
\NewDocumentCommand{\Coeq}{}{\mathsf{Coeq}}
\NewDocumentCommand{\GlueEq}{}{\mathsf{glue}}
\NewDocumentCommand{\Const}{}{\Con{const}}
\newcommand{\Comma}[2]{#1 \mathbin{\downarrow} #2}
\NewDocumentCommand{\COMMA}{mm}{\Comma{#1}{#2}}
\NewDocumentCommand{\IdToIso}{}{\mathsf{IdToIso}}
\NewDocumentCommand{\IdToArr}{}{\mathsf{IdToArr}}
\newcommand{\GM}{\flat}
\newcommand{\SM}{\sharp}
\newcommand{\PM}{\mathfrak{p}}
\newcommand{\AM}{\bar{\mathfrak{p}}}
\newcommand{\OM}{\mathsf{op}}
\newcommand{\TM}{\mathsf{tw}}
\NewDocumentCommand{\Skel}{}{\mathsf{skel}}
\NewDocumentCommand{\Embeds}{}{\hookrightarrow}
\NewDocumentCommand{\Covers}{}{\twoheadrightarrow}
\NewDocumentCommand{\NonDeg}{}{\mathsf{ND}}
\NewDocumentCommand{\Rezkify}{}{\operatorname{\bigcirc_{\Con{cat}}}}
\NewDocumentCommand{\Grpdify}{}{\operatorname{\bigcirc_{\Con{grpd}}}}
\NewDocumentCommand{\LocalLAdj}{}{\Con{LocalLAdj}}
\NewDocumentCommand{\LocalRAdj}{}{\Con{LocalRAdj}}
\NewDocumentCommand{\CoCont}{mm}{\Con{CC}\prn{#1,#2}}
\NewDocumentCommand{\SLICE}{mm}{{#1}_{/#2}}
\NewDocumentCommand{\COSLICE}{mm}{{#1}_{#2/}}
\NewDocumentCommand{\Nv}{}{\mathsf{N}}
\NewDocumentCommand{\ObjInit}{m}{\mathsf{0}\Sub{#1}}

\ExplSyntaxOn
%% This is file `jmsdelim.sty',
\RequirePackage{l3keys2e}
\RequirePackage{ifluatex}
\keys_define:nn { jmsdelim } {
  size~commands .clist_set:N = \l__jmsdelim_size_cmds,
}
\keys_set:nn { jmsdelim } {
  size~commands = {relax,jmsdelim_big:n,jmsdelim_Big:n,jmsdelim_bigg:n,jmsdelim_Bigg:n},
}
\cs_new:Npn \jmsdelim_big:n #1 {
  {\mathchoice{\big #1} {\big #1}{\big #1}{#1}}
}
\cs_new:Npn \jmsdelim_Big:n #1 {
  {\mathchoice{\Big #1} {\Big #1}{\big #1}{#1}}
}
\cs_new:Npn \jmsdelim_bigg:n #1 {
  {\mathchoice{\bigg #1} {\bigg #1}{\big #1}{#1}}
}
\cs_new:Npn \jmsdelim_Bigg:n #1 {
  {\mathchoice{\Bigg #1} {\Bigg #1}{\big #1}{#1}}
}
\int_new:N \g__jmsdelim_size
\int_new:N \g__jmsdelim_size_up
\int_gset:Nn \g__jmsdelim_size {0}
\int_gset:Nn \g__jmsdelim_size_up {0}
\cs_new:Npn \__jmsdelim_clist_item:Nn #1 #2 {
  \clist_item:Nn #1 {
    \int_min:nn { #2 } {\clist_count:N #1}
  }
}
\cs_new:Npn \__jmsdelim_setup_sizes: {
  \int_gset:Nn \g__jmsdelim_size {
    \int_max:nn \g__jmsdelim_size \g__jmsdelim_size_up
  }
  \cs_set_eq:Nc \jmsdelim_size_cmd: {
    \__jmsdelim_clist_item:Nn \l__jmsdelim_size_cmds {
      \g__jmsdelim_size + 1
    }
  }
}
\cs_new:Npn \__jmsdelim_luatex_save_mathstyle:N #1 {
  \ifcase \mathstyle
    \cs_set_eq:NN #1 \displaystyle
  \or
    \cs_set_eq:NN #1 \crampeddisplaystyle
  \or
    \cs_set_eq:NN #1 \textstyle
  \or
    \cs_set_eq:NN #1 \crampedtextstyle
  \or
    \cs_set_eq:NN #1 \scriptstyle
  \or
    \cs_set_eq:NN #1 \crampedscriptstyle
  \or
    \cs_set_eq:NN #1 \scriptscriptstyle
  \or
    \cs_set_eq:NN #1 \crampedscriptscriptstyle
  \fi
}
\cs_new:Npn \__jmsdelim_restore_mathstyle: {
  \SavedStyle
}
\cs_new:Npn \__jmsdelim_save_mathstyle:n #1 {
  \ifluatex
    \__jmsdelim_luatex_save_mathstyle:N \__jmsdelim_restore_mathstyle:
    #1
  \else
    \ThisStyle{#1}
  \fi
}
\cs_new:Npn \jmsdelim_scope:nn #1 #2 {
  \group_begin:
    \int_set:Nn \l_tmpa_int \g__jmsdelim_size_up
    \int_gset:Nn \g__jmsdelim_size_up 0
    \int_gset:Nn \g__jmsdelim_size 0
    \group_begin:
      \__jmsdelim_save_mathstyle:n {
        #1
        \__jmsdelim_setup_sizes:
        #2
      }
    \group_end:
    \int_gset:Nn \g__jmsdelim_size_up {\int_max:nn \g__jmsdelim_size_up \l_tmpa_int}
  \group_end:
}
\cs_new:Npn \jmsdelim_hbox_set:Nn #1 #2 {
   \mode_if_math:TF
     { \hbox_set:Nn #1 {$\mathsurround=0pt\__jmsdelim_restore_mathstyle: #2$} }
     { \hbox_set:Nn #1 { #2 } }
}
\cs_new:Npn \jmsdelim_surround:nnn #1 #2 #3 {
  \jmsdelim_scope:nn {
    \jmsdelim_hbox_set:Nn \l_tmpa_box {#3}
  }{
    \mathopen\jmsdelim_size_cmd: {#1}
    \box_use:N \l_tmpa_box
    \mathclose\jmsdelim_size_cmd: {#2}
  }
}
\cs_new:Npn \jmsdelim_protect:n #1 {
  \group_begin:
    \int_set:Nn \l_tmpa_int \g__jmsdelim_size_up
    \int_set:Nn \l_tmpb_int \g__jmsdelim_size
    \group_begin: #1 \group_end:
    \int_gset:Nn \g__jmsdelim_size_up \l_tmpa_int
    \int_gset:Nn \g__jmsdelim_size \l_tmpb_int
  \group_end:
}
\cs_new:Npn \jmsdelim_between:nnn #1 #2 #3 {
  \jmsdelim_scope:nn {
    \jmsdelim_hbox_set:Nn \l_tmpa_box {#2}
    \jmsdelim_hbox_set:Nn \l_tmpb_box {#3}
  }{
    \box_use:N \l_tmpa_box
    \mathrel{\jmsdelim_size_cmd: {#1}}
    \box_use:N \l_tmpb_box
  }
}
\cs_new:Npn \jmsdelim_between:nnnnn #1 #2 #3 #4 #5 {
  \jmsdelim_scope:nn {
    \jmsdelim_hbox_set:Nn \l_tmpa_box {#4}
    \jmsdelim_hbox_set:Nn \l_tmpb_box {#5}
  }{
    \mathopen\jmsdelim_size_cmd: {#1}
    \box_use:N \l_tmpa_box
    \mathrel{\jmsdelim_size_cmd: {#2}}
    \box_use:N \l_tmpb_box
    \mathclose\jmsdelim_size_cmd: {#3}
  }
}
\NewDocumentCommand\DelimMin{m}{
  \int_gset:Nn \g__jmsdelim_size_up {\int_max:nn \g__jmsdelim_size_up {#1}}
}
\NewDocumentCommand\DelimSurround{mmm}{
  \jmsdelim_surround:nnn {#1} {#2} {#3}
}
\NewDocumentCommand\DelimBetween{mmm}{
  \jmsdelim_between:nnn {#1} {#2} {#3}
}
\NewDocumentCommand\DelimBetweenSurround{mmmmm}{
  \jmsdelim_between:nnnnn {#1} {#2} {#3} {#4} {#5}
}
\newcommand{\DelimProtect}[1]{#1}
\NewDocumentCommand\DelimPrn{m}{
  \jmsdelim_surround:nnn {(} {)} {#1}
}
\NewDocumentCommand\DelimBrk{m}{
  \jmsdelim_surround:nnn {[} {]} {#1}
}
\NewDocumentCommand\DelimBrc{m}{
  \jmsdelim_surround:nnn {\lbrace} {\rbrace} {#1}
}
\NewDocumentCommand\DelimBbrk{m}{
  \jmsdelim_surround:nnn {\llbracket} {\rrbracket} {#1}
}
\NewDocumentCommand\DelimGl{m}{
  \jmsdelim_surround:nnn {\langle} {\rangle} {#1}
}
\ExplSyntaxOff
\let\prn\DelimPrn
\let\brk\DelimBrk
\let\brc\DelimBrc
\let\bbrk\DelimBbrk
\let\gl\DelimGl

\ExplSyntaxOn
%% BIG SYMBOLS
\cs_new_protected:Npn \jms_set_mathsurround_to_zero: {
  \dim_set:Nn \mathsurround { 0pt }
}
\cs_new:Npn \jms_scale_bigop:n #1 {
  \tex_mathop:D {
    \vphantom{\sum}
    \mathpalette\jms_scale_bigop_with_mathstyle:nn {#1}
  }
  \slimits@
}
\cs_new:Npn \jms_scale_bigop_with_mathstyle:nn #1 #2 {
  \tex_vcenter:D {
    \hbox_set:Nn \l_tmpa_box {$#1\sum$}
    \hbox_set:Nn \l_tmpb_box {$\jms_set_mathsurround_to_zero: #2$}
    \box_resize_to_ht:Nn \l_tmpb_box {
      \cs_if_eq:NNT #1 \displaystyle {.9}  \box_ht_plus_dp:N \l_tmpa_box
    }
    \box_use:N \l_tmpb_box
  }
}
\NewDocumentCommand\ReixSymbol{}{\jms_scale_bigop:n {\operatorname{\Delta}}}
\NewDocumentCommand\ForallSymbol{}{\jms_scale_bigop:n {\forall}}
\NewDocumentCommand\ExistsSymbol{}{\jms_scale_bigop:n {\exists}}
\cs_new:Npn \jms_bigop:N #1 #2 {
  \mathchoice
    {\textstyle #1 \Sub{#2}}
    {\textstyle #1 \Sub{#2}}
    {\scriptstyle #1 \Sub{#2}}
    {\scriptscriptstyle #1 \Sub{#2}}
}
\cs_new:Npn \jms_define_big_quantifier:NN #1 #2 {
  \NewDocumentCommand{#1}{m}{
    \DelimMin{1}
    \jms_bigop:N {#2} {##1}
  }
}
\jms_define_big_quantifier:NN \Coprod \coprod
\jms_define_big_quantifier:NN \Sum \sum
\jms_define_big_quantifier:NN \Prod \prod
\jms_define_big_quantifier:NN \DLub \bigsqcup
\jms_define_big_quantifier:NN \Disj \bigvee
\jms_define_big_quantifier:NN \Conj \bigwedge
%\jms_define_big_quantifier:NN \Colim \varinjlim
%\jms_define_big_quantifier:NN \Lim \varprojlim
\newcommand{\Colim}[1]{\mathrm{colim}_{#1}}
\newcommand{\Lim}[1]{\mathrm{lim}_{#1}}
\jms_define_big_quantifier:NN \Union \bigcup
\jms_define_big_quantifier:NN \Isect \bigcap
\jms_define_big_quantifier:NN \Reix \ReixSymbol
\jms_define_big_quantifier:NN \Forall \ForallSymbol
\jms_define_big_quantifier:NN \Exists \ExistsSymbol

\cs_new:Npn \iblock_init_local_state: {
  \cs_set_eq:NN \row \iblock_row:
  \cs_set_eq:NN \mrow \iblock_math_row:n
  \cs_set_eq:NN \mhang \iblock_math_hang:nn
  \bool_set:Nn \l_iblock_inside \c_true_bool
}
\bool_new:N \l_iblock_inside
\cs_new:Npn \iblock_container:n #1 {
  \group_begin:
  \raggedright
  \setlength \hangingleftmargin {
    \bool_if:NTF \l_iblock_inside {0em} {\iblockleftmargin}
  }
  \compacthang
    \iblock_init_local_state:
    #1
  \endcompacthang
  \group_end:
}
\cs_new:Npn \iblock_row: { \item }
\cs_new:Npn \iblock_math_row:n #1 {
  \iblock_row: \ensuremath{#1}
}
\cs_new:Npn \iblock_math_hang:nn #1 #2 {
  \iblock_math_row:n {#1}
  \iblock_container:n {#2}
}
\NewDocumentCommand\iblock{+m}{
  \iblock_container:n {
    #1
  }
}
\ExplSyntaxOff

\usepackage{hang}
\usepackage{tikz}
\usetikzlibrary{matrix,arrows,patterns,fit,positioning,calc,decorations.pathmorphing,babel}
\usepackage{tikz-cd}

\tikzset{
  desc/.style={sloped, fill=white,inner sep=2pt},
  upright desc/.style={fill=white,inner sep=2pt},
  pullback/.style = {
    append after command={
      \pgfextra{
        \draw ($(\tikzlastnode) + (.3cm,-.5cm)$) -- ++(0.2cm,0) -- ++(0,0.2cm);
      }
    }
  },
  dotted pullback/.style = {
    append after command={
      \pgfextra{
        \draw [densely dotted] ($(\tikzlastnode) + (.2cm,-.5cm)$) -- ++(0.3cm,0) -- ++(0,0.3cm);
      }
    }
  },
  pushout/.style = {
    append after command={
      \pgfextra{
        \draw ($(\tikzlastnode) + (-.2cm,.5cm)$) -- ++(-0.3cm,0) -- ++(0,-0.3cm);
      }
    }
  },
  between/.style args={#1 and #2}{
    at = ($(#1)!0.5!(#2)$)
  },
  diagram/.style = {
    on grid,
    node distance=2cm,
    commutative diagrams/every diagram,
    every node/.append style = {
      commutative diagrams/every cell,
    }
  },
  fibration/.style = {
    -{Triangle[open]}
  },
  opfibration/.style = {
    -{Triangle}
  },
  lies over/.style = {
    |-{Triangle[open]}
  },
  op lies over/.style = {
    |-{Triangle}
  },
  embedding/.style = {
    {right hook}->
  },
  open immersion/.style = {
    {right hook}-{Triangle[open]}
  },
  embedding*/.style = {
    {left hook}->
  },
  exists/.style = {
    densely dotted
  },
  unique exists/.style = {
    densely dashed
  }
}
\newlength{\dontworryaboutit}
\tikzset{
  inline diagram/.style = {
    commutative diagrams/every diagram,
    commutative diagrams/cramped,
    every node/.append style = {
      commutative diagrams/every cell,
      anchor = base,
      inner sep = 0pt
    },
    every path/.append style = {
      outer xsep = 2pt
    }
  }
}

\ExplSyntaxOn
\cs_new:Npn \jms_tikz_morphism:nnn #1 #2 #3 {
  \tikz[inline~diagram, baseline = (Y.base)]{
    \node (X) {$#2$};
    \node (Y) [right = 3ex~of~X] {$#3$};
    \path[#1] (X) edge (Y);
  }
}
\cs_new:Npn \jms_tikz_sized_morphism:nnnn #1 #2 #3 #4 {
  \tikz[inline~diagram, baseline = (Y.base)]{%
    \settowidth{\dontworryaboutit}{\scriptsize$#2$}
    \pgfmathsetmacro\len{max(width("$xxx$"),7pt+\dontworryaboutit)}
    \node (X) {$#3$};
    \node (Y) [right = \len~pt~of~X] {$#4$};
    \path[#1] (X) edge node [above,yshift=.3ex] {\scriptsize$#2$} (Y);
  }
}
\cs_new:Npn \jms_tikz_auto_morphism:nnnn #1 #2 #3 #4 {
  \DelimMin{1}
  \IfValueTF{#2}{
    #2 : \jms_tikz_morphism:nnn {#1} {#3} {#4}
  }{
    \jms_tikz_morphism:nnn {#1} {#3} {#4}
  }
}
\NewDocumentCommand\PrimMor{mmmm}{
  \IfValueTF{#2}{
    \DelimMin{1}
    \jms_tikz_sized_morphism:nnnn {#1} {#2} {#3} {#4}
  }{
    \jms_tikz_morphism:nnn {#1} {#3} {#4}
  }
}
% Draw a morphism inline, with dynamically sized arrow
\NewDocumentCommand\Mor{D||{->} o m m}{%
  \jms_tikz_auto_morphism:nnnn {#1} {#2} {#3} {#4}
}
\NewDocumentCommand\FibMor{D||{} o m m}{%
  \jms_tikz_auto_morphism:nnnn {fibration,#1} {#2} {#3} {#4}
}
\NewDocumentCommand\OpFibMor{D||{} o m m}{%
  \jms_tikz_auto_morphism:nnnn {opfibration,#1} {#2} {#3} {#4}
}
\NewDocumentCommand\EmbMor{D||{} o m m}{%
  \jms_tikz_auto_morphism:nnnn {embedding,#1} {#2} {#3} {#4}
}
%% TODO: rectify these
\NewDocumentCommand\Span{m m m}{%
  \tikz[inline diagram, baseline = (Y.base)]{%
    \node (X) {$#2$\hspace*{.5ex}};
    \node (Y) [right = 3ex of X] {\hspace*{.5ex}$#1$\hspace*{.5ex}};
    \node (Z) [right = 3ex of Y] {\hspace*{.5ex}$#3$};
    \path[<-] (X) edge (Y);
    \path[->] (Y) edge (Z);
  }%
}
\NewDocumentCommand\CoSpan{m m m}{%
  \tikz[inline diagram, baseline = (Y.base)]{%
    \node (X) {$#2$\hspace*{.5ex}};
    \node (Y) [right = 3ex of X] {\hspace*{.5ex}$#1$\hspace*{.5ex}};
    \node (Z) [right = 3ex of Y] {\hspace*{.5ex}$#3$};
    \path[->] (X) edge (Y);
    \path[<-] (Y) edge (Z);
  }%
}
\tikzset{
  square/nw/.style = {},
  square/ne/.style = {},
  square/se/.style = {},
  square/sw/.style = {},
  square/north/.style = {},
  square/south/.style = {},
  square/west/.style = {},
  square/east/.style = {},
  square/north/node/.style = {above},
  square/south/node/.style = {below},
  square/west/node/.style = {left},
  square/east/node/.style = {right},
}
\bool_new:N \l_jon_glue_west
\keys_define:nn { jon-tikz/diagram } {
  nw .tl_set:N = \l_jon_tikz_diagram_nw,
  sw .tl_set:N = \l_jon_tikz_diagram_sw,
  ne .tl_set:N = \l_jon_tikz_diagram_ne,
  se .tl_set:N = \l_jon_tikz_diagram_se,
  width .tl_set:N = \l_jon_tikz_diagram_width,
  height .tl_set:N = \l_jon_tikz_diagram_height,
  north .tl_set:N = \l_jon_tikz_diagram_north,
  south .tl_set:N = \l_jon_tikz_diagram_south,
  west .tl_set:N = \l_jon_tikz_diagram_west,
  east .tl_set:N = \l_jon_tikz_diagram_east,
  nw/style .code:n = {\tikzset{square/nw/.style = {#1}}},
  sw/style .code:n = {\tikzset{square/sw/.style = {#1}}},
  ne/style .code:n = {\tikzset{square/ne/.style = {#1}}},
  se/style .code:n = {\tikzset{square/se/.style = {#1}}},
  glue .choice:,
  glue / west .code:n = {\bool_set:Nn \l_jon_glue_west \c_true_bool},
  glue~target .tl_set:N = \l_jon_tikz_glue_target,
  north/style .code:n = {\tikzset{square/north/.style = {#1}}},
  north/node/style .code:n = {\tikzset{square/north/node/.style = {#1}}},
  south/style .code:n = {\tikzset{square/south/.style = {#1}}},
  south/node/style .code:n = {\tikzset{square/south/node/.style = {#1}}},
  west/style .code:n = {\tikzset{square/west/.style = {#1}}},
  west/node/style .code:n = {\tikzset{square/west/node/.style = {#1}}},
  east/style .code:n = {\tikzset{square/east/.style = {#1}}},
  east/node/style .code:n = {\tikzset{square/east/node/.style = {#1}}},
  draft .meta:n = {
    nw = {\__jon_tikz_diagram_fmt_placeholder:n {nw}},
    sw = {\__jon_tikz_diagram_fmt_placeholder:n {sw}},
    se = {\__jon_tikz_diagram_fmt_placeholder:n {se}},
    ne = {\__jon_tikz_diagram_fmt_placeholder:n {ne}},
    north = {\__jon_tikz_diagram_fmt_placeholder:n {north}},
    south = {\__jon_tikz_diagram_fmt_placeholder:n {south}},
    west = {\__jon_tikz_diagram_fmt_placeholder:n {west}},
    east = {\__jon_tikz_diagram_fmt_placeholder:n {east}},
  }
}
\tl_set:Nn \l_jon_tikz_diagram_width { 2cm }
\tl_set:Nn \l_jon_tikz_diagram_height { 2cm }
\cs_new:Npn \__jon_tikz_diagram_fmt_placeholder:n #1 {
  \texttt{\textcolor{red}{#1}}
}
\keys_set:nn { jon-tikz/diagram } {
  glue~target = {},
}
\cs_new:Nn \__jon_tikz_render_square:nn {
  \group_begin:
  \keys_set:nn {jon-tikz/diagram} {#2}
  \bool_if:nTF \l_jon_glue_west {
    \node (#1ne) [right = \l_jon_tikz_diagram_width~of~\l_jon_tikz_glue_target ne,square/ne] {$\l_jon_tikz_diagram_ne$};
    \node (#1se) [below = \l_jon_tikz_diagram_height~of~#1ne,square/se] {$\l_jon_tikz_diagram_se$};
    \path[->,square/north] (\l_jon_tikz_glue_target ne) edge node [square/north/node] {$\l_jon_tikz_diagram_north$} (#1ne);
    \path[->,square/east] (#1ne) edge node [square/east/node] {$\l_jon_tikz_diagram_east$} (#1se);
    \path[->,square/south] (\l_jon_tikz_glue_target se) edge node [square/south/node] {$\l_jon_tikz_diagram_south$} (#1se);
  } {
    \node (#1nw) [square/nw] {$\l_jon_tikz_diagram_nw$};
    \node (#1sw) [below = \l_jon_tikz_diagram_height~of~#1nw,square/sw] {$\l_jon_tikz_diagram_sw$};
    \path[->,square/west] (#1nw) edge node [square/west/node] {$\l_jon_tikz_diagram_west$} (#1sw);
    \node (#1ne) [right = \l_jon_tikz_diagram_width~of~#1nw,square/ne] {$\l_jon_tikz_diagram_ne$};
    \node (#1se) [below = \l_jon_tikz_diagram_height~of~#1ne,square/se] {$\l_jon_tikz_diagram_se$};
    \path[->,square/north] (#1nw) edge node [square/north/node] {$\l_jon_tikz_diagram_north$} (#1ne);
    \path[->,square/east] (#1ne) edge node [square/east/node] {$\l_jon_tikz_diagram_east$} (#1se);
    \path[->,square/south] (#1sw) edge node [square/south/node] {$\l_jon_tikz_diagram_south$} (#1se);
  }
  \group_end:
}
\NewDocumentCommand\SpliceDiagramSquare{D<>{}m}{
  \__jon_tikz_render_square:nn {#1} {#2}
}
\NewDocumentCommand\DiagramSquare{D<>{}O{}m}{
  \begin{tikzpicture}[diagram,#2,baseline=(#1sw.base)]
    \__jon_tikz_render_square:nn {#1} {#3}
  \end{tikzpicture}
}
\ExplSyntaxOff

% *** quiver ***
\usepackage{tikz-cd}
\usepackage{amssymb}
\tikzset{curve/.style={settings={#1},to path={(\tikztostart)
    .. controls ($(\tikztostart)!\pv{pos}!(\tikztotarget)!\pv{height}!270:(\tikztotarget)$)
    and ($(\tikztostart)!1-\pv{pos}!(\tikztotarget)!\pv{height}!270:(\tikztotarget)$)
    .. (\tikztotarget)\tikztonodes}},
    settings/.code={\tikzset{quiver/.cd,#1}
        \def\pv##1{\pgfkeysvalueof{/tikz/quiver/##1}}},
    quiver/.cd,pos/.initial=0.35,height/.initial=0}
\tikzset{tail reversed/.code={\pgfsetarrowsstart{tikzcd to}}}
\tikzset{2tail/.code={\pgfsetarrowsstart{Implies[reversed]}}}
\tikzset{2tail reversed/.code={\pgfsetarrowsstart{Implies}}}
\tikzset{no body/.style={/tikz/dash pattern=on 0 off 1mm}}
\begin{document}

\title{Вкладення Йонеди \\ в сімліціальній теорії типів}
\author{Данієль Гратцер, Джонатан Вайнбергер та Ульрік Бухгольц}

\begin{abstract}
  Ріль і Шульман \textcite{riehl:2017} ввели \emph{симпліціальну теорію типів} (\STT), варіант гомотопічної теорії
  типів, який мав на меті вивчати не лише теорію гомотопій, але і її злиття з теорією категорій:
  теорію $\prn{\infty,1}$-категорій. Хоча це і є технічно складним, маніпулювання $\infty$-категоріями в
  симпліціальній теорії типів часто є простішим, ніж робота зі звичайними категоріями, оскільки теорія типів
  обробляє нескінченні стеки когерентностей у фоновому режимі. Ми спираємося на нещодавню роботу
  \textcite{gratzer:2024}, що визначає $\prn{\infty,1}$-категорію $\infty$-групоїдів у \STT{}, щоб
  визначити категорії передпучків у межах \STT{} та систематично розвинути їхню теорію. Зокрема,
  ми будуємо вкладення Йонеди, доводимо універсальну властивість категорій передпучків, уточнюємо
  теорію спряженостей у \STT{}, вводимо теорію розширень Кана та доводимо теорему А
  Квіллена. На додаток до великого обсягу теорії категорій у \STT{}, ми пропонуємо вагомі
  докази того, що \STT{} може бути використана для отримання складних результатів у теорії $\infty$-категорій з
  набагато меншою складністю.
\end{abstract}

\maketitle

\begin{center}
  \emph{Присвячується світлій пам'яті Томаса Штрайхера (1958--2025)}
\end{center}


\section{Вступ}

Расселл \textcite{russell:1919} відомим чином описав два стилі формалізації математики як різницю
між \emph{крадіжкою} та \emph{чесною працею}. Обидва підходи можна побачити в сучасному використанні
залежної теорії типів. Чесна праця передбачає \emph{аналітичний} підхід: трактування типів як
базових об'єктів, еквівалентних множинам, а також визначення і міркування про такі об'єкти, як дійсні числа,
групи та топологічні простори, як це робиться зазвичай. Це те, що робиться, наприклад, у доведенні Coq
Теореми про непарний порядок~\parencite{gonthier:2013}. Більш швидкий шлях крадіжки передбачає
розгляд теорії типів як спеціалізованої \emph{синтетичної} мови для певного типу математичних
об'єктів і постулювання їх базових властивостей. Це звужує сферу застосування теорії типів, але, водночас,
робить доведення про ці конкретні об'єкти набагато лаконічнішими. Наприклад,
\emph{гомотопічна теорія типів} (\HoTT{})~\parencite{hottbook} постулює різні аксіоми, які гарантують, що типи
поводяться як простори (з точністю до гомотопії), уможливлюючи доведення теорем з алгебраїчної топології
без введення явного опису простору. Насправді, синтетичний підхід
менше схожий на крадіжку, ніж на позику; ви розплачуєтеся за налаштовану теорію типів
семантичною моделлю, яка інтерпретує типи як потрібні об'єкти і валідує додаткові
аксіоми.

У цій роботі ми використовуємо синтетичну методологію для застосування теорії типів до вивчення теорії
категорій. Зокрема, ми додаємо різні аксіоми до гомотопічної теорії типів, щоб побудувати систему,
де гасло \HoTT{} ``всі типи є просторами, а всі функції неперервними'' замінюється на
``(деякі) типи є ($\infty$-)категоріями, а всі функції є функторами''.\footnote{%
  У цій статті під $\infty$-категорією ми розуміємо $\prn{\infty,1}$-категорію.}
Це розширення
теорії типів називається \emph{симпліціальною} теорією типів (\STT{}) і було введено
\textcite{riehl:2017}.

Хоча знання $\infty$-категорій не є обов'язковим для використання нашої теорії, груба інтуїція щодо них
допомагає для розуміння \STT{}. Тому ми нагадаємо наступне нечітке визначення.
$\infty$-категорія $C$ -- це сукупність об'єктів з \emph{простором} стрілок між об'єктами $c$
та $d$, $\Hom{c}{d}$, а не множиною, оснащена неперервною операцією композиції та призначенням
тотожних стрілок. Вирішальним є те, що операція композиції має бути асоціативною та унітальною лише з
точністю до гомотопії, але з обмеженням, що ці гомотопії самі задовольняють закони \emph{когерентності} у
вигляді додаткових гомотопій, і так далі з когерентностями між когерентностями тощо. Як вільна
аналогія, подібно до того, як моноїдальна категорія послаблює моноїди, дозволяючи $\otimes$ бути асоціативним з
точністю до ізоморфізмів, що задовольняють певні рівняння когерентності, $\infty$-категорії послаблюють звичайні категорії
дозволяючи законам категорій виконуватися лише з точністю до (нескінченно когерентних) ізоморфізмів.

Прикметно, що практично кожна теорема, на яку можна сподіватися для звичайних категорій, виконується для
$\infty$-категорій.\footnote{Принаймні, як зауважив один з наших рецензентів, за умови правильного калібрування своїх сподівань.} Однак,
доведення є набагато складнішими, оскільки вони працюють з \emph{моделями} $\infty$-категорій (інструментами, що використовуються для
організації та управління вежею когерентностей~\parencite{bergner:2018}). Метою \STT{} є використання
теорії типів, щоб приховати когерентності від користувача і дозволити доведення, які є не складнішими, ніж
класичні міркування для $1$-категорій.

У цій роботі ми надаємо вагомі докази цієї гіпотези, розвиваючи велику частину
теорії категорій---кілька основних результатів \emph{Категорій для працюючого
  математика}~\parencite{maclane:1978}---виключно в межах \STT{}.

\subsection{Симпліціальна теорія типів}

Щоб побудувати теорію типів для синтетичної теорії категорій, можна було б сподіватися інтерпретувати теорію типів у
категорію категорій ($\infty$ або інакше), щоб гарантувати, що типи реалізують категорії. Однак,
категорія $\CAT$ малих категорій поводиться занадто погано, щоб утворити модель теорії типів
Мартіна-Льофа (\MLTT{}). Замість цього, \textcite{riehl:2017} розширюють $\CAT$ і вкладають її як рефлективну
підкатегорію в категорію ($\infty$-)передпучків на симплекс-категорії $\PSH{\SIMP}$, яка є достатньо
багатою для моделювання \HoTT{}. Тоді \STT{} аксіоматизує частину $\PSH{\SIMP}$, щоб ізолювати $\CAT$ як
рефлективний підуніверсум у межах теорії типів~\parencite{rijke:2020}.

Ми введемо повний набір доповнень у Розділі~\ref{sec:prelims} (зібраних у Додатку~\ref{sec:axioms}
для зручності), але найважливішим серед них є постульований \emph{інтервальний тип}
$\Int : \Uni[0]$. Ми додатково припускаємо, що $\Int$ є обмеженим лінійним порядком з кінцевими точками
$0,1 : \Int$. Інтуїтивно, $\Int$ має фіксувати категорію $\brc{0 \to 1}$---він так
інтерпретується в $\PSH{\SIMP}$---і ми можемо використовувати це для визначення та дослідження типу
\emph{синтетичних морфізмів} у довільному типі $X$: стрілка в $X$ відповідає звичайній
функції $\Int \to X$, причому обчислення в $0,1$ дає домен і кодомен. Наприклад,
тотожна стрілка в $x : X$ задається як $\lambda \_.\,x$.

Однак, подібно до того, як цільова модель $\PSH{\SIMP}$ є строго більшою за $\CAT$, не всі типи в \STT{}
достовірно моделюють категорії. Зокрема, хоча завжди можна побудувати тотожні
морфізми, не всі типи мають оператор композиції. Дивовижно, проте, що оператори композиції
є унікальними, коли вони існують, і їх існування для типу $X$ фіксується відносно коротким
твердженням (Визначення~\ref{def:prelims:category}). Маючи під рукою операцію композиції для $X$, ми можемо
визначити тип ізоморфізмів у $X$, і ми визначаємо категорію як тип, де (1)
операція композиції існує унікально з точністю до гомотопії, і (2) тип ізоморфізмів у $X$ еквівалентний типу ідентичності $=_X$.

\begin{rem}
  Цей останній пункт принципово залежить від \emph{відмови} від припущення унікальності доведень ідентичності, щоб ми
  випадково не заборонили будь-якій синтетичній категорії мати об'єкт з нетривіальним
  автоморфізмом. Однак, припускаючи, що ізоморфізми та доведення ідентичності збігаються, ми можемо
  використовувати підтримку теорією типів заміни рівного на рівне для безперешкодного перенесення доведень уздовж
  ізоморфізмів. Ось чому робота з \HoTT{}/інтенсіональною теорією типів при формулюванні синтетичної
  теорії категорій виявляється зручнішою, ніж екстенсіональна теорія типів, навіть якщо хтось не цікавиться
  $\infty$-категоріями.
\end{rem}


\subsection{Теорія категорій всередині \texorpdfstring{\STT{}}{STT}}

Хоча деякі нещодавні роботи досліджували \STT{} для її застосувань у мовах
програмування~\parencite{weaver:2020,gratzer:2024,weaver:phd}, більшість робіт із симпліціальної теорії типів були
зосереджені на доведенні результатів з теорії категорій всередині теорії
типів~\parencite{riehl:2017,riehl:2023,bardomiano:2021,weinberger:phd,buchholtz:2023,weinberger:twosided:2024,weinberger:sums:2024}.
З цією метою теорія спряженостей, дискретних розшарувань та розшарувань Гротендіка, а також (ко)границь були
введені та вивчені в межах симпліціальної теорії типів. Деякі з цих результатів, наприклад,
лема Йонеди для розшарувань~\parencite{riehl:2017}, були згодом
механізовані~\parencite{kudasov:2024}.

Донедавна, однак, у \STT{} не існувало замкнених типів, які б представляли нетривіальні
категорії. Як наслідок, хоча \textcite{riehl:2017} і наводить відмінне визначення спряженостей,
жодних прикладів навести не можна. У попередній роботі ми змінили це,
розширивши \STT{} для побудови $\Space$, категорії $\mathcal{S}$просторів, яка є гомотопічним аналогом $\SET$~\parencite{gratzer:2024}.
Об'єкти $\Space$ -- це елементи $\Uni[0]$, які кодують $\infty$-групоїди, а морфізми в
$\Space$ відповідають їх функціям. \emph{Там само} $\Space$ використовується як будівельний блок для
відновлення алгебраїчних категорій (груп, кілець), а також інших прикладів (частково впорядкованих множин,
симплекс-категорії тощо).

Наше розширення \STT{} використовувало різні \emph{модальності} поверх \HoTT{} для побудови $\Space$. Тут ми беремо $\Space$ повністю, але деякі з використаних нами модальностей все ще є критично важливими для формулювання природних теорем
у теорії категорій. Відповідно, у цій статті ми також працюємо в межах модального розширення \HoTT{}, заснованого на
\MTT{}~\parencite{gratzer:2020}.

\subsection{Внесок}
Ми переглядаємо базову теорію категорій у світлі побудови $\Space$ і показуємо, що
більшість класичних результатів, з якими стикаються в теорії категорій, тепер знаходяться в межах досяжності симпліціальної
теорії типів. Вперше ми показуємо, що \STT{} може бути використана для доведення життєво важливих теорем у
теорії $\infty$-категорій без звернення до складних моделей. Багато з цих теорем
(наприклад, цілком лояльні та суттєво сюр'єктивні функтори є еквівалентностями) прямо не згадують
$\Space$, але критично спираються на принципи міркування, які уможливлює $\Space$. Ми доводимо два
робочі результати для категорій передпучків $\PSH{C}$:
\begin{itemize}
\item Ми будуємо цілком лояльну функцію $\Yo : C \to \PSH{C}$.
\item Ми доводимо, що $\PSH{C}$ є ``вільним копоповненням $C$''.
\end{itemize}
Ключовою технічною інновацією для цього є категорія \emph{скручених стрілок} (twisted arrow), яку ми інтегруємо в
\STT{} як модальність. Після цього ми можемо вивести різні класичні результати, наприклад:
\begin{itemize}
\item що поточково оборотні відображення в $C \to D$ є оборотними;
\item що поточково ліві спряжені є лівими спряженими;
\item що (ко)границі обчислюються поточково в $C \to D$;
\item теорію та існування поточкових розширень Кана;
\item теорему А Квіллена;
\item власніть кодекартових розшарувань.
\end{itemize}

Синтетичний підхід дає лаконічне формулювання:\begin{restatable}{axiom}{intaxiom}
  \label{ax:int}
  Існує множина $\Int$, яка утворює обмежену дистрибутивну решітку $\prn{0,1,\lor,\land}$, таку що
  виконується $\Prod{i,j : \Int} i \le j \lor j \le i$.
\end{restatable}
Ми розглядаємо $\Int$ як \emph{спрямований інтервал}, і \textcite{riehl:2017} використовують його, щоб оснастити кожен тип поняттям синтетичного морфізму:
\begin{defi}
  Синтетичний морфізм $f : \Hom[X]{x}{y}$, де $x,y : X$, — це функція $f : \Int \to X$ разом із пропозиційними рівностями $f\,0 =_X x$ та $f\,1 =_X y$.
\end{defi}

\begin{rem}
  У \textcite{riehl:2017} синтетичний інтервал визначається як більш примітивна дедуктивна структура, і $\Hom{x}{y}$ використовує типи строгого розширення. Це дає більше дефініційних рівностей: $f\,0$ та $x$ збігалися б дефініційно, коли $f : \Hom{x}{y}$ (і аналогічно для $f\,1$ та $y$). Однак дедуктивний підхід не включає безпосередньо $\Int$ як звичайний тип, а його взаємодія з модальностями (Розділ~\ref{sec:prelims:modal-stt}) є складною. З цих причин ми працюємо з простішим, але менш строгим визначенням $\Hom{x}{y}$. У системі, де доступні обидва підходи, ці два поняття є еквівалентними~\parencite{buchholtz:2023}.
\end{rem}

Будь-яка функція $f : X \to Y$ автоматично має дію на синтетичні морфізми $\alpha : \Int \to X$ через пост-композицію $f \circ \alpha : \Int \to Y$. У цьому випадку ми часто пишемо $f\prn{\alpha}$.

З $\Int$ ми негайно отримуємо $n$-куби $\Int^n$, а з них ми можемо виділити симплекси $\Delta^n$, границі $\partial \Delta^n$ та роги $\Lambda_k^n$. Зокрема, $\Delta^2 \to X$ представляє $2$-клітину в $X$, що свідчить про композицію двох стрілок, а $\Lambda^2_1 \to X$ представляє пару компонованих стрілок (без композиції). Нижче ми наведемо визначення цих типів:
\begin{mathparpagebreakable}
  \Delta^n = \Compr{\prn{i_1, \dots, i_n} : \Int^n}{i_1 \ge i_2 \ge \dots \ge i_n}
  \and
  \Lambda^2_1 = \Compr{\prn{i,j} : \Int^2}{i = 1 \lor j = 0}
\end{mathparpagebreakable}

\begin{nota}
  Ми пишемо $i : \Delta^n$ ($0 \le i \le n$) як скорочення для послідовності з $i$ копій $1$, за якими йдуть $0$: $\prn{1, 1, \dots, 0, \dots}$.
\end{nota}

\begin{figure}
  \begin{tikzpicture}[diagram]
    \node (Ax00) {\tiny $0$};
    \node[right = 0.75cm of Ax00] (Ax10) {\tiny $1$};
    \node[below = 0.75cm of Ax00] (Ax01) {\tiny $1$};
    \path[->] (Ax00) edge node[above] {\tiny $i$} (Ax10);
    \path[->] (Ax00) edge node[left] {\tiny $j$} (Ax01);
    %
    %
    \SpliceDiagramSquare{
      nw/style = {below right = -0.3cm and 1.3cm of Ax01},
      height = 1cm,
      width = 1cm,
      nw = {\bullet},
      ne = {\bullet},
      sw = {\bullet},
      se = {\bullet},
    }
    \fill[pattern=north west lines]
    ([xshift=0.1cm,yshift=-0.1cm]nw.center) --
    ([xshift=-0.1cm,yshift=-0.1cm]ne.center) --
    ([xshift=-0.1cm,yshift=0.1cm]se.center) --
    ([xshift=0.1cm,yshift=0.1cm]sw.center) -- cycle;
    %
    \node[right = 1.5cm of ne] (T0) {$\bullet$};
    \node[right = 1cm of T0] (T1) {$\bullet$};
    \node[below = 1cm of T1] (T2) {$\bullet$};
    \path[->] (T0) edge (T1);
    \path[->] (T0) edge (T2);
    \path[->] (T1) edge (T2);
    \fill[pattern=north west lines]
    ([xshift=0.1cm,yshift=-0.1cm]T0.center) --
    ([xshift=-0.1cm,yshift=-0.1cm]T1.center) --
    ([xshift=-0.1cm,yshift=0.1cm]T2.center) -- cycle;
    %
    \node[right = 1.5cm of T1] (H0) {$\bullet$};
    \node[right = 1cm of H0] (H1) {$\bullet$};
    \node[below = 1cm of H1] (H2) {$\bullet$};
    \path[->] (H0) edge (H1);
    \path[->] (H1) edge (H2);
  \end{tikzpicture}
  \caption{Візуалізація $\Int^2$, $\Delta^2$ та $\Lambda^2_1$.}
  \label{fig:prelims:viz}
\end{figure}

Відображення $f : \Delta^2 \to X$ називається свідченням того, що композиція $f\prn{-,0}$, за якою йде $f\prn{1,-}$, є $\lambda i.\,f\prn{i,i}$. Ми наголошуємо, що це — \emph{дані}; може бути багато різних $f$, що свідчать про одну і ту саму композицію, оскільки $X$ може мати багато нееквівалентних $2$-клітин з однією і тією ж границею. З цієї ж причини, однак, не завжди пара компонованих морфізмів $\Lambda^2_1 \to X$ поширюється до даних композиції $\Delta^2 \to X$. Це саме тому, що не кожен тип у \STT{} можна розглядати як категорію; хоча ми визначили $\Hom[X]{x}{y}$ для кожного $X$, не існує апріорного способу \emph{компонування} цих морфізмів. Прекатегорії — це типи, для яких існують усі композиції:
\begin{defi}
  \emph{Прекатегорія} — це тип $X$, що задовольняє умову Сегала:
  вкладення $\Lambda^2_1 \to \Delta^2$ індукує еквівалентність $\IsEquiv\prn{X^{\Delta^2} \to X^{\Lambda^2_1}}$.
\end{defi}
Грубо кажучи, умова Сегала гарантує, що кожна пара компонованих морфізмів в $X$ поширюється (унікально) до $2$-клітини, що свідчить про їх композицію, і, зокрема, існує індукована функція композиції $\Hom{x}{y} \times \Hom{y}{z} \to \Hom{x}{z}$. Унікальність автоматично гарантує, що ця операція є асоціативною та унітальною. Визначення категорії уточнює це дещо далі. У прекатегорії $X$ ми можемо визначити тип ізоморфізмів $x \cong y$ між $x,y : X$, і тому існує два потенційно різних типи свідчень того, що $x$ та $y$ є ідентичними: $x =_X y$ та $x \cong_X y$. Категорія — це прекатегорія, для якої ці два типи є канонічно еквівалентними.
\begin{defi}
  $\alpha : \Hom{x}{y}$ є ізоморфізмом ($\IsIso\prn{\alpha}$), якщо існують $\beta_0,\beta_1 : \Hom{y}{x}$, такі що $\beta_0 \circ f = \ArrId{}$, $f \circ \beta_1 = \ArrId{}$.\footnote{Це саме HoTT-еквівалентність, але перенесена на синтетичні морфізми.} Ми пишемо $\Iso{x}{y}$ або $x \cong y$ для підтипу ізоморфізмів.
\end{defi}
\begin{defi}
  \label{def:prelims:category}
    Прекатегорія $C$ є \emph{категорією}, якщо вона задовольняє умову Резка:
    \[
      \Prod{\DeclVar{x,y}{}{C}} \IsEquiv\prn{\DeclVar{\mathsf{idtoiso}}{}{\prn{x=y} \to \Iso{x}{y}}}
    \]
    де $\mathsf{idtoiso}(\mathsf{refl}) := \ArrId{}$.
    Якщо кожен морфізм у $C$ є ізоморфізмом, то $C$ є \emph{групоїдом}.
\end{defi}
\begin{exa}
  $\Int,\Delta^n,\Int^n$ є категоріями~\parencite{gratzer:2024}.
\end{exa}

\section{Модальна та симпліціальна теорія типів}
\label{sec:prelims}

У цій статті ми сприймаємо \STT{} значною мірою як належне і зосереджуємося на роботі всередині теорії. Однак, щоб зробити цю статтю більш самодостатньою, ми присвячуємо цей розділ ретельному поясненню нових конструкцій модальної гомотопічної теорії типів та аксіом, що її доповнюють і утворюють симпліціальну теорію типів.

\subsection{Гомотопічна теорія типів}
Ми починаємо з нагадування основних понять і позначень гомотопічної теорії типів, які ми використовуємо в цій статті. Канонічним джерелом є книга \HoTT{}~\parencite{hottbook}. Ми працюємо в межах інтенсіональної теорії типів Мартіна-Льофа і зауважимо, як \HoTT{} розширює її.

\begin{nota}
  Ми пишемо $a =_A b$ для типу ідентичності (часто опускаючи $A$). Дано $p : a =_A b$ та
  $B : A \to \Uni$, ми пишемо $p_!$ для відображення $B\prn{a} \to B\prn{b}$.
\end{nota}

\begin{defi}
  Ми кажемо, що функція $f : A \to B$ є еквівалентністю, якщо $f$ допускає як лівий, так і правий
  обернений:
  \[
    \IsEquiv\prn{f} =
    \Sum{g,h : B \to A} \prn{g \circ f = \ArrId{}} \times \prn{f \circ h = \ArrId{}}
  \]
  Ми пишемо $A \Equiv B$ для суми $\Sum{f : A \to B} \IsEquiv\prn{f}$.
\end{defi}
\HoTT{} є розширенням інтенсіональної теорії типів з ієрархією універсумів, що задовольняють аксіому унівалентності:
\[
  \Con{univ}_i
  :
  \Prod{A, B : \Uni[i]} \IsEquiv\prn{
    \lambda p.\,\prn{p_!, \cdots} : A =_{\Uni[i]} B \to A \Equiv B
  }
\]
Ми будемо опускати $i$ в $\Con{univ}_i$ та $\Uni[i]$ і ігнорувати питання розміру, якщо вони не є суттєвими. Унівалентність породжує велику кількість шляхів у $\Uni$, які відрізняються від $\Refl$. Нас часто цікавлять типи, які є тривіальними, мають лише тривіальні шляхи або тривіальні шляхи між шляхами і так далі. Ці умови організовані в сімейство предикатів, що називаються рівнем усічення ($-2,-1,0,\ldots$) типу. Ми будемо використовувати лише перші три рівні, стверджуючи, що тип є стягуваним, (гомотопічною) пропозицією або множиною:
\begin{mathparpagebreakable}
  \IsContr\prn{A} = \Sum{a : A} \Prod{b : A} a = b
  \and
  \IsProp\prn{A} = \Prod{a,b : A} \IsContr\prn{a = b}
  \and
  \IsSet\prn{A} = \Prod{a,b : A} \IsProp\prn{a = b}
\end{mathparpagebreakable}

\begin{prop}[\textcite{shulman:2019}, див. також \textcite{riehl:2024}]
  \label{thm:prelims:hott-models}
  Усі теоретико-типові модельні топоси
  \textup(і, отже, $\infty$-топоси Гротендіка\textup) моделюють \HoTT{}.
\end{prop}

Ми також матимемо нагоду використовувати різні \emph{вищі індуктивні типи} (ВІТ). Семантика ВІТ є складною і безпосередньо не розглядається наведеним вище результатом~\parencite{lumsdaine:2019}. Зокрема, хоча \textcite{shulman:2019} показує, що вищезгадана модель підтримує всі вищі індуктивні типи, він не показує, що універсуми є строго замкненими щодо цих конструкцій. Хоча робота над отриманням цього результату триває, легко показати, що універсуми є \emph{слабко} замкненими щодо цих конструкцій. Наприклад, існує тип $D : \Uni[0]$, такий що $D \Equiv A \Pushout{C} B$ щоразу, коли $A,B,C : \Uni[0]$. Відповідно, ми будемо припускати, що наші універсуми замкнені щодо вищих індуктивних типів, хоча і лише з пропозиційними $\beta$-правилами.

\subsection{Симпліціальна теорія типів}
Маючи під рукою \HoTT{}, ми переходимо до симпліціальної теорії типів. Це розширення \HoTT{} кількома аксіомами, які дозволяють нам розглядати (певні) типи як $(\infty,1)$-категорії, що надалі будуть називатися просто категоріями. Ми, відповідно, опускатимемо префікс $(\infty,1)$- або $\infty$- скрізь. По-перше і найфундаментальніше, ми додаємо наступне:

\begin{restatable}{axiom}{intaxiom}
  \label{ax:int}
  Існує множина $\Int$, що утворює обмежену дистрибутивну ґратку $\prn{0,1,\lor,\land}$ таку,
  що $\Prod{i,j : \Int} i \le j \lor j \le i$ виконується.
\end{restatable}
Ми розглядаємо $\Int$ як \emph{напрямлений інтервал}, і \textcite{riehl:2017} використовують це, щоб надати кожному типу поняття синтетичного морфізму:
\begin{defi}
  Синтетичний морфізм $f : \Hom[X]{x}{y}$ між $x,y : X$ - це функція $f : \Int \to X$ разом
  з пропозиційними рівностями $f\,0 =_X x$ та $f\,1 =_X y$.
\end{defi}

\begin{rem}
  В \textcite{riehl:2017}, синтетичний інтервал визначається як більш примітивна судженнява структура
  і $\Hom{x}{y}$ використовує строгі типи розширення. Це дає більше визначених рівностей: $f\,0$ та
  $x$ співпадатимуть за означенням, коли $f : \Hom{x}{y}$ (і аналогічно для $f\,1$ та $y$). Проте, підхід
  суджень не дозволяє прямо включити $\Int$ як звичайний тип, і його взаємодії з модальностями
  (Розділ~\ref{sec:prelims:modal-stt}) є складними. З цих причин, ми працюємо з простішим, але менш
  строгим означенням $\Hom{x}{y}$. В системі, де обидва доступні, ці два поняття еквівалентні\parencite{buchholtz:2023}.
\end{rem}

Звісно, ми можемо визначити тотожний морфізм $\ArrId{x} : x \to x$ як $\lambda \_.\,x$. Більше того, кожна
функція $f : X \to Y$ автоматично має дію на синтетичні морфізми $\alpha : \Int \to X$ шляхом
пост-композиції $f \circ \alpha : \Int \to Y$. В такому випадку ми часто пишемо $f\prn{\alpha}$.

З $\Int$ ми негайно отримуємо $n$-куби $\Int^n$ і з них ми можемо виділити симплекси
$\Delta^n$, межі $\partial \Delta^n$ та роги $\Lambda_k^n$. Зокрема, $\Delta^2 \to X$
представляє $2$-комірку в $X$, що свідчить про композицію двох стрілок,
і $\Lambda^2_1 \to X$ представляє пару композиційних стрілок (без композиції). Ми наводимо означення
цих типів нижче:
\begin{mathparpagebreakable}
  \Delta^n = \Compr{\prn{i_1, \dots, i_n} : \Int^n}{i_1 \ge i_2 \ge \dots \ge i_n}
  \and
  \Lambda^2_1 = \Compr{\prn{i,j} : \Int^2}{i = 1 \lor j = 0}
\end{mathparpagebreakable}

\begin{nota}
  Ми пишемо $i : \Delta^n$ ($0 \le i \le n$) як скорочений запис для послідовності з $i$ копій $1$, за якими слідує $0$: $\prn{1, 1, \dots, 0, \dots}$.
\end{nota}

\begin{figure}
  \begin{tikzpicture}[diagram]
    \node (Ax00) {\tiny $0$};
    \node[right = 0.75cm of Ax00] (Ax10) {\tiny $1$};
    \node[below = 0.75cm of Ax00] (Ax01) {\tiny $1$};
    \path[->] (Ax00) edge node[above] {\tiny $i$} (Ax10);
    \path[->] (Ax00) edge node[left] {\tiny $j$} (Ax01);
    %
    %
    \SpliceDiagramSquare{
      nw/style = {below right = -0.3cm and 1.3cm of Ax01},
      height = 1cm,
      width = 1cm,
      nw = {\bullet},
      ne = {\bullet},
      sw = {\bullet},
      se = {\bullet},
    }
    \fill[pattern=north west lines]
    ([xshift=-0.1cm,yshift=0.1cm]nw.south east) --
    ([xshift=0.1cm,yshift=0.1cm]ne.south west) --
    ([xshift=0.1cm,yshift=-0.1cm]se.north west) --
    ([xshift=-0.1cm,yshift=-0.1cm]sw.north east) -- cycle;
    %
    \node[right = 1.5cm of ne] (T0) {$\bullet$};
    \node[right = 1cm of T0] (T1) {$\bullet$};
    \node[below = 1cm of T1] (T2) {$\bullet$};
    \path[->] (T0) edge (T1);
    \path[->] (T0) edge (T2);
    \path[->] (T1) edge (T2);
    \fill[pattern=north west lines]
    ([yshift=0.1cm,]T0.south east) --
    ([yshift=0.1cm,xshift=0.15cm]T1.south west) --
    ([xshift=0.15cm]T2.north west) -- cycle;
    %
    \node[right = 1.5cm of T1] (H0) {$\bullet$};
    \node[right = 1cm of H0] (H1) {$\bullet$};
    \node[below = 1cm of H1] (H2) {$\bullet$};
    \path[->] (H0) edge (H1);
    \path[->] (H1) edge (H2);
  \end{tikzpicture}
  \caption{Visualization of $\Int^2$, $\Delta^2$, and $\Lambda^2_1$.}
  \label{fig:prelims:viz}
\end{figure}

A map $f : \Delta^2 \to X$ is said to witness that the composite of $f\prn{-,0}$
followed by $f\prn{1,-}$ is $\lambda i.\,f\prn{i,i}$. We emphasize that this is \emph{data}; there
can be many distinct $f$'s witnessing the same composition as $X$ may have many non-equivalent
$2$-cells with the same boundary. By the same token however, it is not always the case that a pair of
composable morphisms $\Lambda^2_1 \to X$ extends to a composition datum $\Delta^2 \to X$. This is
precisely because not every type in \STT{} can be regarded as a category; even though we have
defined $\Hom[X]{x}{y}$ for every $X$, there is no \emph{a priori} way of \emph{composing} these morphisms.
Precategories are types for which all composites exist:
\begin{defi}
  A \emph{precategory} is a type $X$ satisfying the Segal condition:
  the inclusion $\Lambda^2_1 \to \Delta^2$ induces an equivalence $\IsEquiv\prn{X^{\Delta^2} \to X^{\Lambda^2_1}}$.
\end{defi}
Roughly, the Segal condition ensures that every pair of composable morphisms in $X$ extends
(uniquely) to a $2$-cell witnessing their composition and, in particular, there is an induced
composition function $\Hom{x}{y} \times \Hom{y}{z} \to \Hom{x}{z}$. Uniqueness automatically ensures
that this operation is associative and unital. The definition of a category refines this
slightly. In a precategory $X$ we are able to define the type of isomorphisms $x \cong y$ between
$x,y : X$ and so there are two potentially distinct types of evidence for $x$ and $y$ being
identical: $x =_X y$ and $x \cong_X y$. A category is a precategory for which these two types are
canonical equivalent.
\begin{defi}
  $\alpha : \Hom{x}{y}$ is an isomorphism ($\IsIso\prn{\alpha}$) if there exist
  $\beta_0,\beta_1 : \Hom{y}{x}$ such that $\beta_0 \circ f =
  \ArrId{}$,$f \circ \beta_1 = \ArrId{}$.\footnote{This is precisely the \HoTT{} equivalence but recast
    into synthetic morphisms.} We write $\Iso{x}{y}$ or $x \cong y$ for the subtype of isomorphisms.
\end{defi}
\begin{defi}
  \label{def:prelims:category}
    A precategory $C$ is a \emph{category} if it satisfies the Rezk condition:
    \[
      \Prod{\DeclVar{x,y}{}{C}} \IsEquiv\prn{\DeclVar{\mathsf{idtoiso}}{}{\prn{x=y} \to \Iso{x}{y}}}
    \]
    where $\mathsf{idtoiso}(\mathsf{refl}) := \ArrId{}$.
    If every morphism in $C$ is an isomorphism, then $C$ is a \emph{groupoid}.
\end{defi}
\begin{exa}
  $\Int,\Delta^n,\Int^n$ are all categories~\parencite{gratzer:2024}.
\end{exa}

\begin{lem}
  $C$ is a groupoid if and only if $\IsEquiv\prn{\DeclVar{C^{!_\Int}}{}{C \to C^\Int}}$
  \textup($C$ is \emph{$\Int$-null} \textup{\cite{rijke:2020})}.
\end{lem}

\textcite{riehl:2017} develop the basic theory of these synthetic categories.
As noted above, every
function has an action on morphisms and \opcit{} shows that this action preserves compositions and
identities and therefore defines a functor. They also show that $C \to D$ is then a category
whenever $D$ is, and that synthetic morphisms $\Hom[D^C]{f}{g}$ are precisely natural
transformations. One can reformulate various classical categorical notions rather directly:
\begin{defiC}[\parencite{bardomiano:2021}]
  A natural transformation $\alpha : \Hom[C^I]{\Const\prn{c}}{F}$ witnesses $c$ as the limit of
  $F : C^I$ if $\alpha$ induces an equivalence $\Hom{c'}{c} \Equiv \Hom{\Const\prn{c'}}{F}$ for all $c'$.
\end{defiC}

\begin{defi}
  An adjunction between two categories $C,D$ consists of a pair of functions $f : C \to D$ and
  $g : D \to C$ with a natural isomorphism
  $\iota : \Prod{c,d} \Hom{f\prn{c}}{d} \Equiv \Hom{c}{g\prn{d}}$.
\end{defi}

While we have given a few examples of categories above, a notable type that is \emph{not} category
is the universe $\Uni$. Maps $A : \Int \to \Uni$ are too unstructured to compose and, in particular,
correspond neither to functions $A\prn{0} \to A\prn{1}$ nor $A\prn{1} \to A\prn{0}$ (consider
$\lambda i.\, i = 0$ or $\lambda i.\, i = 1$). In Section~\ref{sec:prelims:modal-stt}, we shall discuss the
subuniverse $\Space$ constructed previously~\parencite{gratzer:2024}, which is a category of groupoids whose
morphisms correspond to functions.  To properly situate this definition, we recall what it means for
$X : A \to \Uni$ to be covariant~\parencite{riehl:2017}, giving an assignment from morphisms
$\Hom{a_0}{a_1}$ to functions $X\prn{a_0} \to X\prn{a_1}$.

\begin{nota}
  Given $X : A \to \Uni$ we write $\TotalType{X}$ for $\Sum{a : A} X\prn{a}$.
\end{nota}

\begin{defi}
  A family $X : A \to \Uni$ is \emph{covariant}
  if for every $a : \Hom{a_0}{a_1}$ and $x_0 : X\prn{a_0}$,
  the following is contractible:
  \[
    \Con{Lift}\prn{a,x_0} =
    \Sum{x_1 : X\prn{a_1}}
    \Sum{x : \Hom{\prn{a_0,x_0}}{\prn{a_1,x_1}}}
    \Proj[1]\prn{x} =_{\Hom{a_0}{a_1}} a
  \]
  Here, $x$ is a morphism in $\TotalType{X}$. We further say the projection
  $\Sum{a : A} X\prn{a} \to A$ is covariant when $X$ is. For a general map $\pi : X \to A$ we write
  $X_a$ for $\Sum{x : X} \pi\prn{x} = a$ and say $\pi$ is covariant when $\lambda a.\,X_a$ is.
\end{defi}

Since $\Con{Lift}\prn{a,x_0}$ is contractible it has an inhabitant $x_1$. This yields a function
$a,x_0 \mapsto x_1$ which defines $a_! : X\prn{a_0} \to X\prn{a_1}$. The contractibility of
$\Con{Lift}\prn{a,x_0}$ ensures that these functions compose correctly, \etc{}

\begin{lem}
  \label{lem:prelims:covariant}
  A family $\DeclVar{X}{}{A \to \UU}$ is covariant if and only if the map
  $\bar{\pi} := \lambda p.\,\prn{p\prn{0}, \Proj[1] \circ p} : \TotalType X^\Int \to \TotalType X \times_A A^\Int$
  is an equivalence.
\end{lem}

In Sections~\ref{sec:adj:pw} and \ref{sec:kan:quillen-a}, we shall briefly use a weakening of covariance:
\begin{defi}
  A family $X : A \to \Uni$ is cocartesian if
  ${\TotalType{X}}^\Int \to A^\Int \times_{A^{\brc{1}}} {\TotalType{X}}^{\brc{1}}$ is a right
  adjoint $\ell \Adjoint \bar{\pi}$ such that $\bar{\pi} \circ \ell = \ArrId{}$.
\end{defi}
One can give an equivalent characterization in terms of cocartesian morphisms and show that \eg{},
every morphism in $D$ can be factored as a cocartesian morphism followed by a vertical morphism:
\begin{thm}[\textcite{buchholtz:2023}]
  If a map $\pi : D \to C$ is cocartesian then for every $c : \Int \to C$, $x_0 : C_{c(0)}$, the
  category $\Con{Lift}\prn{c,x_0}$ has an initial object.
\end{thm}
Dually, one can consider contravariant and cartesian families and fibrations.

Finally, we note that since categories and groupoids are defined by certain
\emph{orthogonality} conditions, by \textcite{rijke:2020} they define reflective subuniverses.
\begin{prop}
  There are idempotent monads $\Rezkify,\Grpdify$ such that, \eg{}, $\Rezkify X$ is a category and
  $C^{\Rezkify X} \Equiv C^X$ when $C$ is a category.
\end{prop}

\begin{prop}[\textcite{riehl:2017}]
  \label{thm:prelims:stt-models}
  When Theorem~\ref{thm:prelims:hott-models} is specialized to simplicial spaces
  \textup($\PSH{\SIMP}$\textup), the
  resulting model validates Axiom~\ref{ax:int} and in this model categories are realized by
  $\infty$-categories \textup(modeled by complete Segal spaces\textup)
  and groupoids by $\infty$-groupoids.
\end{prop}


\subsection{Modal homotopy type theory}
Many theorems in category theory require the ability to quantify over the objects in a category,
\eg{}, ``if $\alpha : F \to G$ is a natural transformation of functors $\CC \to \DD$ and each
$\alpha_c$ is invertible, then $\alpha$ is invertible''. A version of this is proven by
\textcite{riehl:2017}:
$\prn{\Prod{c : C} \IsIso\prn{\lambda i.\,\alpha\,i\,c}} \to \IsIso\prn{\alpha}$, but this is subtly
different as we discuss below. In fact, as it stands we cannot directly capture the classical
statement in \STT{}.

To understand the divergence between the \STT{} and classical results, note that by working
internally to type theory when proving $\Prod{c : C} \IsIso\prn{\lambda i.\,\alpha\,i\,c}$ we
cannot assume that $c$ is just an object in $C$: since it is an arbitrary element, we have to assume
it is constructed in an arbitrary context which might contain, \eg{}, a copy of $\Int$ such that $c$
represents a synthetic morphism. In fact, if we unfold the above type into the model we find that
constructing $\Prod{c : C} \IsIso\prn{\lambda i.\,\alpha\,i\,c}$ already entails proving that the
chosen inverses are natural. A great deal of the power of simplicial type theory comes from this
implicit naturality, but it makes this particular result weaker. After all, its purpose in standard
category theory was that in this particular situation, \emph{a priori} unnatural choices of inverses
will automatically be natural. Moreover, we shall encounter theorems that are simply false when
naively translated in this way.

Accordingly, to make \STT{} practical we must extend it with \emph{modalities}: unary type
constructors distinguished by their failure to respect substitution or apply in arbitrary
contexts. For instance, we shall eventually equip \STT{} with a modality $\Modify[\GM]{-}$ which
discards all non-invertible synthetic morphisms from a type to produce its \emph{core},
which we then
use to faithfully encode pointwise invertibility (see Example~\ref{ex:prelim:pw-naturality}).

A complete reference to the modal type theory we use---\MTT{}~\parencite{gratzer:2020}---is given by
\textcite{gratzer:phd} and we record formal rules in Section~\ref{sec:mtt-rules}. Fortunately, the rules
for, \eg{}, $\Sum{}$-types are unaffected by the addition of modalities. Accordingly, for brevity we only recall the new rules which must be added to \MLTT{} to extend \HoTT{} with modalities
\`a la \MTT{}.

\MTT{} is parameterized by a \emph{mode theory}: a strict $2$-category describing the collection
of modalities (the morphisms) available along with the natural transformations between them (the
$2$-cells).  We use $\mu,\nu,\xi$ to range over modalities. In the case of simplicial type theory,
our mode theory will have only one object along with a handful of generating modalities and
$2$-cells. There are four generating modalities $\GM,\SM,\OM,\TM$ subject to the following
equations:
\begin{mathparpagebreakable}
  \GM = \GM \circ \GM = \GM \circ \SM = \GM \circ \OM = \TM \circ \GM %
  \and
  \SM = \SM \circ \SM = \SM \circ \GM = \SM \circ \OM %
  \and
  \OM \circ \OM = \ArrId{}
\end{mathparpagebreakable}

We further require the following generating 2-cells:
\begin{mathparpagebreakable}
  \epsilon : \GM \to \ArrId{} \qquad \zeta : \ArrId{} \to \SM
  \and
  \tau : \TM \cong \TM \circ \OM \quad\text{(with $\tau^{-1})$}
  \and
  \pi_0^\TM : \TM \to \OM
  \and
  \pi_1^\TM : \TM \to \ArrId{}
\end{mathparpagebreakable}

These 2-cells are likewise subject to a number of equations. For $\epsilon$ and $\zeta$, we
require $\zeta \Whisker \SM = \SM \Whisker \zeta = \ArrId{}$ viewing all of these as 2-cells
$\SM \to \SM$ and $\GM \Whisker \zeta = \ArrId{} : \GM \to \GM$ (using $\Whisker$ to denote
whiskering).  We require the dual equations on $\epsilon$. For the remaining four 2-cells, we
require that the following diagrams commute:
\[
  \begin{tikzpicture}[diagram]
    \node (T) {$\TM$};
    \node [below = of T] (OT) {$\TM \circ \OM$};
    \node [left = 2.5cm of OT] (OC) {$\OM$};
    \node [right = 2.5cm of OT] (C) {$\ArrId{}$};
    \path[->] (T) edge node[right] {$\tau$} (OT);
    \path[->] (OT) edge node[below] {$\pi^\TM_0 \Whisker \OM$} (C);
    \path[->] (OT) edge node[below] {$\pi^\TM_1 \Whisker \OM$} (OC);
    \path[->] (T) edge node[above] {$\pi^\TM_1$} (C);
    \path[->] (T) edge node[above] {$\pi^\TM_0$} (OC);
  \end{tikzpicture}
  \qquad
  \begin{tikzpicture}[diagram]
    \node (T) {$\GM$};
    \node [below = of T] (OT) {$\TM$};
    \node [left = 2.5cm of OT] (OC) {$\OM$};
    \node [right = 2.5cm of OT] (C) {$\ArrId{}$};
    \path[->] (T) edge node[near end,right] {$\TM\Whisker\epsilon$} (OT);
    \path[->] (OT) edge node[below] {$\pi^\TM_0$} (C);
    \path[->] (OT) edge node[below] {$\pi^\TM_1$} (OC);
    \path[->] (T) edge node[above] {$\epsilon$} (C);
    \path[->] (T) edge node[above=5pt] {$\epsilon\Whisker\OM$} (OC);
  \end{tikzpicture}
\]

Each morphism $\mu$ in the mode theory induces a modal type $\Modify[\mu]{-}$. We will describe the
rules for these modal types in a moment, but first we give some idea of what they are intended to
denote. For now this is merely intuition, though the axioms and model described in
Section~\ref{sec:prelims:modal-stt} will make it so. As already mentioned, $\Modify[\GM]{-}$ removes all
non-identity synthetic morphisms from a type. $\Modify[\SM]{-}$ is the right adjoint to this
operation and so it discards all non-identity morphisms but then freely adds all morphisms so that
an $n$-simplex $\Delta^n \to \Modify[\SM]{X}$ is exactly a collection of $n$ points in
$X$. \footnote{Note that while $\Modify[\GM]{X}$  will always be an $\infty$-category, in fact an
$\infty$-groupoid, the same is not true of $\Modify[\SM]{X}$. In particular, $\Modify[\SM]{X}$
will hardly ever satisfy the Rezk condition.} Next, $\Modify[\OM]{-}$ sends a type to its opposite
and, in particular, reverses the directions of all synthetic morphisms. Finally, $\Modify[\TM]{-}$
sends a type to its corresponding type of \emph{twisted arrows}; we shall analyze it in more depth
in Section~\ref{sec:tw}.

The formation rule for $\Modify[\mu]{-}$ is complex: the entire point of modalities is that
$\Gamma \vdash A$ does not imply $\Gamma \vdash \Modify{A}$. Instead, \MTT{} introduces a novel form
of context operation which acts like a ``left adjoint'' to $\Modify[\mu]{-}$:
\begin{mathparpagebreakable}
  \inferrule{
    \IsCx{\Gamma}
  }{
    \IsCx{\LockCxV{\Gamma}}
  }
  \and
  \inferrule{
    \IsTy[\LockCxV{\Gamma}]{A}
  }{
    \IsTy{\Modify{A}}
  }
  \and
  \inferrule{
    \IsTm[\LockCxV{\Gamma}]{a}{A}
  }{
    \IsTm{\MkMod{a}}{\Modify{A}}
  }
\end{mathparpagebreakable}
We refer to $\brc{\mu}$ as a \emph{modal restriction}. It is helpful to compare $\Modify[\mu]{A}$
with dependent products and, therefore, to see $\LockCxV{-}$ as extending the context by something
akin to a substructural ``$\mu$ variable''~\parencite{bahr:2017,birkedal:2020}. The real force of
modalities comes through how these $\brc{\mu}$s interact with variables. In particular, it is not the
case that $\Gamma, x : A, \brc{\mu} \vdash x : A$; since $\LockCxV{-}$ is intended to model a left
adjoint, we cannot generally assume that there is a weakening substitution
$\LockCxV{\Gamma} \to \Gamma$. Instead, we alter the rule extending a context with a variable so
that each variable is annotated with a modality:
\begin{mathparpagebreakable}
  \inferrule{
    \IsCx{\Gamma}
    \\
    \IsTy[\LockCxV{\Gamma}]{A}
  }{
    \IsCx{\MECxV{\Gamma}{x}{A}}
  }
  \and
  \inferrule{
    \alpha : \mu \to \Con{mods}\prn{\Gamma_1}
  }{
    \IsTm[\Gamma_0, \DeclVar{x}{\mu}{A}, \Gamma_1]{x^\alpha}{A^\alpha}
  }
\end{mathparpagebreakable}
In the above, $x^\alpha$ is the new form of variable rule while $A^\alpha$ is an admissible
operation on the syntax which traverses the term $A$ and appropriately updates all free variables
$y^\beta$ occurring within $A$ and modifying $\beta$ appropriately using $\alpha$. In particular, if
$A$ is closed then $A^\alpha = A$. In the formal syntax, both $x^\alpha$ and $A^\alpha$ are realized
by form of substitution, see \textcite{gratzer:mtt-journal:2021} for further details.

The original context extension is given by taking $\mu = \ArrId{}$. In the second rule,
$\Con{mods}\prn{\Gamma_1}$ is the composite $\nu_0 \circ \nu_1 \circ \cdots$ of all the
$\brc{\nu_i}$s occurring in $\Gamma_1$ (and is $\ArrId{}$ if there are no such occurrences). In
other words, a variable with annotation $\mu$ can be used precisely when it occurs behind a series
of modal restrictions for which there is a 2-cell navigating from $\mu$ to this composite. It is
therefore in the variable rule where the 2-cells comes into play.
\begin{lem}
  If $\IsTy[\Gamma, \DeclVar{x}{\mu}{A}]{B}$, $\IsTm[\Gamma, \DeclVar{x}{\mu}{A}]{b}{B}$,
  and $\IsTm[\LockCxV{\Gamma}]{a}{A}$, then $\IsTy{\Sb{B}{a/x}}$ and $\IsTm{\Sb{b}{a/x}}{\Sb{B}{a/x}}$.
\end{lem}

The final piece of the puzzle is the elimination rule for modalities. Roughly, this rule says that
modal annotations are equivalent to modal types ``from the perspective of a type'', \ie{}, that giving
an element in context $\MECxV{\Gamma}{x}[\nu]{\Modify[\mu]{A}}$ is the same as giving one in
$\MECxV{\Gamma}{x}[\nu\circ\mu]{A}$. This concretely amounts to the following pattern-matching rule
which allows us to assume that $\DeclVar{x}{\nu}{\Modify[\mu]{A}}$ is of the form $\MkMod{y}$ where
$\DeclVar{y}{\nu\circ\mu}{A}$:
\begin{mathparpagebreakable}
  \inferrule{
    \IsTy[\MECxV{\Gamma}{x}[\nu]{\Modify{A}}]{B}
    \\
    \IsTm[\MECxV{\Gamma}{y}[\nu\circ\mu]{A}]{b}{\Sb{B}{\MkMod[\mu]{y}/x}}
    \\
    \IsTm[\LockCxV{\Gamma}[\nu]]{a}{\Modify{A}}
  }{
    \IsTm{\LetMod{a}[y]{b}}{\Sb{B}{a/x}}
  }
  \and
  \inferrule{
    \IsTy[\MECxV{\Gamma}{x}[\nu]{\Modify{A}}]{B}
    \\
    \IsTm[\MECxV{\Gamma}{y}[\nu\circ\mu]{A}]{b}{\Sb{B}{\MkMod[\mu]{y}/x}}
    \\
    \IsTm[\LockCxV{\Gamma}[\nu\circ\mu]]{a}{A}
  }{
    \EqTm{\prn{\LetMod{\MkMod{a}}[y]{b}}}{\Sb{b}{a/y}}{\Sb{B}{\MkMod{a}/x}}
  }
\end{mathparpagebreakable}

While these rules account for all of the necessary extensions to handle modal types, we avail
ourselves of a convenience feature as well, modal $\Prod{}$-types:
\begin{mathparpagebreakable}
  \inferrule{
    \IsTm[\MECxV{\Gamma}{x}{A}]{b}{B}
  }{
    \IsTm{\lambda x.b}{\Prod{\DeclVar{x}{\mu}{A}} B}
  }
  \and
  \inferrule{
    \IsTm{f}{\Prod{\DeclVar{x}{\mu}{A}} B}
    \quad
    \IsTm[\LockCxV{\Gamma}]{a}{A}
  }{
    \IsTm{f\prn{a}}{\Sb{B}{a/x}}
  }
\end{mathparpagebreakable}

\begin{nota}
  ``If $\DeclVar{c}{\GM}{C}$, then $\Phi\prn{c}$'' signifies
  $\Prod{\DeclVar{c}{\GM}{C}} \Phi\prn{c}$.
\end{nota}

\begin{exa}
  \label{ex:prelim:pw-naturality}
  A faithful translation of ``pointwise invertibility implies invertibility'' where
  $\DeclVar{C,D}{\GM}{\Uni}$ and $\alpha : C \times \Int \to D$ is
  $\prn{\Prod{\DeclVar{c}{\GM}{C}} \IsIso\prn{\lambda i.\,\alpha\,i\,c}}
  \to \IsIso\prn{\alpha}$
\end{exa}

Immediately from these rules, we may prove the following:
\begin{prop}[\textcite{gratzer:2020}]
  \label{thm:prelims:mtt-facts}
  \leavevmode
  \begin{itemize}
  \item $\Modify[\mu]{-}$ commutes with $\Sum{}$ and $\ObjTerm{}$
  \item $\Con{comp} : \Modify[\mu]{\Modify[\nu]{-}} \Equiv \Modify[\mu\circ\nu]{-}$ and
    $\Modify[\ArrId{}]{-} \Equiv \ArrId{}$
  \item If $\alpha : \mu \to \nu$, then there is a map $\Coe^{\alpha} : \Modify[\mu]{{-}} \to \Modify[\nu]{-^\alpha}$.
  \item $\Con{transp} : \Modify[\GM]{\Modify[\GM]{A} \to B} \Equiv \Modify[\GM]{A \to \Modify[\SM]{B}}$
  \item $\Con{transp} : \Modify[\GM]{\Modify[\OM]{A} \to B} \Equiv \Modify[\GM]{A \to \Modify[\OM]{B}}$
  \end{itemize}
  The first point yields a function $\prn{\circledast} : \Modify{A \to B} \to \Modify{A} \to \Modify{B}$.
\end{prop}

When it will not confusion, we will suppress the equivalences
$\Modify[\mu]{\Modify[\nu]{A}} \Equiv \Modify[\mu\circ\nu]{A}$ and $\Modify[\ArrId{}]{A} \Equiv
A$. Furthermore, as there is no ambiguity, we suppress $\epsilon$ (and its whiskerings) and simply
write $x$ instead of $x^\epsilon$ and similarly for $\zeta : \ArrId{} \to \SM$. Consequently, if
$\DeclVar{A}{\GM}{\Uni}$ then we are able to simply write $\Modify[\TM]{A}$ rather than
$\Modify[\TM]{A^{\TM\Whisker\epsilon}}$. By convention, we also avoid writing $x^{\ArrId{}}$.

\begin{nota}
  If $\IsTm[\LockCxV{\Gamma}]{f}{A \to B}$, we write $f^\dagger$ for the function
  $\MkMod[\mu]{f} \mathbin{\circledast} -$.
\end{nota}

\begin{rem}\label{rem:prelims:coe-natural}
  Since we shall capitalize on the fact repeatedly, we note that $\Coe^\alpha$ is always suitably
  natural. For instance, fix $\DeclVar{f}{\GM}{A \to B}$. Then we construct a path
  $\alpha : \Coe^{\pi_1^\TM} \circ f^\dagger = f \circ \Coe^{\pi_1^\TM}$ as follows:
  \[
    \alpha = \Con{funext}\prn{\lambda x.\,\LetMod[\TM]{x}[x_0]{\Refl}}
  \]
  Since this path is essentially a commuting conversion (it is given by induction to allow
  $\Coe^{\pi_1^\TM}$ and $f^\dagger$ to reduce) it is fully coherent, with higher paths being
  likewise constructed by induction and then reflexivity.
\end{rem}

In general, $\Modify[\mu]{-}$ need not commute with propositional equality. However, this is true in
our intended models and so we impose it as an axiom:

\begin{restatable}{axiom}{crispidaxiom}
  The map $\MkMod[\mu]{a} = \MkMod[\mu]{b} \to \Modify{a = b}$ sending $\Refl$ to $\MkMod{\Refl}$ is
  an equivalence for all $\DeclVar{a,b}{\mu}{A}$.%
\end{restatable}

To be very precise, this map is defined by path induction in the family of types
$\lambda\prn{x,y:\Modify{A}}.\,\LetMod{x}[a]{\LetMod{y}[b]{\Modify{a = b}}}$. By
\textcite{gratzer:normalization:2022}, there is a computational account of this principle.

\begin{cor}
  Each $\Modify[\mu]{-}$ commutes with pullbacks
  $A \times_C B = \Sum{a : A} \Sum{b : B} f\prn{a} =_C g\prn{b}$.
\end{cor}

\begin{rem}
  For readers familiar with \emph{spatial type theory}~\parencite{shulman:2018}, this modal type
  theory is an extension of spatial type theory to include two additional modalities ($\OM$,
  $\TM$). In particular, the results of \textcite{shulman:2018} that deal with $\GM$ and $\SM$ can
  be reproduced in this setting.
\end{rem}

\subsection{Modalities and simplicial type theory}
\label{sec:prelims:modal-stt}

To connect the modal and simplicial structures, we impose the following axioms motivated by the
intended model, as described in Theorem~\ref{thm:prelims:stt-models} (and more generally
$\mathcal E^{\Op{\SIMP}}$ for an $\infty$-topos $\mathcal E$); see also the work of
\textcite{myers:2023}. First, the opposite map should be an anti-equivalence of $\Int$:

\begin{restatable}{axiom}{intopaxiom}
  \label{ax:op-of-int-is-int}
  There is an equivalence $\neg : \Modify[\OM]{\Int} \to \Int$ which swaps $0$ for $1$ and $\lor$ for $\land$.
\end{restatable}
\begin{cor}
  We can extend $\neg$ to an equivalence $\neg : \Modify[\OM]{\Delta^n} \Equiv \Delta^n$.
\end{cor}

Next, we require the two possible notions of discreteness (being $\Int$-null or $\GM$-modal) to
coincide:

\begin{restatable}{axiom}{intdiscaxiom}
  \label{ax:discrete-iff-crisp}
  If $\DeclVar{A}{\GM}{\Uni}$, then $\Modify[\GM]{A} \to A$ is an equivalence \textup($A$ is \emph{discrete}\textup) if
  and only if $A \to A^\Int$ is an equivalence \textup($A$ is \emph{$\Int$-null}\textup).
\end{restatable}
\begin{restatable}{axiom}{intglobalaxiom}
  \label{ax:global-points}
  The canonical map $\Bool \to \Int$ is injective and induces an equivalence $\Bool \Equiv \Modify[\GM]{\Int}$.
\end{restatable}
Motivated by our intended class of models, we insist that equivalences are jointly detected by
$\Delta^n$:

\begin{restatable}{axiom}{cubessepaxiom}
  \label{ax:cubes-separate}
  $\DeclVar{f}{\GM}{A \to B}$ is an equivalence if and only if
    the following holds:
    \[
      \Prod{\DeclVar{n}{\GM}{\Nat}}
      \IsEquiv\prn{
        \prn{f_*}^\dagger : \Modify[\GM]{\Delta^n \to A}
        \to
        \Modify[\GM]{\Delta^n \to B}
      }
    \]
\end{restatable}
\noindent
Note that since there is a section-retraction pair $\Delta^n \to \Int^n \to \Delta^n$, we can
replace $\Delta^n$ with $\Int^n$ in the above principle.

One useful application is the following:
\begin{lem}
  \label{lem:prelims:covariance-simplices}
  A map $\DeclVar{\pi}{\GM}{X \to A}$ is covariant if and only if the map
  $\Modify[\GM]{X^{\Delta^n}} \to \Modify[\GM]X \times_{\Modify[\GM]A} \Modify[\GM]{A^{\Delta^n}}$
  induced by $\prn{0^*}^\dagger$ and $\prn{\pi_*}^\dagger$ is an equivalence
  for all $\DeclVar n\GM\Nat$.
\end{lem}


\subsubsection*{The axiom for $\Modify[\TM]{-}$}
Finally, we add a new axiom to \STT{} that governs $\TM$. For motivation, we recall some facts
about the external definition of the twisted arrow functor $\PSH{\SIMP} \to \PSH{\SIMP}$ which
$\Modify[\TM]{-}$ is intended to internalize. Classically, $\Con{Tw} : \PSH{\SIMP} \to \PSH{\SIMP}$
is defined as follows:
\[
  \Con{Tw}\prn{X}\prn{\brk{n}} = X\prn{\Op{\brk{n}} \ast \brk{n}}
\]
Here we have written $\Op{\brk{n}} \ast \brk{n}$ instead of the equivalent $\brk{2n + 1}$ to clarify
the action of this functor on morphisms: $f \mapsto \Op{f} \ast f$. (I.e., this corresponds to the
join operation on finite linear orders.)

As this functor is defined by precomposition, it is a right adjoint whose left adjoint is defined by
left Kan extension. In particular, it sends $\Delta^n : \PSH{\SIMP}$ to $\Delta^{2n + 1}$ (again,
with the functorial action given by twisting). As such, there is a universal map
$\Mor[\eta_n]{\Delta^n}{\Con{Tw}\prn{\Delta^{2n+1}}}$ which, when unfolded, is given by the identity
$\Mor{\brk{2n + 1}}{\brk{2n + 1}}$. Universality of $\eta_n$ amounts to the requirement that each
morphism $\Mor[f]{\Delta^n}{\Con{Tw}\prn{C}}$ factors as $\Con{Tw}\prn{\hat{f}} \circ \eta_n$ for
some unique $\Mor[\hat{f}]{\Delta^{2n + 1}}{C}$. Our axiom governing $\Modify[\TM]{-}$ axiomatizes
this $\eta_n$ along with the property that
$\Con{Tw}\prn{-} \circ \eta_n : \Hom{\Delta^{2n+1}}{C} \to \Hom{\Delta^n}{\Con{Tw}\prn{C}}$ is an
equivalence. With this external motivation to hand, we proceed to fix some notation and state the
axiom which governs $\Modify[\TM]{-}$.

\begin{nota}
  If $\DeclVar{n}{\GM}{\Nat}$, we have canonical maps $i_l : \Delta^n \to \Delta^{2n + 1}$,
  $i_r : \Delta^n \to \Delta^{2n + 1}$, and $i_m : \Delta^1 \to \Delta^{2n + 1}$ which picks out
  $\brc{0, \dots, n}$, $\brc{n + 1, \dots, 2n + 1}$, and $\brc{n, n + 1}$ respectively.  For
  convenience, we write
  $\bar{i}_l = i_l^\dagger \circ \neg : \Delta^n \to \Modify[\OM]{\Delta^{2n + 1}}$.
\end{nota}

Finally, in the statement of new axiom, we require a procedure which extends a map
$f : {\Delta^n \to \Delta^m}$ to a map ${\Delta^{2n + 1} \to \Delta^{2m + 1}}$ which acts
appropriately on the images of two inclusions $i_l,i_r : \Delta^n \to \Delta^{2n + 1}$. To justify
this formally, we introduce the \emph{blunt join} $X \BluntJoin Y$:
\[
  X \BluntJoin Y = X \Pushout{X \times \brc{0} \times Y} \prn{X \times \Int \times Y} \Pushout{X \times \brc{1} \times Y} Y
\]
This is the directed version of the join $X \star Y$~\parencite[Ch 6]{hottbook} such that
$X \BluntJoin Y$ is roughly $X \Coprod{} Y$ with morphisms adjoined to connect each $x:X$ to each $y:Y$.
\begin{lem}
  \label{lem:prelims:blunt-join-of-simplices}
  If $C$ is a category, then
  $C^{\Delta^{m+1+n}} \Equiv C^{\Delta^{m} \BluntJoin \Delta^n}$.
\end{lem}
\begin{defi}
  If $\DeclVar{f}{\GM}{\Delta^n \to \Delta^m}$ and we take
  $\Con{twist}\prn{f} : {\Delta^{2n + 1} \to \Delta^{2m + 1}}$ to be the map given given by uniquely
  extending the map
  $f^\dagger \BluntJoin f : \Modify[\OM]{\Delta^n} \BluntJoin {\Delta^n} \to \Modify[\OM]{\Delta^m}
  \BluntJoin {\Delta^m}$ along the categorical equivalences
  $\Modify[\OM]{\Delta^i} \BluntJoin {\Delta^i} \to \Delta^{2i + 1}$.
\end{defi}

\begin{restatable}{axiom}{twarraxiom}
  \label{ax:twisted-arrow-equiv}
  For each $\DeclVar{n}{\GM}{\Nat}$, there is a \textup(necessarily unique\textup) function
  $\DeclVar{\eta_n}{\GM}{\Delta^n \to \Modify[\TM]{\Delta^{2n+1}}}$ such that the following map is an
  equivalence, for each category $\DeclVar{C}{\GM}{\Uni}$:
  \[
    \iota \defeq \lambda \MkMod[\GM]{f}.\,\MkMod[\GM]{f^\dagger \circ \eta_n}
    :
    \Modify[\GM]{\Delta^{2n + 1} \to C}
    \to
    \Modify[\GM]{\Delta^{n} \to \Modify[\TM]{C}}
  \]
  Additionally, we require that
  $\tau = \prn{\Coe^{\neg}}^\dagger : \Modify[\TM]{\Delta^{n}} \to \Modify[\TM]{\Modify[\OM]{\Delta^{n}}}$ and
  that the diagrams in Figure~\ref{fig:prelims:tm-diagrams} commute \textup(these are mere properties---all
  objects are sets since $\Modify[\mu]{-}$ preserves h-level\textup).
\end{restatable}
\begin{figure}
  \centering
  \[
    \begin{tikzpicture}[diagram]
      \node (A) {$\Delta^n$};
      \node [right = 3cm of A] (T) {$\Modify[\TM]{\Delta^{2n + 1}}$};
      \node [above = of T] (B0) {$\Modify[\OM]{\Delta^{2n + 1}}$};
      \node [below = of T] (B1) {$\Delta^{2n + 1}$};
      \path[->] (A) edge node[above] {$\eta_n$} (T);
      \path[->] (T) edge node[right] {$\Coe^{\pi^\TM_0}$} (B0);
      \path[->] (T) edge node[right] {$\Coe^{\pi^\TM_1}$} (B1);
      \path[->] (A) edge node[below] {$i_r$} (B1);
      \path[->] (A) edge node[above] {$\bar{i}_l$} (B0);
    \end{tikzpicture}
    \qquad
    \begin{tikzpicture}[diagram]
      \SpliceDiagramSquare{
        width = 3cm,
        nw = \Delta^n,
        sw = \Delta^m,
        ne = \Modify[\TM]{\Delta^{2n + 1}},
        se = \Modify[\TM]{\Delta^{2m + 1}},
        north = \eta_n,
        south = \eta_m,
        west = f,
        east = \Con{twist}\prn{f}^\dagger
      }
    \end{tikzpicture}
  \]

  \caption{Laws for Axiom~\ref{ax:twisted-arrow-equiv}.}
  \label{fig:prelims:tm-diagrams}
\end{figure}

One may visualize $\iota$ as ensuring that $\Modify[\GM]{\Delta^n \to \Modify[\TM]{C}}$ is
isomorphic to a $2n + 1$ simplex in $C$:
\begin{equation*}
  \begin{tikzpicture}[diagram]
    \node (A0) {$c_n$};
    \node [right = of A0] (A1) {$c_{n-1}$};
    \node [right = of A1] (A2) {$\cdots$};
    \node [right = of A2] (A3) {$c_0$};
    \node [below = 1cm of A0] (B0) {$c_{n+1}$};
    \node [right = of B0] (B1) {$c_{n+2}$};
    \node [right = of B1] (B2) {$\cdots$};
    \node [right = of B2] (B3) {$c_{2n}$};
    \path[<-] (A0) edge (A1);
    \path[<-] (A1) edge (A2);
    \path[<-] (A2) edge (A3);
    \path[->] (B0) edge (B1);
    \path[->] (B1) edge (B2);
    \path[->] (B2) edge (B3);
    \path[->] (A0) edge (B0);
  \end{tikzpicture}
\end{equation*}

Under this correspondence, $\eta$ is the unique map $\Delta^n \to \Modify[\TM]{\Delta^{2n + 1}}$
given by the identity $\ArrId{} : \Delta^{2n + 1} \to \Delta^{2n + 1}$ and is thus the universal
$n$-simplex. The map $\pi^\TM_1$ picks out the bottom row and $\pi^\TM_0$ selects the top but
\emph{twisted} so that it lands in $\Modify[\OM]{C}$ rather than $C$. This axiom will only be used
in the proof of Theorem~\ref{thm:tw:tw}, where we use $\Modify[\TM]{-}$ to construct a bifunctorial version
of $\hom$.

\begin{prop}[\textcite{gratzer:2024}]
  The model constructed in Theorem~\ref{thm:prelims:stt-models} extends to a model of modal \HoTT{}
  validating our axioms.
\end{prop}
\begin{rem}
  While our previous work~\parencite{gratzer:2024} did not handle $\Modify[\TM]{-}$, the methods
  employed there scale directly to this situation. In particular, \textcite{mukherjee:2023} give an
  explicit description of the necessary twisted arrow operation and show it is a Quillen right
  adjoint as required to extend the model.
\end{rem}

With modalities to hand, a number of results from classical category theory can be proven
directly. For instance, the so-called fundamental theorem of $\infty$-category theory:

\begin{restatable}{thm}{ffesssurj}
  \label{thm:prelims:ff-ess-surj-to-equiv}
  If $\DeclVar{C,D}{\GM}{\Uni}$ are categories, then $\DeclVar{F}{\GM}{C \to D}$ is an equivalence if
  (1) the induced map $\Modify[\GM]{C} \to \Modify[\GM]{D}$ is surjective, and (2) for any
  $\DeclVar{c,c'}{\GM}{C}$ the map $\Hom{c}{c'} \to \Hom{F\prn{c}}{F\prn{c'}}$ is an equivalence.
\end{restatable}
\begin{proof}
  Suppose (1) and (2) holds. We prove that $F$ is an equivalence using Axiom~\ref{ax:cubes-separate} and
  fix $\DeclVar{n}{\GM}{\Nat}$ such that it suffices to show
  $\IsEquiv\prn{F_*^\dagger : \Modify[\GM]{\Delta^n \to C} \to \Modify[\GM]{\Delta^n \to D}}$.

  If $n = 0$, then by (1) $F_*^\dagger$ is surjective and by (2) combined with the Rezk condition,
  it is an embedding. Accordingly, $F_*^\dagger$ is an equivalence in this case. The case for
  $n = 1$ is an immediate consequence of the cases for $n = 0$ along with (2). In general,
  since $C^{\Delta^n} \Equiv C^{\Delta^1} \times_{C} \dots \times_C C^{\Delta^1}$ by the Segal
  condition and likewise for $D$, and $\Modify[\GM]{-}$ commutes with pullbacks, the case for
  $n \ge 2$ follows from $n=0,1$.
\end{proof}

\subsection{Basic building blocks for categories}
Finally, we recall two results from our earlier work~\parencite{gratzer:2024} that will be used repeatedly within this work
to construct new categories. The first is a construction of \emph{full subcategories} using $\SM$:

\begin{prop}\label{thm:prelims:full-subcat}
  %
  If $\DeclVar{C}{\GM}{\Uni}$ is a category and $\DeclVar{\phi}{\GM}{\Modify[\GM]{C} \to \Prop}$ is
  a predicate, then 
  \begin{enumerate}
  \item $C_\phi = \Sum{c : C} \Modify[\SM]{\phi\prn{\MkMod[\GM]{c}}}$ is a category,

  \item the projection map $C_\phi \to C$ induces an equivalence on hom-types,

  \item $\Modify[\GM]{C_\phi} \Equiv \Sum{c : \Modify[\GM]{C}} \phi\prn{c}$, and 

  \item a map $\DeclVar{F}{\GM}{D \to C}$ factors through $C_\phi$ if and only if
    $\phi\prn{\MkMod[\GM]{F\prn d}}$ holds for all $\DeclVar{d}{\GM}{D}$.
\end{enumerate}
\end{prop}

\begin{cor}\label{cor:prelims:crisp-ff-is-ff}
  %
  If $\DeclVar{C,D}{\GM}{\Uni}$ are categories and $\DeclVar{F}{\GM}{C \to D}$, then the canonical
  map $\Hom{c}{c'} \to \Hom{F\prn{c}}{F\prn{c'}}$ is an equivalence for all $c,c' : C$ (notice the
  lack of $\GM$!) if and only if it is an equivalence when $\DeclVar{c,c'}{\GM}{C}$.
  %
\end{cor}

Next, we recall the construction of the category of groupoids which plays the role of the
category of sets in simplicial type theory, \eg{}, we shall use this category to define presheaves:

\begin{prop}
  \label{thm:prelims:space} There is a category $\DeclVar{\Space_i}{\GM}{\Uni[i + 1]}$ with an
  embedding $\Space_i \to \Uni[i]$ such that:
  \begin{itemize}
    \item If $X : A \to \Space_i$, then the composite $A \to \Uni[i]$ is covariant.
    \item The converse holds for $\DeclVar{A}{\GM}{\Uni[i]}$, $\DeclVar{X}{\GM}{A \to \Uni[i]}$: if
      $X$ is covariant, then $X$ factors through $\Space_i$.
  \end{itemize}
\end{prop}

\begin{cor}[Directed univalence]
  \label{cor:prelims:dua}
  $\Space^\Int \Equiv \Sum{X_0,X_1 : \Space} X_1^{X_0}$ and composition in $\Space$ is the
  composition of functions.
\end{cor}

\begin{cor}
  If $\DeclVar{X}{\GM}{\Uni}$ is a groupoid, then $X : \Space$.
\end{cor}

\begin{rem}
  We proved~\parencite{gratzer:2024} Proposition~\ref{thm:prelims:space} in a richer variation of \STT{}
  (\emph{triangulated type theory}). Since we only require the result here, we take it as an ``axiom''
  of sorts to work in a simpler type theory and note that one could extend \STT{} to triangulated
  type theory to prove this theorem outright.
\end{rem}

\section{The Yoneda embedding}
\label{sec:tw}

Within this section, we fix a category $\DeclVar{C}{\GM}{\Uni}$. Our goal is to study the type
$\PSH{C} = \Space^{\Modify[\OM]{C}}$ of presheaves on $C$. As $\Space$ is a category, so
is $\PSH{C}$, and by directed univalence:
\begin{lem}
  \label{lem:tw:hom-of-psh}
  If $F,G : \PSH{C}$ then
  $\Hom{F}{G} \Equiv \Prod{\DeclVar{c}{}{\Modify[\OM]{C}}}{F\,c \to G\,c}$
\end{lem}

\begin{rem}
  Just as with \eg{}, completeness, $\PSH{C}$ implicitly fixes a universe level such that
  $\PSH{C} = \Modify[\OM]{C} \to \Space_i$. We may regard $i$ as a parameter or simply take $i =
  0$. Occasionally, we shall need to insist that $C \Equiv C'$ where $C' : \Uni[i]$ and in such
  situations we shall say that $C$ is \emph{small}. We assume all categories are locally
  small---that each $\Hom[C]{c}{c'}$ is small.
\end{rem}

One may recast the \emph{fibrational} Yoneda lemma proven by \textcite{riehl:2017} to take advantage
of $\PSH{C}$ rather than quantifying over contravariant families as in \opcit:
\begin{lem}
  \label{lem:tw:fibrational-yoneda}
  If $F : \PSH{C}$ and $c : \Modify[\OM]{C}$ then
  $F\prn{c} \cong \Prod{c' : \Modify[\OM]{C}} \Hom[\Modify[\OM]{C}]{c}{c'} \to F\prn{c'}$
\end{lem}

\subsection{The twisted arrow category and the Yoneda embedding}

In light of this last result, the natural next step is to define a map $C \to \PSH{C}$
which sends $c : C$ to something like $\Hom{-}{c}$.\footnote{Here
  we see why $C$ must be flat: we wish to discuss both $C$ and $\Modify[\OM]{C}$. It is helpful to
  understand $\DeclVar{C}{\GM}{\Uni}$ as a \emph{closed} type which depends on nothing in the
  context and, in particular, need not be treated functorially.}
However, caution is required: $\Hom{-}{c}$ has type $C \to \Uni$ and not the
required $\Modify[\OM]{C} \to \Space$. Upon reflection, the reader should find it surprising that
$\Hom{-}{-} : C \times C \to \Uni$ at all; if all maps are functorial in \STT{} how can $\Hom{-}{-}$
be covariant in both arguments? In fact, this is a consequence of the strange behavior of synthetic
morphisms in $\Uni$. While $\Hom{-}{-}$ is functorial in both arguments, the lack of directed
univalence for $\Uni$ makes this useless. This strangeness ensures that $\Hom{-}{-}$ does not
restrict to a function into $\Space$.

What is required instead is a function $\Phi : \Modify[\OM]{C} \times C \to \Space$ such that
$\Phi\prn{\MkMod[\OM]{c},-} = \Hom{c}{-}$ whenever $\DeclVar{c}{\GM}{C}$,
\ie{}, a function that agrees on
objects with $\Hom{-}{-}$ and has the same functoriality in the second argument, but takes
$\Modify[\OM]{C}$ as its first argument. In fact, it is highly non-obvious where such a function
should come from; \textcite[p.~xii]{riehl:2022} specifically highlight this construction as
remarkably subtle in $\infty$-category theory. It is for this reason that we introduced
$\Modify[\TM]{-}$. Recall the visualization of $\Modify[\GM]{\Delta^n \to \Modify[\TM]{C}}$:
\begin{equation}
  \label{eq:tw:visual}
  \begin{tikzpicture}[diagram,baseline=(link)]
    \node (A0) {$c_n$};
    \node [right = of A0] (A1) {$c_{n-1}$};
    \node [right = of A1] (A2) {$\cdots$};
    \node [right = of A2] (A3) {$c_0$};
    \node [below = 1cm of A0] (B0) {$c_{n+1}$};
    \node [right = of B0] (B1) {$c_{n+2}$};
    \node [right = of B1] (B2) {$\cdots$};
    \node [right = of B2] (B3) {$c_{2n}$};
    \path[<-] (A0) edge (A1);
    \path[<-] (A1) edge (A2);
    \path[<-] (A2) edge (A3);
    \path[->] (B0) edge (B1);
    \path[->] (B1) edge (B2);
    \path[->] (B2) edge (B3);
    \path[->] (A0) edge coordinate[pos=.5] (link) (B0);
  \end{tikzpicture}
\end{equation}
The projection to $\Modify[\OM]{C}$ gives the top row and the map to $C$ yields the bottom. This
visualization for $n$-simplices is very similar to that of
$C^{\Int} = \Sum{c_0,c_1} \Hom{c_0}{c_1}$, but the top row has been twisted to ensure that
one restriction lands in $\Modify[\OM]{C}$ as required for a bifunctorial version of $\Hom{-}{-}$:

\begin{restatable}{thm}{twlfib}
  \label{thm:tw:tw}
  If $\DeclVar{C}{\GM}{\Uni}$ is a category, then the following holds:
  \begin{itemize}
  \item The map
    $\gl{\Coe^{\pi_0^\TM},\Coe^{\pi_1^\TM}} : \Modify[\TM]{C} \to \Modify[\OM]{C} \times C$
    straightens to $\Phi : \Modify[\TM]{C} \times C \to \Space$.
  \item For every $\DeclVar{c}{\GM}{C}$, the map
    $\alpha_c : \Hom{\Hom{c}{-}}{\Phi\prn{\MkMod[\OM]{c},-}}$ induced by the Yoneda lemma
    (Lemma~\ref{lem:tw:fibrational-yoneda})
    applied to $\iota\prn{\MkMod[\GM]{\ArrId{c}}} : \Phi\prn{\MkMod[\OM]{c},c}$ is
    an equivalence.
  \end{itemize}
\end{restatable}

\begin{lem}
  \label{lem:tw:simple-coherence}
  Given $\DeclVar{f}{\GM}{\Delta^{1} \to C}$ and let
  $\bar{f} = \iota\prn{\MkMod[\GM]{f}} : \Modify[\TM]{C}$, then there exist  paths:
  \begin{gather*}
    \theta\prn{f}_0 :
    \prn{\Coe^{\pi_0^{\TM}} \circ \Con{extract}\prn{\bar{f}}}\prn{\ast}
    =
    \MkMod[\OM]{f\prn{0}}
    \quad
    \theta\prn{f}_1 :
    \prn{\Coe^{\pi_1^{\TM}} \circ \Con{extract}\prn{\bar{f}}}\prn{\ast}
    =
    f\prn{1}
  \end{gather*}
  These paths are natural in $C$ so that \eg{}, the two paths of the following shape induced by
  $\theta\prn{g \circ f}_1$ and $\theta\prn{f}_1$ with the naturality of $\Coe^{\pi_1^\TM}$ at $g$ agree:
  \[
    \prn{\Coe^{\pi_1^{\TM}} \circ \Con{extract}\prn{\iota\prn{\MkMod[\GM]{g \circ f}}}}\prn{\ast}
    =
    g\prn{f\prn{1}}
  \]
  Here $\ast : \Delta^0$ is the unique element of the unit type.
\end{lem}
\begin{proof}
  We show the second, as they are symmetric. We define $\theta_1$ using the naturality of
  $\Coe^{\pi_1^{\TM}}$ and the behavior of $\Coe^{\pi_1^{\TM}}$ on $\eta$
  from Figure~\ref{fig:prelims:tm-diagrams}:
  \begin{align*}
    &\prn{\Coe^{\pi_1^{\TM}} \circ \Con{extract}\prn{\iota\prn{\MkMod[\GM]{f}}}}\prn{\ast}
    \\
    &= \prn{\Coe^{\pi_1^{\TM}} \circ f^\dagger \circ \eta_0}\prn{\ast}
    \\
    &= \prn{f \circ \Coe^{\pi_1^{\TM}} \circ \eta_0}\prn{\ast} && \text{By naturality, Remark \ref{rem:prelims:coe-natural}}
    \\
    &= f\prn{1} && \text{By the first diagram in Figure~\ref{fig:prelims:tm-diagrams}}
  \end{align*}

  To prove that $\theta_1$ is natural in $C$, we observe that the terms agree up to a commuting
  conversion of elimination rules for modal types. Accordingly, we may prove that these two paths
  agree by induction on $\eta_0\prn{\ast}$ and then reflexivity.
\end{proof}

\begin{proof}[Proof of the Theorem~\ref{thm:tw:tw}]
  We begin by showing that the map
  $\gl{\Coe^{\pi^\TM_0},\Coe^{\pi^\TM_1}} : \Modify[\TM]{C} \to \Modify[\OM]{C} \times C$ is a
  covariant family. By Lemma~\ref{lem:prelims:covariance-simplices} the following map induced by
  $\brc{0} : \Delta^0 \to \Delta^n$ is an equivalence:
  \[
    \epsilon : \Modify[\GM]{\Modify[\TM]{C}^{\Delta^n}} \to
    \prn{
      \DelimMin{1}
      \Modify[\GM]{\Modify[\TM]{C}}
      \times_{\Modify[\GM]{\Modify[\OM]{C}} \times \Modify[\GM]{C}}
      \prn{\Modify[\GM]{\Modify[\OM]{C}^{\Delta^n}} \times \Modify[\GM]{C^{\Delta^n}}}
    }
  \]
  For convenience, we begin by applying a few modal transformations (in particular, using
  $\Con{transp} : \Modify[\GM]{A \to \Modify[\OM]{B}} \Equiv \Modify[\GM]{\Modify[\OM]{A} \to B}$)
  such that it suffices to show that the following map is an equivalence:
  \[
    \epsilon' :
    \Modify[\GM]{\Modify[\TM]{C}^{\Delta^n}} \to
    \prn{
      \DelimMin{1}
      \Modify[\GM]{\Modify[\TM]{C}}
      \times_{\Modify[\GM]{C} \times \Modify[\GM]{C}}
      \prn{\Modify[\GM]{C^{\Modify[\OM]{\Delta^n}}} \times \Modify[\GM]{C^{\Delta^n}}}
    }
  \]

  To prove this, we shall construct a commutative diagram:
  \begin{equation}
    \DiagramSquare{
      width = 7cm,
      sw = \Modify[\GM]{\Modify[\TM]{C}^{\Delta^n}},
      nw = \Modify[\GM]{C^{\Delta^{2n+1}}},
      se =
      \Modify[\GM]{\Modify[\TM]{C}}
      \times_{\Modify[\GM]{C} \times \Modify[\GM]{C}}
      \prn{\Modify[\GM]{C^{\Modify[\OM]{\Delta^n}}} \times \Modify[\GM]{C^{\Delta^n}}},
      ne =
      \Modify[\GM]{\Delta^1 \to C}
      \times_{\Modify[\GM]{C} \times \Modify[\GM]{C}}
      \prn{\Modify[\GM]{C^{\Modify[\OM]{\Delta^n}}} \times \Modify[\GM]{C^{\Delta^n}}},
      north = \phi,
      south = \epsilon',
      west = \iota,
      east = \prn{\iota,\ArrId{}}
    }
    \label{diag:key-diagram}
  \end{equation}
  We define $\phi$ momentarily, but we first remark that $\prn{\iota,\ArrId{}}$ is
  well-formed because of Lemma~\ref{lem:tw:simple-coherence}, which ensures that applying $\iota$ and then
  evaluating commutes appropriately with projection.

  Note also that the two vertical maps are equivalences because $\iota$ is an
  equivalence. Accordingly, by 3-for-2, to show $\epsilon'$ is an equivalence,
  it suffices to ensure that $\phi$ is an equivalence making the diagram commute.\footnote{%
    We emphasize that the filler for this square is irrelevant. Any filler
    suffices to show that $\epsilon'$ is an equivalence, which in turn implies that $\epsilon$ is an
    equivalence as required.}
  We now define $\phi$ as follows:
  \[
    \phi\prn{\MkMod[\GM]{f}} \defeq
    \prn{
      \MkMod[\GM]{\Restrict{f}{n \le n + 1}},
      \prn{
        \MkMod[\GM]{\Restrict{f}{0 \le \dots \le n} \circ \neg},
        \MkMod[\GM]{\Restrict{f}{n + 1 \le \dots \le 2n + 1}}
      },
      \Refl
    }
  \]
  $\phi$ is given by restricting along a categorical equivalence
  ($\Modify[\OM]{\Delta^{n}} \BluntJoin \Delta^n \to \Delta^{2n + 1}$),
  so it is an equivalence.

  Next, we note that all four of the maps in this diagram are weakly natural in $C$. For the bottom
  and top maps, this is an easy observation---the top is given by restriction and the bottom uses
  restriction along with $\Coe^{\pi_1^{\TM}}$, which is also natural
  by Remark~\ref{rem:prelims:coe-natural}. For the left-hand map, this is a consequence of the naturality
  of $\iota$. For the right-hand map, the only wrinkle is the paths used to witness that $\iota$
  commutes with evaluating on projections. This requires a filler for a certain path, but this is
  precisely the naturality coherence supplied by the latter part of Lemma~\ref{lem:tw:simple-coherence}.

  Finally, we argue that the diagram commutes.  Fix
  $\MkMod[\GM]{f} : \Modify[\GM]{\Delta^{2n + 1} \to C}$. We wish to show that
  $\epsilon'\prn{\iota\prn{\MkMod[\GM]{f}}}
  = \prn{\iota,\ArrId{}}\prn{\phi\prn{\MkMod[\GM]{f}}}$. By naturality, however, we
  may use $f$ to reduce to the case where $C = \Delta^{2n + 1}$ and $f = \ArrId{}$. In this case,
  everything involved is a set and so it suffices to argue the diagram commutes when we replace each
  pullback with a simple product. With this in place, we now calculate:
  \begin{align*}
    \epsilon'\prn{\iota\prn{\MkMod[\GM]{\ArrId{}}}}
    &= \epsilon'\prn{\MkMod[\GM]{\eta_n}}
    \\
    &=
      \prn{
        \MkMod[\GM]{\eta_n\,0},
        \prn{
          \Con{transp}\prn{\MkMod[\GM]{\Coe^{\pi_0^{\TM}}\circ \eta_n}},
          \MkMod[\GM]{\Coe^{\pi_1^{\TM}} \circ \eta_n}
        }
      }
    \\
    &=
      \prn{
        \MkMod[\GM]{\eta_n\,0},
        \prn{
          \Con{transp}\prn{\MkMod[\GM]{\bar{i}_l}},
          \MkMod[\GM]{i_r}
        }
      }
    \\
    &=
      \prn{
        \MkMod[\GM]{i_m\prn{\eta_0\,\ast}},
        \prn{
          \MkMod[\GM]{i_l \circ \neg},
          \MkMod[\GM]{i_r}
        }
      }
    \\
    &=
      \prn{\iota,\ArrId{}}\prn{
        \MkMod[\GM]{i_m},
        \prn{
          \MkMod[\GM]{i_l \circ \neg},
          \MkMod[\GM]{i_r}
        }
      }
    \\
    &=\prn{\iota,\ArrId{}}\prn{\phi\prn{\MkMod[\GM]{\ArrId{}}}}
  \end{align*}

  This completes the first step of the argument. The second is to show that for each
  $\DeclVar{c}{\GM}{C}$ the map
  $\alpha_c : \Hom[C \to \Space]{\Hom{c}{-}}{\Phi\prn{\MkMod[\OM]{c},-}}$ is an isomorphism. Passing
  to total spaces, it suffices to show the following map is an equivalence:
  \[
    \tilde{\alpha}_c = \lambda \prn{d,f}.\,\prn{d,f_*\prn{\iota\prn{\MkMod[\GM]{\ArrId{c}}}}}
    : \Sum{d : C} \Hom{c}{d} \to \Sum{d : C} \Phi\prn{\MkMod[\OM]{c},d}
  \]
  In the above, $f_*$ is the covariant transport operation on $\Phi\prn{\MkMod[\OM]{c},-}$.
  Since both sides of this map are categories, it suffices to show that this map is
  fully faithful and essentially surjective.

  In fact, $\tilde\alpha_c$ is an equivalence on objects. To this end, we observe that if
  $\DeclVar{f}{\GM}{\Hom{c}{d}}$ for some
  $\DeclVar{d}{\GM}{C}$, then we can construct the transport $f_*$ alternatively as follows.
  Define a path $h : \Delta^1 \to \Modify[\TM]{C}$ by $h \defeq\iota{\prn{\MkMod[\GM]{\lambda \_, \_, k.\, f\prn{k}}}}$,
  i.e.,
  $h$ corresponds to the following doubly degenerate 3-simplex in $C$:
  \[
    \begin{tikzcd}
      c\ar[d,"\ArrId{c}"'] & c\ar[l,"\ArrId{c}"'] \\
      c\ar[r,"f"'] & d
    \end{tikzcd}
  \]
  We then consider the morphism $\lambda i.\,\prn{\pi^\TM_1\prn{h\,i}, h\,i}$ in
  $\Sum{d : C} \Phi\prn{\MkMod[\OM]{c},d}$. Using the definition of $\iota$ and the naturality of
  $\eta$, this is a morphism from $\prn{c,\iota\prn{\MkMod[\GM]{\ArrId{c}}}}$ to
  $\prn{d, \iota\prn{\MkMod[\GM]{f}}}$. Moreover, the naturality of $\Coe^{\pi_1^{\TM}}$ ensures
  that it lies over $f$ in $C$. Consequently,
  $\tilde{\alpha}_c\prn{d,f} = \prn{d,\iota\prn{\MkMod[\GM]{f}}}$ when restricted to
  $\DeclVar{d}{\GM}{C}$ and $\DeclVar{f}{\GM}{\Hom{c}{d}}$, which is an equivalence because $\iota$
  is invertible.

  For fully-faithfulness, it suffices to show that the following map is an equivalence:
  \[
    \tilde{\alpha}_c :
    \Modify[\GM]{\Int \to \Sum{d : C} \Hom{c}{d}}
    \to
    \Modify[\GM]{\Int \to \Sum{d : C} \Phi\prn{\MkMod[\OM]{c},d}}
  \]
  However, as both sides are the total spaces of covariant families, it suffices to show that the
  following map is an equivalence:
  \[
    \tilde{\alpha}_c :
    \Modify[\GM]{\Sum{d : \Int \to C} \Hom{c}{d\,0}}
    \to
    \Modify[\GM]{\Sum{d : \Int \to C} \Phi\prn{\MkMod[\OM]{c},d\,0}}
  \]
  The conclusion now follows from the previous case.
\end{proof}

\begin{nota}
  We write $\Phi_D : \Modify[\OM]{D} \times D \to \Space$ for the same construction applied to some
  category $D$. Within this section, we continue to write $\Phi$ as shorthand for $\Phi_C$.
\end{nota}

\begin{cor}
  \label{cor:tw:tw-of-op}
  If $\DeclVar{c_0}{}{\Modify[\OM]{C}}$ and $c_1 : C$, then
  $\Phi\prn{c_0,c_1} = \Phi_{\Modify[\OM]{C}}\prn{\MkMod[\OM]{\MkMod[\OM]{c_1}},c_0}$.
\end{cor}
\begin{proof}
  Passing to total spaces, it suffices to find an equivalence $\Modify[\TM]{C} \to
  \Modify[\TM]{\Modify[\OM]{C}}$ fitting into the following diagram:
  \[
    \begin{tikzpicture}[diagram]
      \node (T) {$\Modify[\TM]{C}$};
      \node[right = 5cm of T] (OT) {$\Modify[\TM]{\Modify[\OM]{C}}$};
      \node[below right = 1.75cm and 2.5cm of T] (C) {$\Modify[\OM]{C} \times C$};
      \path[->] (T) edge (C);
      \path[->] (OT) edge (C);
      \path[->] (T) edge (OT);
    \end{tikzpicture}
  \]

  The map $\Coe^\tau$ precisely satisfies this role: it is invertible because the 2-cell $\tau$
  is an isomorphism and it fits into the commuting diagram because of the corresponding diagram
  in the mode theory.
\end{proof}

\subsection{The Yoneda lemma}

With a bi-functorial version of $\Hom{-}{-}$ to hand, we can now straightforwardly define the Yoneda
embedding $\Yo$ and leverage Lemma~\ref{lem:tw:fibrational-yoneda} into a result about $\Yo$:
\begin{defi}[Yoneda]
   $\Yo = \lambda c.\,\Phi\prn{-,c} : C \to \PSH{C}$.
\end{defi}
\begin{lem}
  \label{lem:tw:weak-yoneda}
  $\Hom{\Yo{c}}{X} \cong X\prn{\MkMod[\OM]{c}}$ for all $X : \PSH{C}$ and $\DeclVar{c}{\GM}{C}$.
\end{lem}
\begin{proof}
    Since $c$ is $\GM$-annotated, using Theorem~\ref{thm:tw:tw} and Corollary~\ref{cor:tw:tw-of-op} we have the following
    identification
    $\Hom[\Modify[\OM]{C}]{\MkMod[\OM]{c}}{-} = \Phi\prn{-,c}$. Moreover, by Lemma~\ref{lem:tw:hom-of-psh} we
    additionally have the following:
    \[
      \Prod{c' : \Modify[\OM]{C}} \Phi\prn{c',c} \to X\prn{c'} \cong \Hom{\Yo{c}}{X}
    \]
    The conclusion now follows by Lemma~\ref{lem:tw:fibrational-yoneda}.
\end{proof}

A great deal of category theory is contained within Lemma~\ref{lem:tw:weak-yoneda}.
It shows that $\Yo$ is
fully-faithful on $\GM$-annotated elements of $C$ and that $C$ is a full subcategory of $\PSH{C}$:
\begin{lem}
  \label{lem:tw:repr}
  $\Yo : C \to \PSH{C}$ induces an equivalence $C \Equiv \PSH{C}_{\Con{isRepr}}$ where
  $\Con{isRepr} = \lambda X.\,\Sum{c : C} X = \Yo{c}$.\footnote{Note that $\Con{isRepr}\prn{X}$ is a
  proposition due to Lemma~\ref{lem:tw:weak-yoneda} and Corollary \ref{cor:prelims:crisp-ff-is-ff}.}
\end{lem}
While Lemma~\ref{lem:tw:weak-yoneda} follows directly from
Lemma~\ref{lem:tw:fibrational-yoneda}, the above consequence can only be expressed once there exists a
\emph{category} of presheaves---something missing from \textcite{riehl:2017}. This opens up a new
proof strategy: to prove a result of $C$, we first prove that it holds for $\Space$, then
$\PSH{C}$, then that it restricts to the full subcategory. For instance, we may prove the
aforementioned characterization of natural isomorphisms:
\begin{thm}
  \label{thm:tw:nat-iso}
  If $\DeclVar{C,D}{\GM}{\Uni}$ are categories, $\DeclVar{F,G}{\GM}{C \to D}$, and
  $\DeclVar{\alpha}{\GM}{\Hom{F}{G}}$, then $\Prod{c : C} \IsIso\prn{\alpha\,c}$ if
  $\Prod{\DeclVar{c}{\GM}{C}} \IsIso\prn{\alpha\,c}$.
\end{thm}
\begin{proof}
  Note that this theorem is trivial for $C = \Delta^0$ and for $C = \Delta^1, D = \Space$ it is
  a consequence of Corollary~\ref{cor:prelims:dua}. The Segal condition for $\Space$ then implies the theorem
  for $C = \Delta^n,D = \Space$.

  By Lemma~\ref{lem:tw:repr}, it suffices to assume $D = \PSH{D_0}$.  By
  Axiom~\ref{ax:discrete-iff-crisp} and Theorem~\ref{thm:prelims:mtt-facts} it suffices to show that
  $\prn{\Sum{c : C} \IsIso\prn{\alpha\,c}} \to \prn{\Sum{c : C}
    \Modify[\SM]{\IsIso\prn{\alpha\,c}}}$ is an equivalence. By Axiom~\ref{ax:cubes-separate}, it
  suffices to prove for all $n$:
  \[
    \IsEquiv\Bigl(\Modify[\GM]{\prn{\Sum{c : C} \IsIso\prn{\alpha\,c}}^{\Delta^n}}
     \to \Modify[\GM]{\prn{\Sum{c : C} \Modify[\SM]{\IsIso\prn{\alpha\,c}}}^{\Delta^n}}\Bigr)
  \]
  Unfolding and commuting $\GM$ with $\Sum{}$, it suffices to show that
    for every $\DeclVar{c}{\GM}{\Delta^n \to C}$ the following holds:
    \[
      \Sum{\DeclVar{\sigma}{}{\Delta^n}} \IsIso\prn{\alpha\prn{c\,\sigma}}
      \Equiv
      \Sum{\DeclVar{\sigma}{\GM}{\Delta^n}} \IsIso\prn{\alpha\prn{c\,\sigma}}
    \]
  Replacing $\alpha$ with $\alpha \circ c$, however, reduces us
  to the already proven case of $C = \Delta^n, D = \Space$.
\end{proof}

Lemma~\ref{lem:tw:weak-yoneda} is already powerful. However, it does not capture that this
equivalence is \emph{natural} in both $c$ and $X$---or, more precisely, since $c$ is $\GM$-annotated
and the equivalence is in $\Uni$, the naturality it yields is trivial. We are able to prove a far
stronger version of the Yoneda lemma that (1) does not need to assume that $\DeclVar{c}{\GM}{C}$,
and (2) yields the desired functoriality in both $c$ and $X$. To do so, we replace $\Hom{-}{-}$
with $\Phi$:

\begin{restatable}[Functorial Yoneda lemma]{thm}{functorialyoneda}
  \label{thm:tw:strong-yoneda}
  There is a natural isomorphism
  $\Phi_{\PSH{C}}\prn{\Yo^\dagger\prn{-}, -} \cong \Con{eval} : \Modify[\OM]{C} \times \PSH{C} \to \Space$.
\end{restatable}
\begin{rem}
  This result uses a handful of results from Section~\ref{sec:adj}. These forward references are justified:
  we do not use Theorem~\ref{thm:tw:strong-yoneda} till Section~\ref{sec:adj:psh}. We present the proof here for
  conceptual coherence.
\end{rem}
\begin{proof}
  The central difficulty in this proof is to find a map
  $\Phi_{\PSH{C}}\prn{\Yo^\dagger\prn{-}, -} \to \Con{eval}$ which can then be checked to be an
  equivalence. To construct this map, we use the presentation of
  $\Modify[\OM]{C} \times \PSH{C} \to \Space$ as covariant families over
  $\Modify[\OM]{C} \times \PSH{C}$. In particular, we consider the following pullback diagrams:
  \[
    \begin{tikzpicture}[diagram]
      \node (PhiPsh) {$\Modify[\TM]{\PSH{C}}$};
      \node [below = 2cm of PhiPsh] (PshPsh) {$\Modify[\OM]{\PSH{C}} \times \PSH{C}$};
      \node [pullback,left = 3.6cm of PhiPsh] (V) {$V$};
      \node [left = 3.6cm of PshPsh] (CPsh) {$\Modify[\OM]{C} \times \PSH{C}$};
      \node [pullback,left = 3.6cm of V] (PhiC) {$\tilde{\Phi}_C$};
      \node [left = 3.6cm of CPsh] (CC) {$\Modify[\OM]{C} \times C$};
      \path[->] (PhiPsh) edge (PshPsh);
      \path[->] (V) edge (PhiPsh);
      \path[->] (V) edge (CPsh);
      \path[->] (PhiC) edge node[above] {$v$} (V);
      \path[->] (PhiC) edge (CC);
      \path[->] (CPsh) edge node[below] {$\Yo^\dagger \times \ArrId{}$} (PshPsh);
      \path[->] (CC) edge node[below] {$\ArrId{} \times \Yo$} (CPsh);
    \end{tikzpicture}
  \]
  \[
    \begin{tikzpicture}[diagram]
      \node (Spt) {$\Sum{A : \Space} A$};
      \node [below = 2cm of Spt] (S) {$\Space$};
      \node [pullback,left = 3.75cm of Spt] (V) {$W$};
      \node [left = 3.75cm of S] (CPsh) {$\Modify[\OM]{C} \times \PSH{C}$};
      \node [pullback,left = 3.75cm of V] (PhiC) {$\Modify[\TM]{C}$};
      \node [left = 3.75cm of CPsh] (CC) {$\Modify[\OM]{C} \times C$};
      \path[->] (Spt) edge (S);
      \path[->] (V) edge (Spt);
      \path[->] (V) edge (CPsh);
      \path[->] (PhiC) edge node[above] {$w$} (V);
      \path[->] (PhiC) edge (CC);
      \path[->] (CPsh) edge node[below] {$\Con{eval}$} (S);
      \path[->] (CC) edge node[below] {$\ArrId{} \times \Yo$} (CPsh);
    \end{tikzpicture}
  \]
  The claim is then that $V \Equiv W$. To show this, we argue that if we replace the composite
  $\Modify[\TM]{C} \to \Modify[\OM]{C} \times C \to \Modify[\OM]{C} \times \PSH{C}$ with the free
  covariant family, then the maps $\bar{v},\bar{w}$ induced by $v$ and $w$ are both equivalences. The
  conclusion then shows $\bar{v} \circ \bar{w}^{-1}$ is the desired equivalence.

  Using \textcite[Theorem 2.41]{rijke:2020} there is a free covariant fibration
  $Z : \Modify[\OM]{C} \times \PSH{C} \to \Uni$ and if $\DeclVar{c}{\GM}{C}$ and
  $\DeclVar{X}{\GM}{\PSH{C}}$ then by Corollary~\ref{cor:adj:push-pull-formula} we have the following:
  \begin{align*}
    &Z\prn{c, X} =
    \\
    &\ \Grpdify\prn{\Sum{c_0,c_1} \Hom{c_0}{c} \times \Hom{\Yo{c_1}}{X} \times \Phi\prn{c_0,c_1}}
  \end{align*}
  To show that \eg{}, $v$ induces an equivalence, we must show that the following map is an
  equivalence:
  \[
    Z\prn{c,X} \to \Hom{\Yo^\dagger\prn{c}}{X}
  \]
  We may use Theorem~\ref{thm:tw:nat-iso} and assume that there exists $\DeclVar{c'}{\GM}{C}$ such that
  $c = \MkMod[\OM]{c'}$ and that $\DeclVar{X}{\GM}{\PSH{C}}$. Moreover, since the right-hand side is
  a groupoid, this map is uniquely induced by extending the canonical map of the following type:
  \begin{align*}
    &\prn{\Sum{c_0,c_1} \Hom{c_0}{c} \times \Hom{\Yo{c_1}}{X} \times \Phi\prn{c_0,c_1}}
    \\
    &\ \to \Phi\prn{\Yo^\dagger\prn{c},X} \Equiv X\prn{c'}
  \end{align*}
  This map sends $\prn{c_0,c_1,f,\alpha,t}$ to $\alpha\,c\,\prn{\Phi\prn{f,\ArrId{}}\,t}$ and one
  may check directly that the assignment $x \mapsto \eta\prn{c,c',\ArrId{},\ArrId{},F_x}$ is a
  quasi-inverse to this map where $F_x : \Hom{\Yo\prn{c'}}{X}$ corresponds to
  $x : X\prn{\MkMod[\OM]{c'}}$ under Lemma~\ref{lem:tw:weak-yoneda}. The case for $w$ is similar.
\end{proof}

\section{Revisiting adjunctions}
\label{sec:adj}

With presheaves and the Yoneda embedding available, we now revisit the theory of adjoint functors
introduced by \textcite{riehl:2017} in \STT{}. They define a pair of functions $f : C \to D$ and
$g : D \to C$ to be adjoint when equipped with
$\iota : \Prod{c,d} \Hom{f\prn{c}}{d} \Equiv \Hom{c}{g\prn{d}}$. While they produce several
equivalent reformulations using a unit and counit natural transformations, no non-trivial examples
of adjunctions are given---unsurprisingly, since concrete examples of categories in \STT{} are relatively
recent. Even with $\Space$ available it is quite difficult to produce examples of such
adjunctions.

It is far more feasible to construct only $f$ and then show that
$\Phi\prn{f^\dagger\prn{-},d} : \PSH{C}$ is representable for every $\DeclVar{d}{\GM}{D}$. This is
comparable to Theorem~\ref{thm:tw:nat-iso}: we wish to give a functorial definition of either $f$ or $g$
and a \emph{non-functorial} definition of the other, and then show that this can be upgraded to a
full adjunction. In this section, we show that this is indeed possible, and we observe that a number
of important adjunctions and results are then immediately within reach. In particular, we shall use
this technique to prove that $\PSH{C}$ is cocomplete and, moreover, is the free cocompletion of $C$.

\subsection{Pointwise adjunctions to adjunctions}
\label{sec:adj:pw}

Let us begin by formalizing the notion of pointwise adjoints:
\begin{defi}
  We say that $\DeclVar{f}{\GM}{C \to D}$ is a pointwise left adjoint if the following type is
  inhabited:
  \[
    \Prod{\DeclVar{d}{\GM}{D}} \Con{isRepr}\prn{\Phi\prn{f^\dagger\prn{-}, d}}
  \]
  Dually, $f$ is a pointwise right adjoint if $f^\dagger : \Modify[\OM]{C} \to \Modify[\OM]{D}$ is a
  pointwise left adjoint.
\end{defi}

Our main theorem relies on two crucial preliminary results. The first shows that any pointwise left
adjoint $f$ gives rise to a function in the other direction picking out the various (necessarily
unique) representing objects for $\Phi\prn{f^\dagger\prn{-}, d}$.
\begin{lem}
  If $\DeclVar{f}{\GM}{C \to D}$ is a pointwise left adjoint, then the type of morphisms
  $\DeclVar{g}{\GM}{D \to C}$ equipped with a natural isomorphism
  $\iota : \Phi\prn{f^\dagger\prn{-}, -} \cong \Yo \circ g$ is contractible.
\end{lem}
\begin{proof}
  Since $\Yo$ is an embedding, this type is a proposition. It therefore suffices to show that it is
  inhabited. By assumption, $\bar{g} = \Phi\prn{f^\dagger\prn{-}, d}$ is representable for all
  $\DeclVar{d}{\GM}{D}$, and thus it factors through $\PSH{C}_{\Con{isRepr}}$. Post-composing
  with the equivalence $\PSH{C}_{\Con{isRepr}} \Equiv C$ yields the desired $g : D \to C$.
\end{proof}

Using this, we prove a universal case of the theorem improving a pointwise adjoint to an adjoint:
every $\DeclVar{g}{\GM}{D \to C}$ that is a cartesian fibration~\parencite{buchholtz:2023}
such that the fiber
over every $\DeclVar{c}{\GM}{C}$ has an initial object~\parencite{bardomiano:2021} admits a left
adjoint.

\begin{lem}
  \label{lem:adj:functoriality-of-init-objects}
  %
  If $\DeclVar{g}{\GM}{D \to C}$ is cartesian and for
  each $\DeclVar{c}{\GM}{C}$ the fiber $D_c$ has an initial object, then there exists $f : C \to D$
  such that $f\prn{c}$ is initial in $D_c$ for all $c : C$.
\end{lem}

\begin{proof}
  Note that $\Con{hasInitialObj}\prn{D_c}$ is a proposition and, therefore, by
  Axiom~\ref{ax:discrete-iff-crisp} we may assume $\Modify[\GM]{\Con{hasInitialObj}\prn{D_c}}$ holds
  for each $\DeclVar{c}{\GM}{C}$. With this observation to hand, we can show that $g$ is a pointwise
  right adjoint: if $\DeclVar{c}{\GM}{C}$, $\DeclVar{d}{}{D}$:
  \begin{align*}
    \Phi_C\prn{\MkMod[\OM]{c},g\prn{d}}
    &\Equiv \Hom{c}{g\prn{d}}
    \\
    &\Equiv \Hom{\ObjInit{D_c}}{d} && \text{$g$ is cartesian}
    \\
    &\Equiv \Phi_D\prn{\MkMod[\OM]{\ObjInit{D_c}},d}
  \end{align*}

  In this last step, we use our observation that $\Modify[\GM]{\Con{hasInitialObj}\prn{D_c}}$ and,
  crucially, that not only $\Con{hasInitialObj}\prn{D_c}$ holds. In particular, we rely on the fact
  that $\DeclVar{\ObjInit{D_c}}{\GM}{D_c}$.

  Accordingly, we obtain a function $\DeclVar{f}{\GM}{C \to D}$ which sends $\DeclVar{c}{\GM}{C}$ to
  $\ObjInit{D_c}$. It remains to show that $f\prn{c}$ is initial in $D_c$ for all $c : C$.
  Since $D = \Sum{c:C}D_c$, this amounts to the following map being an equivalence:
  $\prn{\Sum{d : D} \Hom[D_{g\prn{d}}]{f\prn{g\prn{d}}}{d}} \to D$.

  To prove this, we use Theorem~\ref{thm:prelims:ff-ess-surj-to-equiv} which allows us to reduce to the
  $\GM$-annotated case, where the conclusion follows from the fact that $f\prn{c}$ is then initial
  in $D_c$.
\end{proof}

\begin{thm}
  \label{thm:adj:pw-to-full-adj}
  Pointwise right adjoints are right adjoints.
\end{thm}
\begin{proof}
  Given a map $\DeclVar{g}{\GM}{D \to C}$, consider the cartesian family
    \[
      \pi : \COMMA{C}{g} = \prn{\Sum{c : C}\Sum{d : D} \Hom{c}{g\prn{d}}} \to C
    \]
  Since $g$ is a
  pointwise right adjoint, each fiber of $\pi$ over $\DeclVar{c}{\GM}{C}$ has an initial object. We
  then apply Lemma~\ref{lem:adj:functoriality-of-init-objects} to obtain $\bar{f} : C \to
  \COMMA{C}{g}$. Finally, the composite $\Proj[2] \circ \bar{f}$ is the desired left adjoint to $g$:
  \begin{align*}
    &\Hom[C]{c}{g\prn{d}}
    \\
    &\quad\Equiv \Sum{\alpha : \Hom[C]{c}{g\prn{d}}}
      \Hom[\COMMA{C}{g}_c]{\bar{f}\prn{c}}{\prn{c, d, \alpha}}
    \\
    &\quad\Equiv
      \Sum{\alpha : \Hom[C]{c}{g\prn{d}}} \Sum{\beta : \Hom[D]{f\prn{c}}{d}}
      g\prn{\beta} \circ \Proj[3]\prn{\bar{f}\prn{c}} = \alpha
    \\
    &\quad\Equiv \Hom[D]{f\prn{c}}{d}
  \end{align*}
  The first step uses the initiality of $\bar{f}\prn{c}$ in the fiber over $c$ and the second
  unfolds the definition of a morphism in $\COMMA{C}{g}$.
\end{proof}

\subsection{Examples of adjunctions}

We take advantage of Theorem~\ref{thm:adj:pw-to-full-adj} to produce vital examples of
adjoints.%
The most important is the
following:

\begin{thm}
  \label{thm:adj:precomp-is-right-adj}
  If $\DeclVar{f}{\GM}{D \to C}$ and $D$ is small, then
  $\PSH{f} \defeq (f^\dagger)^* : \PSH{C} \to \PSH{D}$ is a right adjoint with left adjoint $f_!$.
\end{thm}
\begin{proof}
  For notational simplicity, we replace $C$ and $D$ with $\Modify[\OM]{C}$ and $\Modify[\OM]{D}$.
  By Theorem~\ref{thm:adj:pw-to-full-adj}, it suffices to assume $\DeclVar{X}{\GM}{D \to \Space}$ and to
  construct $f_!\prn{X} : C \to \Space$ along with a natural bijection
  $\Prod{Y} \Hom{f_!\prn{X}}{Y} \Equiv \Hom{X}{f^*\prn{Y}}$. This is an immediate consequence of
  \textcite[Theorem 2.41]{rijke:2020} after localizing the composite $\Sum{c : C} X\prn{C} \to D$
  against the map $\brc{0} \to \Int$.
\end{proof}

\begin{cor}
  The left adjoint $f_! \Adjoint \PSH{f}$ satisfies $f_! \circ \Yo \cong \Yo \circ f$.
\end{cor}

\begin{cor}
  $\Space$ is small cocomplete: $\Const : \Space \to \Space^C$ is a right adjoint with left
  adjoint $\Colim{}$ for small categories $\DeclVar{C}{\GM}{\Uni}$. Explicitly, if
  $\DeclVar{X}{\GM}{C \to \Space}$, then $\Colim{C} X = \Grpdify\prn{\Sum{c : C} X\prn{c}}$.
\end{cor}
\begin{rem}
  One can prove $\Space$ is complete (that $\Const \Adjoint \Lim{}$) by a result of
  \textcite{gratzer:2024}. In particular, they show $\Prod{c : C} X\prn{c} : \Space$ whenever
  $X : C \to \Space$ and $\DeclVar{C}{\GM}{\Uni}$. Corollary~\ref{cor:prelims:dua} and Lemma~\ref{lem:tw:hom-of-psh} then
  imply that $\Hom{A}{\Prod{c : C} X\prn{c}} \Equiv \Hom{\Con{const}\,A}{X}$.
\end{rem}

\begin{lem}
  \label{lem:adj:eval-at-pt-is-radj}
  If $\DeclVar{c}{\GM}{C}$ then $\PSH{c} : \PSH{C} \to \Space$ is a left adjoint.
\end{lem}
\begin{proof}
  For simplicity, we once more replace $C$ with $\Modify[\OM]{C}$ and invoke
  Theorem~\ref{thm:adj:pw-to-full-adj}. We then assume that we are given $\DeclVar{X}{\GM}{\Space}$ and
  define $c_*X$ to be $\lambda c'.\,X^{\Hom{c'}{c}}$. This is covariant in $c'$
  \parencite{gratzer:2024}. It then suffices to show that the map $X^{\Hom{c}{c}} \to X$ given by
  evaluation at the identity map induces an equivalence
  $\Hom[C \to \Space]{Y}{c_*X} \Equiv \Hom[\Space]{Y\prn{c}}{X}$. This, in turn, is a consequence of
  the fact that the map $Y\prn{c} \to \Grpdify\prn{\Sum{c' : C} \Hom{c'}{c} \times Y\prn{c'}}$ is an
  equivalence, so we may decompose the evaluation map as follows:
  \begin{align*}
    &\Hom[C \to \Space]{Y}{c_*X}
    \\
    &\Equiv \Prod{c' : C} Y\prn{c'} \to X^{\Hom{c'}{c}}
    \\
    &\Equiv \Prod{c' : C} \Hom{c'}{c} \times Y\prn{c'} \to X
    \\
    &\Equiv \Grpdify \prn{\Sum{c' : C} \Hom{c'}{c} \times Y\prn{c'}} \to X
    \\
    &\Equiv Y\prn{c} \to X \qedhere
  \end{align*}
\end{proof}

Combining Theorem~\ref{thm:adj:precomp-is-right-adj} and Lemma~\ref{lem:adj:eval-at-pt-is-radj}, we obtain the
following characterization of $f_!X$:

\begin{cor}
  \label{cor:adj:push-pull-formula}
  %
  If $\DeclVar{X}{\GM}{\PSH{C}}$, $\DeclVar{f}{\GM}{C \to D}$, and $\DeclVar{d}{\GM}{D}$ then we may
  explicitly identify
  $\prn{f_!X}\,d$ as the following type:
  \[
    \Grpdify \prn{\Sum{c : \Modify[\OM]{C}} X\prn{c} \times \Hom{f^\dagger c}{d}}
  \]
\end{cor}
\begin{proof}
  We first observe that $\PSH{d}f_!\prn{X} = \prn{f_!X}\,d$. Transposing, we have
  $\Hom{\PSH{d}f_!X}{Z} \Equiv \Hom{X}{\PSH{f}d_*Z}$ for every $\DeclVar{Z}{\GM}{\Space}$. We
  calculate $\Hom{X}{\PSH{f}d_*Z}$ using the definition of $d_*$ provided above:
  \begin{align*}
    &\Hom[\PSH{C}]{X}{f^*d_*Z}
    \\
    &\Equiv \Prod{\DeclVar{c}{}{\Modify[\OM]{C}}} X\prn{c} \to \Hom{f^\dagger c}{d} \to Z
    \\
    &\Equiv \prn{\Sum{\DeclVar{c}{}{\Modify[\OM]{C}}} X\prn{c} \times \Hom{f^\dagger c}{d}} \to Z
  \end{align*}
  Therefore $\PSH{d}f_!X$ satisfies the universal property of
  $\Grpdify \prn{\Sum{\DeclVar{c}{}{\Modify[\OM]{C}}} X\prn{c} \times \Hom{f^\dagger c}{d}}$
\end{proof}

The following lemma does not require Theorem~\ref{thm:adj:pw-to-full-adj}, but is merely a consequence of
manipulating natural transformations:
\begin{lem}
  \label{lem:adj:comp}
  If $f : C \to D$ is an adjoint so is $f_* : C^A \to D^A$.
\end{lem}

\begin{cor}
  \label{cor:adj:limits-are-pointwise}
  If $C$ is (co)complete so is $C^D$ and (co)limits are computed pointwise. In particular, $\PSH{C}$
  is (co)complete.
\end{cor}
\begin{cor}
  \label{cor:adj:yo-preserves-limits}
  The Yoneda embedding preserves all limits.
\end{cor}
\begin{proof}
  If $\DeclVar{F}{\GM}{I \to C}$ and $\Lim{} F$ exists, then functoriality of $\Yo$ induces a map
  $\Yo{\Lim{} F} \to \Lim{}\prn{\Yo \circ F}$, so it suffices to check that this map is
  invertible at all $\DeclVar{c}{\GM}{C}$. Unfolding, we must argue that
  $\Hom{c}{\Lim{} F} \Equiv \Lim{} \Hom{c}{F}$ is an equivalence, but this is immediate by
  Lemma~\ref{lem:tw:hom-of-psh}.
\end{proof}

Finally, we show the full subcategories $\Space_{\le n}$ of $\Space$ defined by $n$-truncated types
form reflective subcategories of $\Space$. The idea is simple: use the truncation HITs. However, it
is not automatic that they restrict to $\Trunc{-}_{n} : \Space \to \Space_{\le n}$. We prove this
alongside with the reflectivity of $\Space_{\le n}$ using Theorem~\ref{thm:adj:pw-to-full-adj}:
\begin{cor}
  The inclusion $\Space_{\le n} \to \Space$ is a right adjoint.
\end{cor}
\begin{cor}
  $\Space_{\le n}$ is (co)complete.
\end{cor}
The same methodology applies to the subcategory of modal types associated to an idempotent
monad~\parencite{rijke:2020}.

\begin{exa}[Isbell conjugation]
  If $\DeclVar{C}{\GM}{\Uni}$, then the Isbell conjugation map $\phi$ is a left adjoint:
  \begin{align*}
    &\phi : \PSH{C} \to \Modify[\OM]{C \to \Space}
    \\
    &\phi\prn{X} = \MkMod[\OM]{\lambda c.\,\Phi\prn{X,\Yo{c}}}
  \end{align*}
\end{exa}

\subsection{The universal property of presheaf categories}
\label{sec:adj:psh}

Next, we generalize Theorem~\ref{thm:adj:precomp-is-right-adj} to show that if $\DeclVar{f}{\GM}{C \to E}$
where $C$ is a small category and $E$ is a cocomplete category, then
$\Phi\prn{f^\dagger\prn{-},-} : E \to \PSH{C}$ is a right adjoint loosely following the argument
given by \textcite{cisinski:2019}. We begin with a few general lemmas. In what follows, fix $C$
and $E$ as above.

First, as a corollary of the proof of Theorem~\ref{thm:adj:precomp-is-right-adj}:
\begin{lem}
  The colimit of $\Yo : C \to \PSH{C}$ is $\ObjTerm{\PSH{C}} = \lambda \_.\,\ObjTerm{}$.
\end{lem}

From the above, and further inspection of colimits, we are able to derive a result of independent
interest: Every presheaf is the colimit of representable presheaves.
\begin{lem}[Density of $\Yo$]
  \label{lem:adj:density}
  If $\DeclVar{X}{\GM}{\PSH{C}}$, then
  $X \cong \Colim{\Modify[\OM]{\TotalType{X}}} \Yo \circ \pi^\dagger$,
  where $\TotalType X=\Sum{c:\Modify[\OM]C}X\prn c$.
\end{lem}
\begin{proof}
  We begin with the following computation where $\pi : \TotalType{X} \to \Modify[\OM]{C}$
  and $\pi^\dagger_! : \Space^{\TotalType{X}} \to \PSH C$:
  \[
    \pi^\dagger_!\ObjTerm{}
    \cong
    \pi^\dagger_!\prn{\Colim{\Modify[\OM]{\TotalType{X}}} \Yo}
    \cong
    \Colim{\Modify[\OM]{\TotalType{X}}} \pi^\dagger_! \circ \Yo
    \cong
    \Colim{\Modify[\OM]{\TotalType{X}}} \Yo \circ \pi^\dagger
  \]
  We have used the fact that $\pi^\dagger_!$, a left adjoint, commutes with
  colimits~\parencite{bardomiano:2021}.
  To show $\pi^\dagger_!\ObjTerm{} \cong X$, we note that for all $Z : \PSH{C}$:
  \begin{align*}
    \Hom{\pi^\dagger_!\ObjTerm{}}{Z}
    &\Equiv \Hom[{\TotalType{X}} \to \Space]{\ObjTerm{}}{Z \circ \pi}
    \\
    &\Equiv \Prod{\prn{c, x} : \Sum{c : \Modify[\OM]{C}} X\prn{c}} Z\prn{c}
    \\
    &\Equiv \Prod{c : \Modify[\OM]{C}} X\prn{c} \to Z\prn{c}
    \\
    &\Equiv \Hom{X}{Z}
  \end{align*}
  The conclusion now follows from the Yoneda lemma.
\end{proof}

\begin{lem}
  $\Nv_f = \Phi\prn{f^\dagger\prn{-},-} : E \to \PSH{C}$ is a right adjoint.
\end{lem}
\begin{proof}
  We will prove that $\Nv_f$ is a pointwise right adjoint. Accordingly, fixing
  $\DeclVar{X}{\GM}{\PSH{C}}$ we must construct $\DeclVar{e}{}{\Modify[\OM]{E}}$ such that
  $\Phi\prn{e,-} \cong \Phi\prn{\MkMod[\OM]{X},\Nv_f\prn{-}}$. Since
  $X \cong \Colim{\Modify[\OM]{\TotalType{X}}} \Yo \circ \pi^\dagger$ and $E$ is cocomplete, by the
  dual of Corollary~\ref{cor:adj:yo-preserves-limits} it suffices to assume
  $\MkMod[\OM]{X} = \Yo^\dagger\prn{c}$ with $\DeclVar{c}{}{\Modify[\OM]{C}}$.\footnote{Note the
    lack of $\GM$-annotation here: we must ensure that we are functorial in $c$ in order to obtain a
    \emph{diagram} in $E$.} Finally, take $e = f^\dagger\prn{c}$ and
  $\Phi\prn{f^\dagger\prn{c},-} \cong \Phi\prn{\Yo^\dagger\prn{c},\Nv_f\prn{-}}$ by
  Theorem~\ref{thm:tw:strong-yoneda}.
\end{proof}

We are now able to prove, as promised, the universal property of $\PSH{C}$. If we write
$\CoCont{\PSH{C}}{E}$ for the full subcategory of $\PSH{C} \to E$ spanned by functors preserving
all colimits, then $\Yo^* : \CoCont{\PSH{C}}{E} \to \prn{C \to E}$ is an equivalence. To prove this,
we essentially argue that there is a map sending $f$ to the left adjoint to $\Nv_f$ and that this is
the inverse to $\Yo^*$.

\begin{restatable}{thm}{universalpropofpsh}
  \label{thm:adj:cc}
  $\Yo^* : \CoCont{\PSH{C}}{E} \to \prn{C \to E}$ is an equivalence.
\end{restatable}
\begin{proof}
  We use Theorem~\ref{thm:prelims:ff-ess-surj-to-equiv}. If $\DeclVar{f}{\GM}{C \to E}$, then
  $f_! : \PSH{C} \to E$ satisfies $f_! \circ \Yo = f$ and so $\Yo^*$ is essentially surjective:
  \[
    \Phi\prn{\prn{f_!\circ \Yo}^\dagger\prn{-},-} \cong \Phi\prn{\Yo^\dagger\prn{-},\Nv_f\prn{-}}
    \cong \Nv_f = \Phi\prn{f^\dagger\prn{-},-}
  \]

  Moreover, if $\DeclVar{F}{\GM}{\CoCont{\PSH{C}}{E}}$, then $\prn{F\circ\Yo}_! \cong F$, so that
  $\Yo^*$ is a bijection on $\GM$-elements. Let us write $f = F \circ \Yo$. We first construct a
  comparison map $\Hom{f_!}{F}$ by constructing a natural transformation
  $\Hom{\ArrId{}}{\Nv_{f}\prn{F\prn{-}}}$. Currying, this is equivalent to constructing a natural
  transformation between maps $\Modify[\OM]{C} \times \PSH{C} \to \Space$ and, in this form,
  $\ArrId{}$ is given by evaluation $\epsilon$ and $\Nv_F\prn{F\prn{-}}$ is
  $\Phi\prn{f\prn{-},F\prn{-}}$. We can replace $\epsilon$ with $\Phi\prn{\Yo\prn{-},-}$ by
  Theorem~\ref{thm:tw:strong-yoneda} and
  $\Phi\prn{f\prn{-},F\prn{-}} = \Phi\prn{F\prn{\Yo\prn{-}}, F\prn{-}}$ by definition. Accordingly,
  the relevant map is supplied by $\Phi_F$. It is routine to check that this is pointwise an
  equivalence by Lemma~\ref{lem:adj:density}.

  Finally, we now show that $\Yo^*$ is fully faithful. To show that it is fully faithful, we must
  show that if $\DeclVar{f,g}{\GM}{C \to E}$, then $\Hom{f_!}{g_!} \cong \Hom{f}{g}$. Both sides are
  groupoids, so it suffices to consider $\GM$-annotated elements. If
  $\DeclVar{\alpha}{\GM}{\Hom{f}{g}}$, then by transposing we may regard $\alpha$ as an element of
  $\Modify[\GM]{C \to E^\Int}$ and the previous observation ensures that this type is equivalent to
  $\Modify[\GM]{\CoCont{\PSH{C}}{E^\Int}}$ which yields the desired conclusion after transposing.
\end{proof}

\section{The theory of Kan extensions}
\label{sec:kan}
A unifying concept in category theory are \emph{Kan extensions}, which are universal extensions of
functors along functors on the same domain. Mac Lane, one of the founders of category theory,
famously stated: ``The notion of Kan extensions subsumes all the other fundamental concepts of
category theory,'' such as (co)limits and adjunctions~\parencite{maclane:1978,riehl:2014}.

\begin{defi}[Kan extensions]
  Given a map $f : C \to D$ and a category $E$, the left (right) Kan extension $\Lan{f}$
  ($\Ran{f}$) is the left (right) adjoint to $f^* : E^D \to E^C$.
\end{defi}

While the definition makes sense in general, to use the results of the previous sections, we shall
assume $\DeclVar{f}{\GM}{C \to D}$ and $\DeclVar{E}{\GM}{\Uni}$.  In Section~\ref{sec:kan:existence} we
show that Kan extensions exist whenever $E$ is (co)complete and in
Sections~\ref{sec:kan:cofinal} and \ref{sec:kan:quillen-a} we put this to work by deducing two important results:
Quillen's theorem A and the properness of cocartesian fibrations. Our arguments for the existence of
Kan extensions and Quillen's theorem A adapt the (model-agnostic) $\infty$-categorical arguments of
\textcite{ramzi:bousfield-kan}.

\subsection{Existence and characterization of Kan extensions}
\label{sec:kan:existence}

We can prove that Kan extensions can be computed in an expected way. For $\DeclVar{d}{}{D}$, we
write $\SLICE{C}{d} \defeq C \times_D \SLICE{D}{d}$ and
$\COSLICE{C}{d} \defeq C \times_D \COSLICE{D}{d}$.
We assume that $C$ and $D$ are both small so each $\SLICE Cd$ is also small.
By Theorem~\ref{thm:adj:precomp-is-right-adj} and Lemma \ref{lem:adj:comp}:
\begin{lem}
  \label{lem:kan:lan-exists-for-psh}
  If $E = \PSH{A}$ for some category $\DeclVar A\GM\Uni$, then $\Lan{f}$ exists.
  Moreover, if $\DeclVar{X}{\GM}{C \to E}$ and $\DeclVar{d}{\GM}{D}$, then
  $\Lan{f}\,X\,d = \Colim{} \prn{\SLICE{C}{d} \to C \to E} = \Grpdify\prn{\Sum{\prn{c,\_} : \SLICE{C}{d}} X\prn{c}}$.
\end{lem}
This yields more generally:
\begin{thm}
  \label{thm:kan:left-kan-exists}
  If $E$ is cocomplete, then $\Lan{f}$ exists, and if $\DeclVar{X}{\GM}{C \to E}$,
  $\DeclVar{d}{\GM}{D}$, then $\Lan{f}\,X\,d = \Colim{} \prn{\SLICE{C}{d} \to C \to E}$.
\end{thm}
\begin{proof}
  It suffices to argue that precomposition is pointwise a right adjoint and so we fix
  $\DeclVar{X}{\GM}{C \to E}$. By Theorem~\ref{thm:adj:cc}, we may view $X$ as the composition
  $\bar X \circ \Yo$, where $\bar X : \PSH{C} \to E$ is the left adjoint to $\Nv_X$.
  Next, we observe by~Lemma~\ref{lem:kan:lan-exists-for-psh} that $\Yo: C \to \PSH C$
  admits an extension to $D$ along $f$, namely $\Lan{f} \Yo : D \to \PSH C$,
  and we claim that
  $\bar{X} \circ \Lan{f} \Yo$ is our desired extension of $f$.
  Fixing $Z : D \to E$, we calculate:
  \begin{align*}
    \Hom[D \to E]{\bar{X} \circ \Lan{f} \Yo}{Z}
    &\Equiv \Hom[D \to \PSH{C}]{\Lan{f} \Yo}{\Nv_X \circ Z}
    \\
    &\Equiv \Hom[C \to \PSH{C}]{\Yo}{\Nv_X \circ Z \circ f}
    \\
    &\Equiv \Hom[C \to E]{\bar X\circ\Yo}{Z \circ f}
    \\
    &= \Hom[C \to E]{X}{Z \circ f}
  \end{align*}
  The expected colimit formula continues to hold as a consequence of
  Lemma~\ref{lem:kan:lan-exists-for-psh} and the cocontinuity of $\bar{X}$.
\end{proof}

By duality, we obtain the following variant:
\begin{thm}
  \label{thm:kan:right-kan-exists}
  If $E$ is complete, then $\Ran{f}$ exists and is specified by the dual limit formula:
  $\Ran{f}\,X\,d = \Lim{} \prn{\COSLICE{C}{d} \to C \to E}$.
\end{thm}

\subsection{Final and initial functors}
\label{sec:kan:cofinal}

It is frequently useful to show that the limit of a complex diagram $D$ can be calculated by first
restricting to a simpler diagram using $f : C \to D$ and calculating the limit there \eg{},
restricting from $\mathbb{Z}$ to $\mathbb{Z}_{\le 0}$. When this approach is valid, $f$ is said to
be initial:
\begin{defi}
  A functor $\DeclVar{f}{\GM}{C \to D}$ is \emph{initial} if for every
  $\DeclVar{X}{\GM}{D \to \Space}$ the map $\Lim{D} X \to \Lim{C} X \circ f$ is an
  equivalence. A map is \emph{final} if its opposite is initial.
\end{defi}
While this definition is asymmetrical in its treatment of initiality and finality,
we shall restore the symmetry as a consequence of Quillen's Theorem~A in the next section,
see~Corollary~\ref{cor:kan:right-cofinal-char}.

Recall that $\Lim{D} X = \Prod{d : D} X\prn{d}$ and so the definition of
initiality equivalently states that the
restriction map $\prn{\Prod{d : D} X\prn{d}} \to \prn{\Prod{c : C} X\prn{f\prn{c}}}$ is an
equivalence.

\begin{exa}
  The $\brc{0}$/$\brc{1}$ inclusion ${\Unit \to \Int}$ is initial/final.
\end{exa}

\begin{lem}
  \label{lem:kan:left-cofinal-preserves-limit}
  If $\DeclVar{f}{\GM}{C \to D}$ is initial and
  $\DeclVar{X}{\GM}{D \to E}$, then $\Lim{C} \prn{X \circ f}$ and $\Lim{D} X$ both exist
  whenever either exists and are canonically isomorphic.
\end{lem}
\begin{proof}
  By Corollary~\ref{cor:adj:yo-preserves-limits}, we replace $E$ with $\PSH{E}$ and by
  Corollary~\ref{cor:adj:limits-are-pointwise} we reduce to $\Space$ where the result is immediate.
\end{proof}

\begin{restatable}{lem}{locop}
  \label{lem:app:localization-op}
  If $\DeclVar{C}{\GM}{\Uni}$, then $\Grpdify{C} \Equiv \Grpdify{\Modify[\OM]{C}}$.
\end{restatable}
\begin{proof}
  We observe that $\Grpdify{C} \Equiv \Modify[\GM]{\Grpdify{C}}$ and likewise for
  $\Modify[\OM]{C}$. Accordingly, we note that:
  \begin{gather*}
    \Modify[\GM]{
      \Modify[\OM]{C} \to \Modify[\GM]{X}
    }
    \Equiv
    \Modify[\GM]{
      \Grpdify{\Modify[\OM]{C}} \to \Modify[\GM]{X}
    }
    \\
    \Modify[\GM]{
      C \to \Modify[\GM]{X}
    }
    \Equiv
    \Modify[\GM]{
      \Grpdify{C} \to \Modify[\GM]{X}
    }
    \\
    \Modify[\GM]{
      C \to \Modify[\GM]{X}
    }
    \Equiv
    \Modify[\GM]{
      \Modify[\OM]{C} \to \Modify[\GM]{X}
    }
  \end{gather*}
  Finally, the result follows from a simple Yoneda argument.
\end{proof}

\begin{lem}
  \label{lem:kan:localization-is-cofinal}
  For every $C$, the canonical map $C \to \Grpdify{C}$ is both initial and final.
\end{lem}
\begin{proof}
  By Lemma~\ref{lem:app:localization-op}, it suffices to argue that this map is initial. To this
  end, we must show the following map to be an equivalence for every $X : \Grpdify{C} \to \Space$:
  \[
    \prn{\Prod{d : \Grpdify{C}} X\prn{d}} \to \prn{\Prod{c : C} X\prn{\eta\prn{c}}}
  \]
  However, $X\prn{d}$ is discrete for every $d : \Grpdify{C}$ and so this is simply the universal
  property of $\Grpdify$.
\end{proof}
\begin{cor}
  \label{lem:kan:limit-over-contractible}
  If $\Grpdify{C} = \ObjTerm{}$, then $\Lim{C} A = A$ for $A : \Space$.
\end{cor}

\subsection{Quillen's Theorem A}
\label{sec:kan:quillen-a}

Our next goal is to prove the $\infty$-categorical version of Quillen's theorem A. Unlike
traditional proofs, we follow \textcite{ramzi:bousfield-kan} and rely on having already established
the basic apparatus of Kan extensions to simplify our argument.

\begin{defi}
  A functor $\DeclVar{f}{\GM}{C \to D}$ is Quillen final if
  $\Grpdify(\COSLICE{C}{d}) \Equiv \Unit$ for all $\DeclVar{d}{\GM}{D}$
\end{defi}
\begin{thm}
  \label{thm:kan:theorem-a}
  A functor $\DeclVar{f}{\GM}{C \to D}$ is final if and only if
  it is Quillen final.
\end{thm}

\begin{rem}
  This result shows that, in particular, finality doesn't depend on the particular universe $\Space$
  chosen.
\end{rem}

\begin{lem}
  \label{lem:kan:theorem-a-for-psh}
  If $\DeclVar{f}{\GM}{C \to D}$ is Quillen final and $\DeclVar{X}{\GM}{D \to \PSH{A}}$,
  then ${\Colim D} X \Equiv {\Colim C} X \circ f$.
\end{lem}
\begin{proof}
  This statement is pointwise, so we quickly reduce to $\Space$ instead of $\PSH{A}$. In this
  situation, we wish to show that the following commutes:
  \[
    \begin{tikzcd}
      \Space^D \ar[rr,"f^*"]\ar[dr,"{\mathrm{colim}_D}"'] && \Space^C\ar[dl,"{\mathrm{colim}_C}"] \\
      & \Space &
    \end{tikzcd}
%
%
%
%
%
%
%
%
  \]
  Note that all three morphisms are left adjoints,
  and so it suffices to compare their right adjoints:
  the constant functors $\Delta_C$ and $\Delta_D$,
  along with the right Kan extension $\Ran f$.
  We next note that there is at least a comparison map
  $\Delta_D \to \Ran{f} \circ \Delta_C$ given by transposing the
  identity map  $f^* \circ \Delta_D \to \Delta_C$.
  We must argue that this map is pointwise invertible,
  and so we reduce to considering $\DeclVar{X}{\GM}{\Space}$
  and $\DeclVar{d}{\GM}{D}$,
  and we must show the following, using~Theorem~\ref{thm:kan:right-kan-exists}:
  $X \Equiv \Lim{\COSLICE{C}{d}} X$.
  This now follows from our assumption and~Lemma~\ref{lem:kan:limit-over-contractible}.
\end{proof}

\begin{lem}
  \label{lem:kan:theorem-a}
  If $\DeclVar{f}{\GM}{C \to D}$ is Quillen final, $E$ a cocomplete category, and
  $\DeclVar{X}{\GM}{D \to E}$, then ${\Colim D} X \Equiv {\Colim C} X \circ f$.
\end{lem}
\begin{proof}
  We reduce to the case where $E = \PSH{D}$ (and therefore Lemma~\ref{lem:kan:theorem-a-for-psh}) by
  factoring $X$ as $\bar{X} \circ \Yo$ and noting that $\bar{X}$ preserves colimits by construction.
\end{proof}

\begin{proof}[Proof of Theorem~\ref{thm:kan:theorem-a}]
  To see that Quillen finality implies finality, we apply
  Lemma~\ref{lem:kan:theorem-a} to $\Modify[\OM]{\Space}$,
  and calculate:
  \[
    \Lim{\Modify[\OM]D} X
    \Equiv {\Colim D} X^\dagger
    \Equiv {\Colim C} X^\dagger \circ f
    \Equiv \Lim{\Modify[\OM]C} X \circ f^\dagger
  \]
  For the reverse, suppose that $f$ is final. We note that by the dual of
  Lemma~\ref{lem:kan:left-cofinal-preserves-limit} (again applied to $\Modify[\OM]{\Space}$), the
  canonical map $\Colim{C} X \circ f \to \Colim{D} X$ is an equivalence for any
  $\DeclVar{X}{\GM}{D \to \Space}$.  Fix $\DeclVar d\GM D$ and choose
  $X = \hom(d,-) = \Phi\prn{\MkMod[\OM]d,-}$ such that the colimits in question are precisely
  $\Grpdify{\COSLICE{D}{d}}$ and $\Grpdify{\COSLICE{C}{d}}$,
  using~Theorem~\ref{thm:adj:precomp-is-right-adj}.  This completes the proof since
  $\Grpdify{\COSLICE{D}{d}} = \ObjTerm{}$.
\end{proof}

\begin{cor}
  \label{cor:kan:right-cofinal-char}
  A functor $\DeclVar{f}{\GM}{C \to D}$ is final if and only if, for every
  $\DeclVar{X}{\GM}{D \to \Space}$ the map $\Colim{D} X \to \Colim{C} X \circ f$
  is an equivalence.
\end{cor}
This restores the symmetry between initial and final functors, as promised.  We offer another
symmetric definition of initiality and finality, informed by \textcite[Ch.~8]{cisinski:2024}.
\begin{defi}[Covariant equivalences]
  Fix $\DeclVar{p}{\GM}{C \to A}$ and $\DeclVar{q}{\GM}{D \to A}$ between categories $A,C,D$. Let
  $\DeclVar{f}{\GM}{C \to D}$ be a fibered map as follows:
  \[
    \begin{tikzcd}
     C && D \\
     & A
     \arrow["f", from=1-1, to=1-3]
     \arrow["p"', from=1-1, to=2-2]
     \arrow["q", from=1-3, to=2-2]
    \end{tikzcd}
  \]
  We call $f$ a \emph{covariant equivalence} if for all families $\DeclVar{X}{\GM}{A \to \Space}$
  reindexing gives rise to an equivalence, \ie,:
  \[
    f^* : \prn{\DelimMin{1}\Prod{a:A} D_a \to X_a} \to \prn{\DelimMin{1}\Prod{a:A} C_a \to X_a}
  \]
  Dually, $f$ is called \emph{contravariant equivalence} precomposition with respect to all contravariant families is an equivalence.
\end{defi}

\begin{lem}\label{lem:cov-eq-comp}
  Let $f$ as below be a covariant equivalence with respect to $p$ and $q$. Then, for any functor
  $\DeclVar{r}{\GM}{B \to A}$ it is also a covariant equivalence with respect to $rp$ and $rq$:
  \[
    \begin{tikzcd}
      C && D \\
      & A \\
      & B
      \arrow["f", from=1-1, to=1-3]
      \arrow["p"', from=1-1, to=2-2]
      \arrow["q", from=1-3, to=2-2]
      \arrow["r", from=2-2, to=3-2]
    \end{tikzcd}
  \]
\end{lem}

\begin{proof}
  We get the following induced square:
  \[
    \begin{tikzcd}
      {\big( \prod_{b:B} D_b \to X_b \big)} && {\big( \prod_{b:B} C_b \to X_b \big)} \\
      {\big( \prod_{a:A} D_a \to X_{r(a)} \big)} && {\big( \prod_{a:A} C_a \to X_{r(a)} \big)}
      \arrow["{f^*}", "\simeq"', from=1-1, to=1-3]
      \arrow["\simeq"', from=1-1, to=2-1]
      \arrow["\simeq", from=1-3, to=2-3]
      \arrow["{f^*}"', from=2-1, to=2-3]
    \end{tikzcd}
  \]
  The upper horizontal map is an equivalence by the preconditions. The goal is to show that the
  lower horizontal map is an equivalence, too. But this follows from $3$-for-$2$ for equivalences.
\end{proof}

\begin{lem}[Characterizations of initiality]
  \label{lem:kan:characterization-initiality}
  Let $\DeclVar{f}{\GM}{C \to D}$ be a functor. Then the following are equivalent:
  \begin{enumerate}[label=\textup{(\arabic*)}]
  \item\label{it:prop:char-left-cof-i} $f$ is initial.
  \item\label{it:prop:char-left-cof-ii} Let $\DeclVar{X}{\GM}{A \to \Space}$ be a family with
    associated left fibration $\DeclVar{\pi}{\GM}{\TotalType{X} \to A}$. Then any square of the
    following form has a filler $\overline{\varphi}$, uniquely up to homotopy:
    \[
      \begin{tikzpicture}
        \node (nw) {$C$};
        \node[below = 1.5cm of nw] (sw) {$D$};
        \node[right = 4cm of nw] (ne) {$\TotalType{X}$};
        \node[right = 4cm of sw] (se) {$A$};
        \path[->] (nw) edge node[above] {$\varphi$} (ne);
        \path[->] (nw) edge node[left] {$f$} (sw);
        \path[->] (ne) edge node[right] {$\pi$} (se);
        \path[->] (sw) edge node[below] {$\alpha$} (se);
        \path[->, unique exists] (sw) edge node[upright desc] {$\overline{\varphi}$} (ne);
      \end{tikzpicture}
    \]
  \item\label{it:prop:char-left-cof-iii} For any family $\DeclVar{X}{\GM}{A \to \Space}$ the
    following square is a pullback:
    \[
      \DiagramSquare{
        nw = \TotalType{X}^D,
        ne = \TotalType{X}^C,
        sw = A^D,
        se = A^C,
        nw/style = pullback,
        height = 1.5cm,
      }
    \]
  \item\label{it:prop:char-left-cof-iv} $f$ is a covariant equivalence with respect to any
    $\DeclVar{\alpha}{\GM}{D \to A}$.
  \end{enumerate}
  The analogous characterization holds for contravariant equivalences and final functores.
\end{lem}
\begin{proof}
  Conditions~\ref{it:prop:char-left-cof-ii} and~\ref{it:prop:char-left-cof-iii} are readily seen
  to be equivalent by commuting $\Prod{}$ and $\Sum{}$. Condition~\ref{it:prop:char-left-cof-iv}
  unfolds to the following: for any $\DeclVar{X}{\GM}{A \to \Space}$ reindexing along $f$ is an
  equivalence, namely
  \[
    f^* :
    \prn{\Prod{a:A} D_a \to X_a}
    \xrightarrow{\sim}
    \prn{\Prod{a:A} (\Sum{d:D_a} C_{a,d}) \to X_a}
  \]
  This, again, is readily seen to be equivalent to~\ref{it:prop:char-left-cof-ii}.

  We turn to the implication
  $\text{\ref{it:prop:char-left-cof-iv}} \implies \text{\ref{it:prop:char-left-cof-i}}$.
  But this is clear,
  since~\ref{it:prop:char-left-cof-i} says that $f$ is a covariant equivalence with respect to
  itself and $\ArrId{D}$.

  For the converse direction
  $\text{\ref{it:prop:char-left-cof-i}} \implies \text{\ref{it:prop:char-left-cof-iv}}$
  we use the insight
  just made together with Lemma~\ref{lem:cov-eq-comp}.
\end{proof}

The following alternative characterization is also often useful:

\begin{restatable}{lem}{charcofinal}
  \label{prop:kan:charcofinal}
  $\DeclVar{f}{\GM}{C \to D}$ is initial \textup(resp. final\textup) if and only if for every
  covariant \textup(resp. contravariant\textup) family $\DeclVar{\pi}{\GM}{X \to Y}$, $f$ is left
  orthogonal to $\pi$, \ie{}, $\IsEquiv\prn{X^D \to X^C \times_{Y^C} Y^D}$.
\end{restatable}
\begin{proof}
  Immediate by Proposition~\ref{it:prop:char-left-cof-iv} for the initial case
  and by duality and~Corollary~\ref{cor:kan:right-cofinal-char} for the final case.
\end{proof}

\begin{cor}
  \label{cor:initial-covariant-ofs}
  There is an orthogonal factorization system in the sense of \cite{rijke:2020} with the left class given by the initial functors and the right class by covariant fibrations.
\end{cor}

As another consequence we get the dual of~Theorem~\ref{thm:kan:theorem-a}:
\begin{cor}
  \label{thm:kan:theorem-a-dual}
  A functor $\DeclVar{f}{\GM}{C \to D}$ is initial if and only if
  $\Grpdify(\SLICE{C}{d}) \Equiv \Unit$ for all $\DeclVar{d}{\GM}{D}$
  \textup(\emph{Quillen initial}\textup).
\end{cor}

We demonstrate the utility of Theorem~\ref{thm:kan:theorem-a} by giving a new and far simpler
proof that cocartesian fibrations are \emph{proper}.
\begin{defi}
  \label{def:proper-maps}
  A functor $\DeclVar{\pi}{\GM}{E \to B}$ between categories is \emph{proper} if for all pullbacks
  (of $\GM$-functors) of the following form, $v$ is final if $u$ is final:
  \[
    \begin{tikzpicture}[diagram]
      \SpliceDiagramSquare{
        height = 1.35cm,
        ne = E,
        se = B,
        nw = E',
        sw = B',
        east = \pi,
        west = \pi',
        nw/style = pullback,
      }
      \node [pullback, left = of nw] (A) {$E''$};
      \node [left = of sw] (B) {$B''$};
      \path[->] (A) edge node[above] {$v$} (nw);
      \path[->] (B) edge node[below] {$u$} (sw);
      \path[->] (A) edge (B);
    \end{tikzpicture}
  \]
  We call $\pi$ \emph{smooth} if $\pi^\dagger : \Modify[\OM]{E} \to \Modify[\OM]{B}$ is proper.
\end{defi}
\begin{lem}
  Smooth and proper functors are closed under composition and pullback.\footnote{The definition of
    properness is formulated specifically to bake in the latter.}
\end{lem}
\begin{thm}
  \label{thm:kan:cocartesian-is-proper}
  Cartesian fibrations are smooth and cocartesian fibrations are proper.
\end{thm}
\begin{proof}
  It suffices to treat the proper case. Fix a cocartesian fibration $\DeclVar{\pi}{\GM}{E \to B}$
  and note that since cocartesian fibrations are stable under pullbacks, it suffices show to that $v$ is
  final in the following pullback diagram if $u$ is final:
  \[
    \DiagramSquare{
      height = 1.5cm,
      width = 3cm,
      nw = { A \times_B E},
      sw = A,
      ne = E,
      se = B,
      nw/style = pullback,
      north = v,
      south = u,
      east = \pi,
    }
  \]
  We now use Theorem~\ref{thm:kan:theorem-a}. For $\DeclVar{e}{\GM}{E}$ we compute the fiber:
  \begin{align*}
    & (A \times_B E) \times_E \COSLICE{E}{e}
    \\
    &\quad
    \Equiv A \times_B \COSLICE{E}{e}
    \\
    &\quad
    \Equiv A \times_B \prn{\Sum{\DeclVar{b'}{}{B},\DeclVar{f}{}{\Hom{\pi\prn{e}}{b'}}}{\prn{E_{b'}}^{\Int}}}
    && \text{$\pi$ is cocartesian}
    \\
    &\quad
    \Equiv \Sum{\DeclVar{(a,f)}{}{A \times_B \COSLICE{B}{\pi\prn{e}}}}{\COSLICE{\prn{E_{u\prn{a}}}}{f_!e}}
  \end{align*}
  Applying $\Grpdify$ to each fiber yields $\Grpdify{\COSLICE{\prn{E_{u\prn{a}}}}{f_!e}} \Equiv \Unit$ (as
  coslices have initial elements) and $\Grpdify(A \times_B \COSLICE{B}{\pi\prn{e}}) \Equiv \Unit$
  since $u$ is final by assumption. This implies that applying $\Grpdify$ to the entire
  $\Sum{}$-type produces $\Unit$~\parencite{rijke:2020}.
\end{proof}

\begin{cor}
  If $\DeclVar{\pi}{\GM}{E \to B}$ is cocartesian and $\DeclVar{X}{\GM}{E \to D}$, then the left Kan
  extension $\Lan{\pi} X$ sends $\DeclVar{b}{\GM}{B}$ to $\Colim{}\prn{E_b \to E \to D}$.
\end{cor}

\subsection{Smooth and proper base change}

We want to show that smooth and proper functors satisfy the \emph{Beck--Chevalley condition}. We follow \cite[Section 8.4]{cisinski:2024}; see also \cite{anel:2024} for a general discussion.

First, using the initial-covariant factorization (see \ref{cor:initial-covariant-ofs}) of a functor with small fibers we can compute the action of precomposition for (co)presheaves:
\begin{prop}[\protect{\cite[Theorem~8.1.18]{cisinski:2024}}]
  \label{prop:action-postcomp}
  Let $\DeclVar{u}{\GM}{A \to B}$ be a functor with small fibers. Then the left adjoint $\DeclVar{u_! \dashv u^*}{\GM}{\Space^B \to \Space^A}$ acts as follows: for $\DeclVar{F}{\GM}{A \to \Space}$, if
  \[\begin{tikzcd}
	E & {\Space_*} \\
	A & \Space
	\arrow[from=1-1, to=1-2]
	\arrow["\phi"', from=1-1, to=2-1]
	\arrow["\lrcorner"{anchor=center, pos=0.125}, draw=none, from=1-1, to=2-2]
	\arrow[from=1-2, to=2-2]
	\arrow["F", from=2-1, to=2-2]
  \end{tikzcd}\]
  denote by $E \stackrel{j}{\to} X \stackrel{\psi}{\to} B$ the initial-covariant factorization of $u \circ \phi$. Then $u_!(F)$ is the straightening of $\psi$:
  \[\begin{tikzcd}
	E & X & {\Space_*} \\
	A & B & \Space
	\arrow["j", from=1-1, to=1-2]
	\arrow["\phi"', from=1-1, to=2-1]
	\arrow[from=1-2, to=1-3]
	\arrow["\psi", from=1-2, to=2-2]
	\arrow["\lrcorner"{anchor=center, pos=0.125}, draw=none, from=1-2, to=2-3]
	\arrow[from=1-3, to=2-3]
	\arrow["u"', from=2-1, to=2-2]
	\arrow["{u_!(F)}"', from=2-2, to=2-3]
\end{tikzcd}\]
In particular, the unit $\DeclVar{\eta}{\GM}{F \to u^*u_!F}$ corresponds under directed univalence to the map $\widetilde{\eta}$:
\[\begin{tikzcd}
	E & Y & X \\
	& A & B
	\arrow["{\widetilde{\eta}}", dashed, from=1-1, to=1-2]
	\arrow["j"{description}, curve={height=-20pt}, from=1-1, to=1-3]
	\arrow["\phi"{description}, from=1-1, to=2-2]
	\arrow[from=1-2, to=1-3]
	\arrow["{u^*\psi}", from=1-2, to=2-2]
	\arrow["\lrcorner"{anchor=center, pos=0.125}, draw=none, from=1-2, to=2-3]
	\arrow["\psi", from=1-3, to=2-3]
	\arrow["u", from=2-2, to=2-3]
\end{tikzcd}\]
\end{prop}

\begin{proof}
  We want to show that an equivalence is given by the map
  \[ \DeclVar{\lambda \alpha.u^*\alpha \circ \eta}{\GM}{\hom_{B \to \Space}(u_!F, G) \to \hom_{A \to \Space}(F, u^*G)}.\]
  We write the unstraightening of $G$ as
  \[\begin{tikzcd}
	Z & {\Space_*} \\
	B & \Space
	\arrow[from=1-1, to=1-2]
	\arrow["\gamma"', from=1-1, to=2-1]
	\arrow["\lrcorner"{anchor=center, pos=0.125}, draw=none, from=1-1, to=2-2]
	\arrow[from=1-2, to=2-2]
	\arrow["G", from=2-1, to=2-2]
  \end{tikzcd}\]
  The above transposition map is equivalent to a map
  \[ \prn{\prod_{b:B} \hom_\Space(X_b, Z_b)} \to \prn{\prod_{a:A} \hom_\Space(E_a, (A \times_B Z)_a)} \simeq \prn{\prod_{b:B} \hom_\Space(E_b,Z_b)}.\]
  By directed univalence, this map corresponds to the map
  \[ \prn{\prod_{b:B} (X_b \to Z_b)} \to\prn{\prod_{b:B} (E_b \to Z_b)} \]
  which in turn acts as precomposition of fibered functors by $j$:
  \[\begin{tikzcd}
	E & X & Z \\
	& B
	\arrow["j", from=1-1, to=1-2]
	\arrow["{u \circ \varphi}"', from=1-1, to=2-2]
	\arrow[from=1-2, to=1-3]
	\arrow["\psi"', from=1-2, to=2-2]
	\arrow["\gamma", from=1-3, to=2-2]
\end{tikzcd}\]
Finally, due to Proposition~\ref{prop:kan:charcofinal}\ref{it:prop:char-left-cof-iv} this map is an equivalence.
\end{proof}

Let $\DeclVar{f}{\GM}{A \to B}$ be a functor with small fibers. Consider the induced pair of adjoints:
\[\begin{tikzcd}
	{\Space^B} & {\Space^A}
	\arrow[""{name=0, anchor=center, inner sep=0}, "{f^*}"', curve={height=12pt}, from=1-1, to=1-2]
	\arrow[""{name=1, anchor=center, inner sep=0}, "{f_!}"', curve={height=12pt}, from=1-2, to=1-1]
	\arrow["\dashv"{anchor=center, rotate=-90}, draw=none, from=1, to=0]
\end{tikzcd}\]
Thus, a square
\[\begin{tikzcd}
	{A'} & A \\
	{B'} & B
	\arrow["u", from=1-1, to=1-2]
	\arrow["{f'}"', from=1-1, to=2-1]
	\arrow["f", from=1-2, to=2-2]
	\arrow["v"', from=2-1, to=2-2]
\end{tikzcd}\]
induces the \emph{Beck--Chevalley square}
\[\begin{tikzcd}
	{\Space^{A'}} & {\Space^{A}} \\
	{\Space^{B'}} & {\Space^{B}}
	\arrow["{f'_!}"', from=1-1, to=2-1]
	\arrow["\beta"{description}, Rightarrow, from=1-1, to=2-2]
	\arrow["{u^*}"', from=1-2, to=1-1]
	\arrow["{f_!}", from=1-2, to=2-2]
	\arrow["{v^*}", from=2-2, to=2-1]
\end{tikzcd}\]
with the $2$-cell $\beta$ (a morphism in $\Space^A \to \Space^{B'}$) being the transpose of the map:
\[\begin{tikzcd}
	{u^*} & {u^*f^*f_!} & {f'^* v^* f_!}
	\arrow["{u^* \eta}", Rightarrow, from=1-1, to=1-2]
	\arrow["\simeq"{description}, draw=none, from=1-2, to=1-3]
\end{tikzcd}\]

Analogously to \ref{prop:action-postcomp}, we give a description of the Beck--Chevalley map as in \cite[Section 8.4]{cisinski:2024}.

Let $\DeclVar{F}{\GM}{A \to \Space}$ and $\varphi$ its unstraightening:
\[\begin{tikzcd}
	X & {\Space_*} \\
	A & \Space
	\arrow[from=1-1, to=1-2]
	\arrow["\varphi"', from=1-1, to=2-1]
	\arrow["\lrcorner"{anchor=center, pos=0.125}, draw=none, from=1-1, to=2-2]
	\arrow[from=1-2, to=2-2]
	\arrow["F", from=2-1, to=2-2]
\end{tikzcd}\]
Recall, that $f_!F$ is given by
%
\[\begin{tikzcd}
	X & Y & {\Space_*} \\
	A & B & \Space
	\arrow["j", from=1-1, to=1-2]
	\arrow["\varphi"', from=1-1, to=2-1]
	\arrow[from=1-2, to=1-3]
	\arrow["\psi"', from=1-2, to=2-2]
	\arrow["\lrcorner"{anchor=center, pos=0.125}, draw=none, from=1-2, to=2-3]
	\arrow[from=1-3, to=2-3]
	\arrow["f", from=2-1, to=2-2]
	\arrow["{f_!F}", from=2-2, to=2-3]
\end{tikzcd}\]
and that the unit $\DeclVar{\eta}{\GM}{F \to f^*f_!F}$ is given by the mediating map:
\[\begin{tikzcd}
	X & Z & Y \\
	& A & B
	\arrow["{\widetilde{\eta}}", dashed, from=1-1, to=1-2]
	\arrow["j", curve={height=-19pt}, from=1-1, to=1-3]
	\arrow["\varphi"', from=1-1, to=2-2]
	\arrow[from=1-2, to=1-3]
	\arrow[from=1-2, to=2-2]
	\arrow["\lrcorner"{anchor=center, pos=0.125}, draw=none, from=1-2, to=2-3]
	\arrow["\psi", from=1-3, to=2-3]
	\arrow["f", from=2-2, to=2-3]
\end{tikzcd}\]
Pulling back along $u$ and $v$, resp., and factoring the upper horizontal map yields the square
\[\begin{tikzcd}
	& {Y'} \\
	{A' \times_A X} && {B' \times_B Y} \\
	{A'} && {B'}
	\arrow["r", from=1-2, to=2-3]
	\arrow["{{j'}}", from=2-1, to=1-2]
	\arrow["{j '':={f' \times_fj}}", dashed, from=2-1, to=2-3]
	\arrow["{{u^*\varphi}}"', from=2-1, to=3-1]
	\arrow["{{v^*\psi}}", from=2-3, to=3-3]
	\arrow["{{f'}}", from=3-1, to=3-3]
\end{tikzcd}\]
where $j'$ is initial and $r$ is a covariant fibration. We consider the composite $\DeclVar{\psi' := v^*\psi \circ r}{\GM}{Y' \to B}$ whence
\[\begin{tikzcd}
	{Y'} & {\Space_*} & {B' \times_B Y} & Y & {\Space_*} \\
	B & {\Space_*} & {B'} & B & \Space
	\arrow[from=1-1, to=1-2]
	\arrow["{\psi'}"', from=1-1, to=2-1]
	\arrow["\lrcorner"{anchor=center, pos=0.125}, draw=none, from=1-1, to=2-2]
	\arrow[from=1-2, to=2-2]
	\arrow[from=1-3, to=1-4]
	\arrow["{v^*\psi}"', from=1-3, to=2-3]
	\arrow["\lrcorner"{anchor=center, pos=0.125}, draw=none, from=1-3, to=2-4]
	\arrow[from=1-4, to=1-5]
	\arrow["\psi"', from=1-4, to=2-4]
	\arrow["\lrcorner"{anchor=center, pos=0.125}, draw=none, from=1-4, to=2-5]
	\arrow[from=1-5, to=2-5]
	\arrow["{f'_!u^*(F)}", from=2-1, to=2-2]
	\arrow["v", from=2-3, to=2-4]
	\arrow["{v^*f_!F}"', curve={height=12pt}, from=2-3, to=2-5]
	\arrow["{f_!F}", from=2-4, to=2-5]
\end{tikzcd}\]
and the Beck--Chevalley transformation corresponds to the fibered map $\DeclVar{r}{\GM}{Y' \to_B B' \times_B Y}$ over $B$.

\begin{thm}[Smooth base change, \protect{\cite[Theorem 8.4.1]{cisinski:2024}}]
  Consider a pullback square
  \[\begin{tikzcd}
	{A'} & A \\
	{B'} & B
	\arrow["u", from=1-1, to=1-2]
	\arrow["{{f'}}"', from=1-1, to=2-1]
	\arrow["\lrcorner"{anchor=center, pos=0.125}, draw=none, from=1-1, to=2-2]
	\arrow["f", from=1-2, to=2-2]
	\arrow["v"', from=2-1, to=2-2]
\end{tikzcd}\]
of small categories where $v$ is smooth. Then, the Beck--Chevalley transformation is an equivalence
\[ f'_! u^* \simeq v^* f_!.\]
\end{thm}

\begin{proof}
We will show that the map $r$ from the preceding description is an equivalence by showing that it is also initial (and it is a covariant fibration by construction). We have the following commutative cube:
\[\begin{tikzcd}
	{A' \times_A X} && {B' \times_B Y} \\
	& X && Y \\
	{A'} && {B'} \\
	& A && B
	\arrow["{j''}"{description}, from=1-1, to=1-3]
	\arrow[from=1-1, to=2-2]
	\arrow["{u^*\phi}"{description, pos=0.3}, from=1-1, to=3-1]
	\arrow["\lrcorner"{anchor=center, pos=0.125}, draw=none, from=1-1, to=4-2]
	\arrow["{v'}"{description},from=1-3, to=2-4]
	\arrow["{v^*\psi}"{description, pos=0.3}, from=1-3, to=3-3]
	\arrow["\lrcorner"{anchor=center, pos=0.125}, draw=none, from=1-3, to=4-4]
	\arrow["j"{description, pos=0.4}, from=2-2, to=2-4, crossing over]
	\arrow["\psi"{description, pos=0.3}, from=2-4, to=4-4]
	\arrow["{f'}"{description, pos=0.7}, from=3-1, to=3-3]
	\arrow["u"{description}, from=3-1, to=4-2]
	\arrow["v"{description}, from=3-3, to=4-4]
	\arrow["f"{description, pos=0.6}, from=4-2, to=4-4]
  \arrow["\phi"{description, pos=0.3}, from=2-2, to=4-2, crossing over]
  \arrow["\lrcorner"{anchor=center, pos=0.125}, draw=none, from=3-1, to=4-4]
\end{tikzcd}\]
By the pullback lemma, the top square is a pullback, too. Since $v$ is smooth and $j$ is initial, $j''$ is initial. But $j'' = r \circ j'$ so $r$ is also initial due to cancelation.
\end{proof}

One can also prove the analogous \emph{proper} base change formula as in \cite[Theorem 8.4.6]{cisinski:2024}.

%
%

%

\section{Conclusions and future work}
\label{sec:conc}

We have introduced and studied the impact of the $\infty$-categorical Yoneda embedding in
\STT{}. This includes the development of classical concepts (Kan extensions, adjoints, (co)limits,
\etc{}), all in the synthetic $\infty$-categorical setting. While some of the basic theory had been
investigated in \STT{} already, we were able to produce the first non-trivial concrete examples of,
\eg{}, adjunctions (Theorem~\ref{thm:adj:precomp-is-right-adj}) and give several more refined versions of
existing theorems (Theorem~\ref{thm:tw:nat-iso}) which more closely match their standard counterparts.

\subsection{Related work}

There are several closely related type-theoretic approaches to synthetic ($\infty$-)category
theory. We may roughly divide these into (1) directed type theory, where every type is a category
but various operations ($\Prod{}$) must be restricted, and (2) variations on simplicial type
theory. For instance, many directed type theories have been proposed and studied over the years~%
\parencite{%
  licata:2011,%
  warren:2013,%
  nuyts:2015,%
  kavvos:directed:2019,%
  buchholtz:2019,%
  kolomatskaia:2023,%
  weaver:2020,%
  weaver:phd,%
  ahrens:2023,%
  north:2018,%
  nuyts:2020,%
  neumann:2024,%
  }.
In general, while these type theories are a promising approach to formalize category theory in type
theory, none of them have thus far received as much attention as \STT{} and, consequently, none have
developed category theory to the extent of this work. Furthermore, it is substantially harder to
design a directed type theory in this style (as it is a more radical alteration of the basic rules
of type theory) and most proposals handle only $1$-category theory rather than
$\prn{\infty,1}$-categories. We note, however, that some of these type theories do include a version
of Theorem~\ref{thm:tw:strong-yoneda} in the form of directed path
induction~\parencite{north:2018,nuyts:2020,neumann:2024}. Given, however, that few of our arguments
rely on types which are not categories, we expect many of them to transfer to sufficiently rich
future variants of directed type theory.

Other variations of simplicial type theory have been considered in the literature. For instance,
several papers use additional judgmental structure (extension types) to get more definitional
equalities around
hom-types~\parencite{riehl:2017,bardomiano:2021,weinberger:phd,buchholtz:2023,weinberger:twosided:2024,weinberger:sums:2024}
at the cost of making the interval a second-class type similar to two-level type
theory~\parencite{annenkov:2023,voevodsky:2013}. Other versions have favored a cubical
interval~\parencite{gratzer:2024} or even a cubical interval atop a cubical version of
\HoTT{}~\parencite{weaver:2020,weaver:phd}. Aside from the addition of modalities, our version of
\STT{} is deliberately minimalistic: we use only ordinary \HoTT{} with a handful of
postulates. Accordingly, our results can be interpret into essentially any incarnation of modal
\STT{} and does not rely on extra definitional equalities.

Finally, there are many attempts to formulate more conceptual and synthetic foundations for
$\infty$-category theory which do not rely on type theory. For instance, the $\infty$-cosmos program
of \textcite{riehl:2022} aims to give a systematic account of the formal category theory and
model-independence using 2-category theory. On the other hand, most practitioners in the field
attempt to give looser ``model independent'' arguments which avoid relying on explicit computations
as much as possible. We have successfully translated some of these arguments into our framework,
proving that this informal discipline is effective (\eg{}, Section~\ref{sec:kan}). More recently,
\textcite{cisinski:2024} have begun to redevelop $\infty$-category theory in a deliberately informal
and high-level language, splitting the difference between a formal theory like \STT{} and the usual
``model-independent'' discipline of practitioners. We expect that their arguments can be translated
into \STT{} and we have shown that some of their primitive axioms are \emph{provable}
in \STT{} (\eg{},
Theorems~\ref{thm:prelims:space} and \ref{thm:prelims:full-subcat} and Lemma~\ref{lem:adj:functoriality-of-init-objects}).

\subsection{Future work}

Many promising avenues for future work remain to be explored. While we have focused on presheaf
categories and immediate consequences of their theory, we plan to port other foundational results
from category theory (presentable and accessible categories, Bousfield localizations, topos theory,
\etc{}) into \STT{}.  It would also be desirable to adapt more parts of the internal
$\infty$-category theory and $\infty$-topos theory of Martini and
Wolf~\parencite{martini:cocartesianfibrations:2022,martini:colimitscocompletions:2021,martini:phd,martini:presentable:2022,martini:propermorphisms:2023,martini:yonedaslemma:2021,wolf:phd}
to \STT{}.
Additionally, we hope to extend a proof assistant like Agda~\parencite{agda}
with the necessary support for modalities to give machine-checked versions of the proofs in this
paper. On the foundational side, \STT{} presently relies on a handful of axioms
(Appendix~\ref{sec:axioms}) and therefore satisfies only normalization and not canonicity. In future work,
we hope to examine which of these principles can be given computational interpretations and to what
extent one can `compute' with synthetic $\infty$-categories.

\appendix
%
\section{The formal rules of \texorpdfstring{\MTT{}}{MTT}}
\label{sec:mtt-rules}

The formal syntax of \MTT{} is comprised of four judgments: $\IsCx{\Gamma}$,
$\IsSb{\delta}{\Delta}$, $\IsTm{a}{A}$, and $\IsTy{A}$. We list the relevant novel rules for these
judgments below:%
\def\MathparLineskip{\lineskip=11pt}
\begin{mathparpagebreakable}
  \JdgFrame{\IsCx{\Gamma}}
  \\
  \inferrule{ }{
    \IsCx{\EmpCx}
  }
  \and
  \inferrule{
    \IsCx{\Gamma}
  }{
    \IsCx{\LockCx{\Gamma}}
  }
  \and
  \inferrule{
    \IsCx{\Gamma}
    \\
    \IsTy[\LockCx{\Gamma}]{A}
  }{
    \IsCx{\MECx{\Gamma}{A}}
  }
  \\
  \inferrule{
    \IsCx{\Gamma}
  }{
    \EqCx{\LockCx{\Gamma}[\ArrId{}]}{\Gamma}
    \\
    \EqCx{\LockCx{\LockCx{\Gamma}[\mu]}[\nu]}{\LockCx{\Gamma}[\mu \circ \nu]}
  }
  \\
  \JdgFrame{\IsSb{\delta}{\Delta}}
  \\
  \inferrule{ }{
    \IsSb{\EmpSb}{\EmpCx}
  }
  \and
  \inferrule{
    \IsCx{\Gamma}
    \\
    \IsTy[\LockCx{\Gamma}]{A}
  }{
    \IsSb[\MECx{\Gamma}{A}]{\Wk}{\Gamma}
  }
  \and
  \inferrule{
    \IsSb{\delta}{\Delta}
    \\
    \IsTm[\LockCx{\Gamma}]{a}{A}
  }{
    \IsSb{\ESb{\delta}{a}}{\MECx{\Delta}{A}}
  }
  \and
  \inferrule{
    \IsSb{\delta}{\Delta}
  }{
    \IsSb[\LockCx{\Gamma}]{\LockSb{\delta}}{\LockCx{\Delta}}
  }
  \and
  \inferrule{
    \IsCx{\Gamma}
    \\
    \Mor[\alpha]{\mu}{\nu}
  }{
    \IsSb[\LockCx{\Gamma}[\nu]]{\KeySb{\alpha}{\Gamma}}{\LockCx{\Gamma}[\mu]}
  }
  \\
  \inferrule{
    \IsSb{\gamma}{\EmpCx}
  }{
    \EqSb{\EmpSb}{\gamma}{\EmpCx}
  }
  \and
  \inferrule{
    \IsSb{\delta}{\Delta}
    \\
    \IsTm[\LockCx{\Gamma}]{a}{A}
  }{
    \EqSb{\Wk \circ \prn{\ESb{\delta}{a}}}{\delta}{\Delta}
  }
  \and
  \inferrule{
    \IsSb{\delta}{\MECx{\Delta}{A}}
  }{
    \EqSb{\ESb{\prn{\Wk \circ \delta}}{\Sb{\Var}{\delta}}}{\delta}{\MECx{\Delta}{A}}
  }
  \and
  \inferrule{
    \IsSb{\delta}{\Delta}
  }{
    \EqSb{\LockSb{\delta}[\ArrId{}]}{\delta}{\Delta}
    \\
    \EqSb[
      \LockCx{\Gamma}[\nu\circ\mu]
    ]{
      \LockSb{\delta}[\nu\circ\mu]
    }{
      \LockSb{\LockSb{\delta}[\nu]}
    }{
      \LockCx{\Delta}[\nu\circ\mu]
    }
  }
  \and
  \inferrule{
    \IsCx{\Gamma}
  }{
    \EqSb[\LockCx{\Gamma}]{\KeySb{\ArrId{}}{\Gamma}}{\ArrId{}}{\LockCx{\Gamma}}
    \\
    \EqSb[
      \LockCx{\Gamma}[\xi]
    ]{
      \KeySb{\alpha}{\Gamma} \circ \KeySb{\beta}{\Gamma}
    }{
      \KeySb{\alpha \circ \beta}{\Gamma}
    }{
      \LockCx{\Gamma}[\mu]
    }
  }
  \and
  \inferrule{
    \IsSb{\delta}{\Delta}
    \\
    \mu \le \nu
  }{
    \EqSb[
      \LockCx{\Gamma}[\nu]
    ]{
      \KeySb{\alpha}{\Delta} \circ \LockCx{\delta}
    }{
      \LockSb{\delta} \circ \KeySb{\alpha}{\Gamma}
    }{
      \LockCx{\Delta}[\mu]
    }
  }
  \and
  \inferrule{
    \IsCx{\Gamma}
    \\
    \Mor[\alpha]{\mu_0}{\nu_0}
    \\
    \Mor[\beta]{\mu_1}{\nu_1}
  }{
    \LockCx{\Gamma}[\nu_1 \circ \nu_0] \vdash
    \substack{
      \LockCx{\KeySb{\beta}{\Gamma}}[\mu_0]
      \circ
      \KeySb{\alpha}{\Gamma}
      \\
      =
      \KeySb{\beta \HComp \alpha}{\Gamma}
    }
    : \LockCx{\Gamma}[\mu_1 \circ \mu_0]
  }
  \\
  \JdgFrame{\IsTy{A}}
  \\
  \inferrule{
    \IsTy[\LockCx{\Gamma}]{A}
  }{
    \IsTy{\Modify{A}}
  }
  \and
  \inferrule{
    \IsSb{\delta}{\Delta}
    \\
    \IsTy[\LockCx{\Delta}]{A}
  }{
    \EqTy{\Sb{\Modify{A}}{\delta}}{\Modify{\Sb{A}{\LockSb{\delta}}}}
  }
  \\
  \JdgFrame{\IsTm{a}{A}}
  \\
  \inferrule{
    \IsTy[\LockCx{\Gamma}]{A}
  }{
    \IsTm[\LockCx{\MECx{\Gamma}{A}}]{\Var}{\Sb{A}{\LockSb{\Wk}}}
  }
  \and
  \inferrule{
    \IsTm[\LockCx{\Gamma}]{a}{A}
  }{
    \IsTm{\MkMod{a}}{\Modify{A}}
  }
  \and
  \inferrule{
    \IsTy[\MECx{\Gamma}[\nu]{\Modify{A}}]{B}
    \\
    \IsTm[\MECx{\Gamma}[\nu\circ\mu]{A}]{b}{\Sb{B}{\ESb{\Wk}{\MkMod[\mu]{\Var}}}}
    \\
    \IsTm[\LockCx{\Gamma}[\nu]]{a}{\Modify{A}}
  }{
    \IsTm{\LetMod{a}{b}}{\Sb{B}{\ESb{\ISb}{a}}}
  }
  \\
  \inferrule{
    \IsTy[\MECx{\Delta}[\nu]{\Modify{A}}]{B}
    \\
    \IsTm[\MECx{\Delta}[\nu\circ\mu]{A}]{b}{\Sb{B}{\ESb{\Wk}{\MkMod[\mu]{\Var}}}}
    \\
    \IsTm[\LockCx{\Delta}[\nu]]{a}{\Modify{A}}
    \\
    \IsSb{\delta}{\Delta}
  }{
    \IsTm{
      \substack{
        \LetMod{\Sb{a}{\LockSb{\delta}[\nu]}}{\Sb{b}{\ESb{\prn{\delta \circ \Wk}}{\Var}}}
        \\
        {}={}
        \Sb{\prn{\LetMod{a}{b}}}{\delta}
      }
    }{
      \Sb{B}{\ESb{\delta}{a}}
    }
  }
  \and
  \inferrule{
    \IsTm[\LockCx{\Gamma}]{a}{A}
    \\
    \IsSb{\delta}{\Delta}
  }{
    \EqTm{\Sb{\MkMod{a}}{\delta}}{\MkMod{\Sb{a}{\LockSb{\delta}}}}{\Modify{\Sb{A}{\LockSb{\delta}}}}
  }
  \and
  \inferrule{
    \IsSb{\delta}{\Delta}
    \\
    \IsTm[\LockCx{\Gamma}]{a}{\Sb{A}{\LockSb{\delta}}}
    \\
    \IsTy[\LockCx{\Delta}]{A}
  }{
    \EqTm[\LockCx{\Gamma}]{\Sb{\Var}{\LockCx{\ESb{\delta}{a}}}}{a}{\Sb{A}{\LockSb{\delta}}}
  }
  \and
  \hspace*{-5pt}\inferrule{
    \IsTy[\MECx{\Gamma}[\nu]{\Modify{A}}]{B}
    \\
    \IsTm[\MECx{\Gamma}[\nu\circ\mu]{A}]{b}{\Sb{B}{\ESb{\Wk}{\MkMod[\mu]{\Var}}}}
    \\
    \IsTm[\LockCx{\Gamma}[\nu]]{a}{\Modify{A}}
  }{
    \EqTm{\prn{\LetMod{\MkMod{a}}{b}}}{\Sb{b}{\ESb{\ISb}{a}}}{\Sb{B}{\ESb{\ISb}{\MkMod{a}}}}
  }
\end{mathparpagebreakable}

\section{The complete list of axioms}
\label{sec:axioms}


\intaxiom*
\crispidaxiom*
\intopaxiom*
\intdiscaxiom*
\intglobalaxiom*
\cubessepaxiom*
\twarraxiom*

The following \emph{duality axiom} was first studied by \textcite{blechschmidt:2023} and implies
that, \eg{}, $\Int$ is a category. Closely related axioms and consequences are considered by
\textcite{pugh:2025} and \textcite{cherubini:2023}. We did not introduce it in the main body of the
paper as it was not explicitly invoked in any of our proofs.
\begin{axiom}
  If $A$ is a finitely presented $\Int$-algebra \textup(\ie{}, $A$ is a bounded distributive lattice
  equivalent to $\Int[x_1,\ldots,x_n]$ quotiented by finitely many relations\textup) and
  $\Hom[\Int\mathrm{Alg}]{A}{\Int}$ is the type of $\Int$-algebra homomorphisms, then the map
  $\lambda a\,f.\,f\prn{a} : A \to \prn{\Hom[\Int\mathrm{Alg}]{A}{\Int} \to \Int}$ is an equivalence.
\end{axiom}

%
%
%
%
 

\begin{thebibliography}{99}

\providecommand{\url}[1]{{#1}}
\providecommand{\urlprefix}{URL }
\expandafter\ifx\csname urlstyle\endcsname\relax
  \providecommand{\doi}[1]{DOI~\discretionary{}{}{}#1}\else
  \providecommand{\doi}{DOI~\discretionary{}{}{}\begingroup
  \urlstyle{rm}\Url}\fi

\bibitem{annenkov:2023}
Danil Annenkov, Paolo Capriotti, Nicolai Kraus, and Christian Sattler.
\newblock Two-level type theory and applications.
\newblock \emph{Mathematical Structures in Computer Science}, 33(8):688--743, 2023.
\newblock \doi{10.1017/S0960129523000130}.

\bibitem{ahrens:2023}
Benedikt Ahrens, Paige Randall North, and Niels van der Weide.
\newblock Bicategorical type theory: semantics and syntax.
\newblock \emph{Mathematical Structures in Computer Science}, 33(10):868--912, 2023.
\newblock \doi{10.1017/s0960129523000312}.

\bibitem{anel:2024}
Mathieu Anel and Jonathan Weinberger.
\newblock Smooth and proper maps with respect to a fibration.
\newblock \emph{Mathematical Structures in Computer Science}, 34(9):971--984, 2024.

\bibitem{bardomiano:2021}
César Bardomiano Martínez.
\newblock Limits and colimits of synthetic $\infty$-categories.
\newblock arXiv:2202.12386, 2022.

\bibitem{bergner:2018}
Julia Bergner.
\newblock \emph{The Homotopy Theory of $(\infty,1)$-Categories}.
\newblock Cambridge University Press, 2018.
\newblock \doi{10.1017/9781316181874}.

\bibitem{bahr:2017}
Patrick Bahr, Hans Bugge Grathwohl, and Rasmus Ejlers Møgelberg.
\newblock The clocks are ticking: No more delays!
\newblock In \emph{2017 32nd Annual ACM/IEEE Symposium on Logic in Computer Science (LICS)}, 2017.
\newblock \doi{10.1109/LICS.2017.8005097}.

\bibitem{birkedal:2020}
Lars Birkedal, Ranald Clouston, Bassel Mannaa, Rasmus Ejlers Møgelberg, Andrew M. Pitts, and Bas Spitters.
\newblock Modal dependent type theory and dependent right adjoints.
\newblock \emph{Mathematical Structures in Computer Science}, 30(2):118--138, 2020.
\newblock \doi{10.1017/S0960129519000197}.

\bibitem{blechschmidt:2023}
Ingo Blechschmidt.
\newblock A general Nullstellensatz for generalized spaces.
\newblock Draft, 2023.
\newblock \url{https://rawgit.com/iblech/internal-methods/master/paper-qcoh.pdf}.

\bibitem{buchholtz:2019}
Ulrik Buchholtz.
\newblock Higher structures in homotopy type theory.
\newblock In \emph{Reflections on the Foundations of Mathematics}, pages 151--172. Springer, 2019.
\newblock \doi{10.1007/978-3-030-15655-8_7}.

\bibitem{buchholtz:2023}
Ulrik Buchholtz and Jonathan Weinberger.
\newblock Synthetic fibered $(\infty,1)$-category theory.
\newblock \emph{Higher Structures}, 7(1):74--165, 2023.
\newblock \doi{10.21136/HS.2023.04}.

\bibitem{cherubini:2023}
Felix Cherubini, Thierry Coquand, and Matthias Hutzler.
\newblock A foundation for synthetic algebraic geometry.
\newblock \emph{Mathematical Structures in Computer Science}, 34(9):1008--1053, 2024.
\newblock \doi{10.1017/S0960129524000239}.

\bibitem{cisinski:2024}
Denis-Charles Cisinski, Bastiaan Cnossen, Kim Nguyen, and Tashi Walde.
\newblock Formalization of Higher Categories.
\newblock Lecture notes, 2024.

\bibitem{cisinski:2019}
Denis-Charles Cisinski.
\newblock \emph{Higher Categories and Homotopical Algebra}.
\newblock Cambridge University Press, 2019.
\newblock \doi{10.1017/9781108588737}.

\bibitem{gonthier:2013}
Georges Gonthier et al.
\newblock A Machine-Checked Proof of the Odd Order Theorem.
\newblock In \emph{Interactive Theorem Proving}, pages 163--179. Springer, 2013.
\newblock \doi{10.1007/978-3-642-39634-2_14}.

\bibitem{gratzer:2020}
Daniel Gratzer, G. A. Kavvos, Andreas Nuyts, and Lars Birkedal.
\newblock Multimodal Dependent Type Theory.
\newblock In \emph{Proceedings of the 35th Annual ACM/IEEE Symposium on Logic in Computer Science}, 2020.
\newblock \doi{10.1145/3373718.3394736}.

\bibitem{gratzer:mtt-journal:2021}
Daniel Gratzer, G. A. Kavvos, Andreas Nuyts, and Lars Birkedal.
\newblock Multimodal Dependent Type Theory.
\newblock \emph{Logical Methods in Computer Science}, 17(3), 2021.
\newblock \doi{10.46298/lmcs-17(3:11)2021}.

\bibitem{gratzer:normalization:2022}
Daniel Gratzer.
\newblock Normalization for Multimodal Type Theory.
\newblock In \emph{Proceedings of the 37th Annual ACM/IEEE Symposium on Logic in Computer Science}, 2022.
\newblock \doi{10.1145/3531130.3532398}.

\bibitem{gratzer:phd}
Daniel Gratzer.
\newblock Syntax and semantics of modal type theory.
\newblock PhD thesis, Aarhus University, 2023.

\bibitem{gratzer:2024}
Daniel Gratzer, Jonathan Weinberger, and Ulrik Buchholtz.
\newblock Directed univalence in simplicial homotopy type theory.
\newblock arXiv:2407.09146, 2024.

\bibitem{hyland:1991}
J. M. E. Hyland.
\newblock First steps in synthetic domain theory.
\newblock In \emph{Category Theory}, pages 131--156. Springer, 1991.
\newblock \doi{10.1007/bfb0084217}.

\bibitem{kavvos:directed:2019}
G. A. Kavvos.
\newblock A quantum of direction, 2019.
\newblock \url{https://seis.bristol.ac.uk/~tz20861/papers/meio.pdf}.

\bibitem{kock:2006}
Anders Kock.
\newblock \emph{Synthetic Differential Geometry}.
\newblock Cambridge University Press, 2 edition, 2006.

\bibitem{kudasov:2024}
Nikolai Kudasov, Emily Riehl, and Jonathan Weinberger.
\newblock Formalizing the $\infty$-Categorical Yoneda Lemma.
\newblock In \emph{Proceedings of the 13th ACM SIGPLAN International Conference on Certified Programs and Proofs}, pages 274--290, 2024.
\newblock \doi{10.1145/3636501.3636945}.

\bibitem{kolomatskaia:2023}
Astra Kolomatskaia and Michael Shulman.
\newblock Displayed Type Theory and Semi-Simplicial Types.
\newblock arXiv:2311.18781, 2023.

\bibitem{licata:2011}
Daniel R. Licata and Robert Harper.
\newblock 2-Dimensional Directed Type Theory.
\newblock \emph{Electronic Notes in Theoretical Computer Science}, 276:263--289, 2011.
\newblock \doi{10.1016/j.entcs.2011.09.026}.

\bibitem{lumsdaine:2019}
Peter LeFanu Lumsdaine and Michael Shulman.
\newblock Semantics of higher inductive types.
\newblock \emph{Mathematical Proceedings of the Cambridge Philosophical Society}, 169(1):159--208, 2019.
\newblock \doi{10.1017/s030500411900015x}.

\bibitem{lurie:2009}
Jacob Lurie.
\newblock \emph{Higher Topos Theory}.
\newblock Princeton University Press, 2009.

\bibitem{maclane:1978}
Saunders Mac Lane.
\newblock \emph{Categories for the Working Mathematician}.
\newblock Springer, 1978.

\bibitem{martini:cocartesianfibrations:2022}
Louis Martini.
\newblock Cocartesian fibrations and straightening internal to an $\infty$-topos.
\newblock arXiv:2204.00295, 2022.

\bibitem{martini:yonedaslemma:2021}
Louis Martini.
\newblock Yoneda's lemma for internal higher categories.
\newblock arXiv:2103.17141, 2022.

\bibitem{martini:phd}
Louis Martini.
\newblock Internal Higher Category Theory.
\newblock PhD thesis, NTNU, 2024.

\bibitem{mukherjee:2023}
Chirantan Mukherjee and Nima Rasekh.
\newblock Twisted Arrow Construction for Segal Spaces.
\newblock arXiv:2203.01788, 2023.

\bibitem{myers:2023}
David Jaz Myers and Mitchell Riley.
\newblock Commuting Cohesions.
\newblock arXiv:2301.13780, 2023.

\bibitem{martini:presentable:2022}
Louis Martini and Sebastian Wolf.
\newblock Presentability and topoi in internal higher category theory.
\newblock arXiv:2209.05103, 2022.

\bibitem{martini:colimitscocompletions:2021}
Louis Martini and Sebastian Wolf.
\newblock Colimits and cocompletions in internal higher category theory.
\newblock \emph{Higher Structures}, 8(1):97--192, 2024.
\newblock \doi{10.21136/HS.2024.03}.

\bibitem{martini:propermorphisms:2023}
Louis Martini and Sebastian Wolf.
\newblock Proper morphisms of $\infty$-topoi.
\newblock arXiv:2311.08051, 2024.

\bibitem{neumann:2024}
Jacob Neumann and Thorsten Altenkirch.
\newblock The Category Interpretation of Directed Type Theory, 2024.
\newblock \url{https://jacobneu.github.io/research/preprints/catModel-2024.pdf}.

\bibitem{north:2018}
Paige Randall North.
\newblock Towards a Directed Homotopy Type Theory.
\newblock \emph{Electronic Notes in Theoretical Computer Science}, 347:223--239, 2019.
\newblock \doi{10.1016/j.entcs.2019.09.012}.

\bibitem{nuyts:2015}
Andreas Nuyts.
\newblock Towards a Directed Homotopy Type Theory based on 4 Kinds of Variance.
\newblock Master's thesis, KU Leuven, 2015.

\bibitem{nuyts:2020}
Andreas Nuyts.
\newblock A Vision for Natural Type Theory, 2020.
\newblock \url{https://anuyts.github.io/files/nattt-vision.pdf}.

\bibitem{pugh:2025}
Leoni Pugh and Jonathan Sterling.
\newblock When is the partial map classifier a Sierpiński cone?
\newblock arXiv:2504.06789, 2025. To appear at LICS 2025.

\bibitem{ramzi:bousfield-kan}
Maxime Ramzi.
\newblock Deducing the Bousfield-Kan formula for homotopy (co)limits from first principles, 2021.
\newblock \url{https://sites.google.com/view/maxime-ramzi-en/notes/bousfield-kan}.

\bibitem{riehl:2014}
Emily Riehl.
\newblock \emph{Categorical homotopy theory}.
\newblock Cambridge University Press, 2014.

\bibitem{riehl:2023}
Emily Riehl.
\newblock Could $\infty$-Category Theory Be Taught to Undergraduates?
\newblock \emph{Notices of the American Mathematical Society}, 70(5), 2023.
\newblock \doi{10.1090/noti2692}.

\bibitem{riehl:2024}
Emily Riehl.
\newblock On the $\infty$-topos semantics of homotopy type theory.
\newblock \emph{Bulletin of the London Mathematical Society}, 56(2):461--517, 2024.
\newblock \doi{10.1112/blms.12997}.

\bibitem{riehl:2017}
Emily Riehl and Michael Shulman.
\newblock A type theory for synthetic $\infty$-categories.
\newblock \emph{Higher Structures}, 1(1):147--224, 2017.
\newblock \doi{10.21136/HS.2017.06}.

\bibitem{rijke:2020}
Egbert Rijke, Michael Shulman, and Bas Spitters.
\newblock Modalities in homotopy type theory.
\newblock \emph{Logical Methods in Computer Science}, 16(1), 2020.
\newblock arXiv:1706.07526.

\bibitem{russell:1919}
Bertrand Russell.
\newblock \emph{Introduction to Mathematical Logic}.
\newblock George Allen \& Unwin, 1919.

\bibitem{riehl:2022}
Emily Riehl and Dominic Verity.
\newblock \emph{Elements of $\infty$-Category Theory}.
\newblock Cambridge University Press, 2022.
\newblock \doi{10.1017/9781108936880}.

\bibitem{shulman:2018}
Michael Shulman.
\newblock Brouwer's fixed-point theorem in real-cohesive homotopy type theory.
\newblock \emph{Mathematical Structures in Computer Science}, 28(6):856--941, 2018.
\newblock \doi{10.1017/S0960129517000147}.

\bibitem{shulman:2019}
Michael Shulman.
\newblock All $(\infty,1)$-toposes have strict univalent universes.
\newblock arXiv:1904.07004, 2019.

\bibitem{agda}
Agda Development Team.
\newblock Agda User Manual.
\newblock \url{https://agda.readthedocs.io/en/v2.7.0.1/}.

\bibitem{hottbook}
The Univalent Foundations Program.
\newblock \emph{Homotopy Type Theory: Univalent Foundations of Mathematics}.
\newblock Institute for Advanced Study, 2013.
\newblock \url{https://homotopytypetheory.org/book}.

\bibitem{voevodsky:2013}
Vladimir Voevodsky.
\newblock A simple type system with two identity types, 2012.
\newblock \url{https://www.math.ias.edu/vladimir/sites/math.ias.edu.vladimir/files/HTS.pdf}.

\bibitem{warren:2013}
Michael Warren.
\newblock Directed type theory, 2013.
\newblock Seminar talk. \url{https://www.ias.edu/video/univalent/1213/0410-MichaelWarren}.

\bibitem{weaver:phd}
Matthew Weaver.
\newblock Bicubical Directed Type Theory.
\newblock PhD thesis, Princeton University, 2024.

\bibitem{weinberger:phd}
Jonathan Weinberger.
\newblock A Synthetic Perspective on $(\infty,1)$-Category Theory: Fibrational and Semantic Aspects.
\newblock PhD thesis, Technische Universität Darmstadt, 2022.
\newblock \doi{10.26083/tuprints-00020716}.

\bibitem{weinberger:sums:2024}
Jonathan Weinberger.
\newblock Internal sums for synthetic fibered $(\infty,1)$-categories.
\newblock \emph{Journal of Pure and Applied Algebra}, 228(9):107659, 2024.
\newblock \doi{10.1016/j.jpaa.2024.107659}.

\bibitem{weinberger:twosided:2024}
Jonathan Weinberger.
\newblock Two-sided cartesian fibrations of synthetic $(\infty,1)$-categories.
\newblock \emph{Journal of Homotopy and Related Structures}, 2024.
\newblock \doi{10.1007/s40062-024-00348-3}.

\bibitem{weaver:2020}
Matthew Z. Weaver and Daniel R. Licata.
\newblock A Constructive Model of Directed Univalence in Bicubical Sets.
\newblock In \emph{Proceedings of the 35th Annual ACM/IEEE Symposium on Logic in Computer Science}, 2020.
\newblock \doi{10.1145/3373718.3394794}.

\bibitem{wolf:phd}
Sebastian Wolf.
\newblock Internal Higher Categories and Applications.
\newblock PhD thesis, Universität Regensburg, 2025.
\newblock \url{https://epub.uni-regensburg.de/76465/}.

\end{thebibliography}

\end{document}
